%% ======================================================================
%%  Minicurso: CFD Aplicado à Segurança de Processos
%%  Dispersão de gases, incêndios e explosões --- do zero ao avançado
%%  Compilar com: pdflatex (2x)
%% ======================================================================
\documentclass[10pt,a4paper,twocolumn,openany]{report}

\usepackage[utf8]{inputenc}
\usepackage[T1]{fontenc}
\usepackage[brazil]{babel}
\usepackage{lmodern}
\usepackage{microtype}
\usepackage[a4paper,margin=1.7cm,columnsep=0.65cm,headheight=14pt]{geometry}
\usepackage{amsmath,amssymb,bm}
\usepackage{booktabs,tabularx,array,longtable}
\usepackage{enumitem}
\usepackage[dvipsnames]{xcolor}
\usepackage[most]{tcolorbox}
\usepackage{listings}
\usepackage{fancyhdr}
\usepackage{titlesec}
\usepackage{truncate}
\usepackage[colorlinks=true,linkcolor=azul,urlcolor=azul,citecolor=azul]{hyperref}

% ---------------------------------------------------------------- cores
\definecolor{azul}{RGB}{25,70,115}
\definecolor{azulclaro}{RGB}{232,240,250}
\definecolor{verde}{RGB}{30,110,60}
\definecolor{verdeclaro}{RGB}{233,246,237}
\definecolor{laranja}{RGB}{190,95,10}
\definecolor{laranjaclaro}{RGB}{253,242,228}
\definecolor{roxo}{RGB}{105,45,130}
\definecolor{roxoclaro}{RGB}{243,234,248}
\definecolor{cinza}{RGB}{90,90,90}
\definecolor{cinzaclaro}{RGB}{242,242,242}
\definecolor{vermelho}{RGB}{160,30,30}
\definecolor{vermelhoclaro}{RGB}{251,234,234}
\definecolor{amarelo}{RGB}{160,120,0}
\definecolor{amareloclaro}{RGB}{255,249,225}

% ---------------------------------------------------------------- estilos
% Única caixa do documento: as tarefas (entregáveis). O resto é texto corrido.
\tcbset{enhanced,breakable,fonttitle=\bfseries,sharp corners,boxrule=0.5pt,
        left=2mm,right=2mm,top=1.2mm,bottom=1.2mm,before skip=10pt,after skip=10pt}
\newtcolorbox[auto counter,number within=chapter]{tarefa}[1]{colback=verdeclaro,
  colframe=verde,title={Tarefa~\thetcbcounter\ --- #1}}

\titleformat{\section}{\normalfont\large\bfseries\color{azul}}{\thesection}{0.6em}{}
\titleformat{\subsection}{\normalfont\normalsize\bfseries\color{azul}}{\thesubsection}{0.6em}{}
\titlespacing*{\section}{0pt}{2.2ex plus 1ex minus .3ex}{1ex}
\titlespacing*{\subsection}{0pt}{1.8ex plus .8ex minus .3ex}{0.6ex}

\newcommand{\lead}[2]{\par\medskip\noindent{\color{#1}\bfseries #2}\ }
\newenvironment{conceito}[1]{\section{#1}}{\par}
\newenvironment{professor}{\lead{laranja}{Nota do professor.}\itshape}{\par\medskip}
\newenvironment{armadilha}{\lead{vermelho}{Armadilha clássica.}}{\par\medskip}
\newenvironment{dicionario}[1][Dicionário da equação]{\lead{azul}{#1.}\small}{\par\medskip}
\newenvironment{critico}{\subsection*{\textcolor{roxo}{Pensamento crítico}}}{\par}
\newenvironment{chaves}{\lead{cinza}{Palavras-chave (PT/EN):}\small}{\par}
\newenvironment{refs}{\subsection*{Referências do módulo}\small}{\par}
\newenvironment{aprofundar}{\subsection*{\textcolor{amarelo}{Para se aprofundar}}\small}{\par}
\newcommand{\origem}{\par\smallskip\noindent\textcolor{azul}{\textbf{De onde vem?}}\ }

\newcommand{\enunciado}{\noindent\textbf{Enunciado.}\ }
\newcommand{\entregaveis}{\par\medskip\noindent\textbf{Entregáveis:}}
\newcommand{\criterio}{\par\medskip\noindent\textbf{Critério de sucesso:}\ }
\newcommand{\reflexao}{\par\medskip\noindent\textbf{Reflexão obrigatória (escreva no diário):}}
\newcommand{\objetivos}{\noindent\textbf{Ao final deste módulo você será capaz de:}}

% ---------------------------------------------------------------- código
\lstdefinelanguage{Julia}{
  morekeywords={function,end,for,in,if,else,elseif,while,return,using,import,
    module,struct,mutable,begin,let,const,local,global,do,true,false},
  sensitive=true, morecomment=[l]{\#}, morestring=[b]"}
\lstset{language=Julia,basicstyle=\ttfamily\footnotesize,
  keywordstyle=\color{azul}\bfseries,commentstyle=\color{gray}\itshape,
  stringstyle=\color{BrickRed},frame=leftline,rulecolor=\color{azul!40},framerule=1.2pt,
  breaklines=true,columns=fullflexible,backgroundcolor=\color{gray!5},
  numbers=left,numberstyle=\tiny\color{gray},xleftmargin=1.3em,
  showstringspaces=false}

% ---------------------------------------------------------------- layout
\setlist{itemsep=2pt,topsep=4pt}
\raggedbottom
\makeatletter\setlength{\@fptop}{0pt}\setlength{\@dblfptop}{0pt}\makeatother
\setlength{\emergencystretch}{2.5em}
\addto\captionsbrazil{\renewcommand{\chaptername}{Módulo}}
\pagestyle{fancy}
\fancyhf{}
\fancyhead[L]{\small\color{cinza}\leftmark}
\fancyhead[R]{\small\color{cinza}CFD para Segurança de Processos}
\fancyfoot[C]{\thepage}
\renewcommand{\headrulewidth}{0.4pt}
\renewcommand{\chaptermark}[1]{\markboth{\chaptername\ \thechapter\ --- #1}{}}
\fancyhead[L]{\small\color{cinza}\nouppercase{\truncate{0.62\textwidth}{\leftmark}}}

\hypersetup{pdftitle={Minicurso: CFD Aplicado à Segurança de Processos},
            pdfauthor={Minicurso autodirigido}}

\setcounter{tocdepth}{0}
\begin{document}

% ====================================================================== CAPA
\begin{titlepage}
\centering
\vspace*{2cm}
{\color{azul}\rule{\linewidth}{1.5pt}}\\[0.6cm]
{\Huge\bfseries\color{azul} CFD Aplicado à\\[0.2cm] Segurança de Processos\par}
\vspace{0.4cm}
{\Large Dispersão de gases, incêndios e explosões\\ --- do zero ao avançado ---\par}
{\color{azul}\rule{\linewidth}{1.5pt}}\\[1.2cm]
{\large Minicurso autodirigido, orientado a entregáveis\par}
\vspace{0.6cm}
{\normalsize Teoria $\cdot$ Análise numérica $\cdot$ Programação $\cdot$
Validação $\cdot$ Prática industrial offshore\par}
\vfill
\begin{minipage}{0.8\linewidth}\small
\textcolor{azul}{\textbf{Como ler este documento.}} Cada módulo traz objetivos,
teoria, dicionário das equações, palavras-chave, referências, leituras
comentadas, notas do professor, perguntas de pensamento crítico e
\textbf{tarefas com entregáveis}. Não avance de módulo sem fechar os
entregáveis do anterior: CFD é uma pirâmide, e a base é análise numérica.
\end{minipage}
\vspace{1cm}
{\small Versão 1.0 --- outubro de 2026\par}
\end{titlepage}

\tableofcontents

% ====================================================================== PREFÁCIO
\chapter*{Antes de começar: uma conversa franca}
\addcontentsline{toc}{chapter}{Antes de começar: uma conversa franca}
\markboth{Antes de começar}{}

Depois de muitos anos ensinando CFD para engenheiros de segurança de processos,
aprendi que existem dois tipos de alunos que fracassam. O primeiro acredita que
CFD é ``apertar o botão do software'': monta a geometria, escolhe o
$k$-$\varepsilon$ porque ``todo mundo usa'', gera um contorno colorido e o
coloca no relatório. O segundo se perde na matemática, passa meses
otimizando um solver de Poisson e nunca chega a perguntar \emph{qual decisão de
engenharia aquele número vai suportar}.

Este minicurso foi desenhado para evitar os dois fracassos. Você vai
\textbf{programar} os blocos fundamentais com as próprias mãos (porque só
entende um esquema numérico quem já o viu explodir), e vai \textbf{conectar
cada bloco a uma pergunta real de segurança}: até onde vai a nuvem
inflamável? Onde coloco os detectores? Qual a sobrepressão de projeto de um
módulo de processo de FPSO? Qual o fluxo térmico que a parede corta-fogo
precisa suportar?

\section*{Filosofia do curso}
\begin{enumerate}
  \item \textbf{Entregável ou não aconteceu.} Cada teoria termina em uma tarefa
        com produto concreto: código, gráfico, tabela, nota técnica curta.
  \item \textbf{Do analítico ao numérico ao industrial.} Sempre que possível,
        compare sua simulação com uma solução analítica ou experimental antes
        de confiar nela.
  \item \textbf{O erro é o professor.} Você será instruído a quebrar seus
        próprios códigos de propósito (instabilidade, malha grossa, condição de
        contorno errada). Registre o que viu.
  \item \textbf{Ceticismo saudável.} Todo resultado de CFD é uma hipótese
        sobre a realidade, condicionada a dezenas de escolhas. Aprender a
        listar essas escolhas é metade da competência profissional.
\end{enumerate}

\section*{Estrutura de cada módulo}
\begin{itemize}
  \item \textbf{Objetivos} --- o que você saberá fazer ao final.
  \item \textbf{Teoria} em seções numeradas --- o mínimo rigoroso.
  \item \textbf{Palavras-chave} em português e inglês --- a literatura é quase
        toda em inglês; aprenda o vocabulário para buscar artigos.
  \item \textbf{Tarefas} --- enunciado, entregáveis, critério de sucesso e
        reflexão obrigatória. São o único elemento em caixa, para que você as
        encontre rapidamente.
  \item \textbf{Notas do Professor} e \textbf{Armadilhas clássicas} --- o que
        os livros raramente dizem.
  \item \textbf{Pensamento Crítico} --- perguntas sem resposta única, para o
        diário de bordo e para discussão com colegas e orientadores.
  \item \textbf{Dicionário da equação} --- o significado e a unidade de cada
        símbolo, e \emph{de onde veio} a equação (quem, quando, a partir de
        qual raciocínio).
  \item \textbf{Referências} --- livros-texto, artigos seminais e normas.
  \item \textbf{Para se aprofundar} --- leituras comentadas: o que ler, em que
        ordem, e por quê.
\end{itemize}

\section*{Ferramentas recomendadas}
\begin{itemize}
  \item \textbf{Linguagem principal:} Julia --- se você nunca a usou, faça
        antes a \emph{Oficina J} (logo após o Módulo 0). Julia oferece desempenho próximo de C com
        sintaxe de alto nível; ótimo para os solvers próprios). Python
        (NumPy/SciPy/Matplotlib) é aceitável. Rust é recomendado no projeto
        final para os pós-processadores que você queira transformar em
        ferramenta de produção.
  \item \sloppy\textbf{Pacotes Julia úteis:} \texttt{Plots.jl} ou \texttt{Makie.jl},
        \texttt{LinearAlgebra}, \texttt{SparseArrays},
        \texttt{DifferentialEquations.jl}, \texttt{JuMP.jl} + \texttt{HiGHS.jl},
        \texttt{BenchmarkTools.jl}.
  \item \textbf{Código aberto industrial:} OpenFOAM (dispersão, escoamentos
        reativos), FDS --- \emph{Fire Dynamics Simulator} do NIST (incêndios),
        Cantera (cinética e propriedades de chama), ParaView (visualização).
  \item \textbf{Comerciais de referência (para conhecer):} FLACS (Gexcon),
        KFX (DNV), Ansys Fluent/CFX, PHAST (modelos integrais, DNV).
  \item \textbf{Gestão:} Git desde o primeiro dia; um \emph{diário de bordo}
        (Markdown ou caderno) com data, o que tentou, o que deu errado e o que
        aprendeu.
\end{itemize}

\section*{Estrutura de repositório sugerida}
\begin{lstlisting}[language={},numbers=none]
cfd-seguranca/
  Project.toml     # ambiente Julia
  README.md        # mapa e status
  diario/          # 1 entrada/sessao
  moduloJ_julia/
  modulo01_edp/    # codigo+figuras
  ...
  modulo16_projeto/
  lib/             # funcoes comuns
  dados/           # validacao
\end{lstlisting}

% ====================================================================== MAPA
\section*{Mapa do curso e cronograma sugerido}
O mapa completo está na tabela de página inteira a seguir. Cada linha traz a
pergunta de segurança que o módulo ajuda a responder --- volte a ela sempre
que perder o fio da meada.

\begin{table*}[!t]
\centering\small
\begin{tabularx}{\textwidth}{@{}clXcc@{}}
\toprule
Mód. & Tema & Pergunta de segurança que ele responde & Nível & Semanas\\
\midrule
0  & Preparação & Que física mata pessoas em acidentes de processo? & Básico & 1\\
J  & Oficina de Julia & Como transformar equações em código rápido e confiável? & Básico & 1--2\\
1  & EDPs e diferenças finitas & Posso confiar num número que saiu de um computador? & Básico & 1--2\\
2  & Advecção e difusão numérica & Por que minha pluma ``espalha'' demais? & Básico & 1--2\\
3  & Volumes finitos & Meu código conserva massa de gás vazado? & Básico & 2\\
4  & Navier--Stokes incompressível & Como o vento contorna equipamentos? & Interm. & 2--3\\
5  & Turbulência & Quanto o modelo de turbulência muda a nuvem? & Interm. & 2--3\\
6  & Camada limite atmosférica & Qual vento e estabilidade simular? & Interm. & 2\\
7  & Termodinâmica da fonte & Quanto gás sai, e como representar o jato? & Interm. & 2\\
8  & Jatos, gás denso, nuvens & Qual o volume inflamável equivalente? & Avanç. & 3\\
9  & Detecção de gás & Onde posicionar detectores? & Avanç. & 2\\
10 & Fundamentos de combustão & Quão reativa é esta mistura? & Interm. & 2\\
11 & Incêndios e radiação & Qual fluxo térmico chega ao alvo? & Avanç. & 3\\
12 & Explosões: fenomenologia & Que sobrepressão esperar (estimativa)? & Avanç. & 2--3\\
13 & CFD de explosões & Qual a carga acidental de projeto? & Avanç. & 3\\
14 & V\&V e incerteza & Quão errado posso estar? & Avanç. & 2\\
15 & Do CFD ao QRA & Quando o CFD vale o custo? & Avanç. & 1--2\\
16 & Projeto final integrado & Tudo junto, num módulo de FPSO & Síntese & 4--6\\
\bottomrule
\end{tabularx}
\par\smallskip{\small\textbf{Mapa do curso e cronograma sugerido.}}
\end{table*}

\begin{professor}
O cronograma assume de 6 a 10 horas semanais. Se você trabalha em tempo
integral, não se culpe por estender os módulos 4, 5 e 13: são os que mais
exigem maturação. O que eu \emph{não} aceito é pular os módulos 1 a 3. Já vi
engenheiros experientes defendendo resultados de dispersão em auditoria sem
saber explicar o que é difusão numérica --- e vi a credibilidade do estudo
inteiro desabar em cinco minutos.
\end{professor}

% ======================================================================
\part{Fundamentos numéricos}
% ======================================================================

\setcounter{chapter}{-1}
\chapter{Preparação: a física dos acidentes}

\objetivos
\begin{itemize}
  \item Montar o ambiente de trabalho e o hábito do diário de bordo.
  \item Reconhecer, em acidentes reais, a sequência
        \emph{liberação $\rightarrow$ dispersão $\rightarrow$ ignição
        $\rightarrow$ incêndio/explosão $\rightarrow$ escalonamento}.
  \item Mapear cada elo dessa cadeia a um módulo do curso.
\end{itemize}

\begin{conceito}{A cadeia de eventos de um acidente de processo}
Quase todo grande acidente com hidrocarbonetos segue o mesmo esqueleto:
\begin{enumerate}
  \item \textbf{Perda de contenção} (furo, ruptura, falha de flange) --- termodinâmica
        e escoamento compressível (Módulo 7).
  \item \textbf{Dispersão} do material, formando uma nuvem inflamável ou tóxica ---
        transporte turbulento, camada limite atmosférica, efeitos de
        flutuabilidade (Módulos 3--8).
  \item \textbf{Detecção e ação} (ou falha delas) --- posicionamento de
        detectores, tempos de resposta (Módulo 9).
  \item \textbf{Ignição} imediata (incêndio em jato, poça) ou retardada
        (\emph{flash fire}, explosão de nuvem) --- combustão (Módulo 10).
  \item \textbf{Efeitos}: radiação térmica (Módulo 11), sobrepressão
        (Módulos 12--13).
  \item \textbf{Escalonamento}: falha de outras barreiras, novos vazamentos,
        efeito dominó.
\end{enumerate}
A Análise Quantitativa de Risco (QRA) combina \emph{frequência} e
\emph{consequência} de muitos cenários dessa cadeia (Módulo 15).
\end{conceito}

\begin{tarefa}{Ambiente e primeiro gráfico}
\enunciado Instale Julia (ou Python), um pacote de gráficos, Git, ParaView e
--- em ambiente Linux ou contêiner --- OpenFOAM e FDS. Plote a solução
analítica da difusão em meio semi-infinito com temperatura imposta na face:
\[
  \frac{T(x,t)-T_0}{T_s-T_0} = \operatorname{erfc}\!\left(\frac{x}{2\sqrt{\alpha t}}\right)
\]
para quatro instantes de tempo.
\entregaveis
\begin{itemize}
  \item Repositório Git criado com a estrutura sugerida e um README.
  \item Script que gera a figura (eixos com unidades, legenda).
  \item Primeira entrada no diário: seus objetivos com o curso.
\end{itemize}
\criterio a figura é reproduzida por qualquer pessoa clonando o repositório
e executando um único comando.
\reflexao\ Por que a solução depende só da combinação $x/\sqrt{\alpha t}$
(variável de similaridade)? O que isso diz sobre como uma frente de difusão
avança no tempo?
\end{tarefa}

\begin{tarefa}{Autópsia de acidentes}
\enunciado Leia os resumos (e, se possível, trechos dos relatórios oficiais)
de \textbf{Piper Alpha} (1988) e \textbf{Buncefield} (2005). Para cada um,
preencha uma tabela com: material liberado, mecanismo de liberação, como a
nuvem se formou, fonte de ignição provável, tipo de efeito (jato de fogo,
poça, explosão, \emph{flash fire}) e que barreiras falharam.
\entregaveis
\begin{itemize}
  \item Tabela comparativa (1 página).
  \item Para cada linha da tabela, indique qual módulo deste curso trata a
        física correspondente.
\end{itemize}
\criterio você consegue explicar a sequência de cada acidente em dois minutos,
usando os termos técnicos corretos.
\reflexao\ Em Buncefield, a sobrepressão observada foi muito superior ao
que os métodos de avaliação da época previam para uma área considerada
``pouco congestionada''. O que isso ensina sobre as hipóteses embutidas em
modelos de consequência?
\end{tarefa}

\begin{critico}
\begin{itemize}
  \item Se CFD fosse perfeito, ainda precisaríamos de QRA? E se a QRA fosse
        perfeita, ainda precisaríamos de CFD?
  \item O que é mais perigoso: um modelo simples cujas limitações você
        conhece, ou um modelo sofisticado cujas limitações você ignora?
\end{itemize}
\end{critico}

\begin{chaves}
perda de contenção / \emph{loss of containment}; cenário acidental /
\emph{accident scenario}; escalonamento / \emph{escalation}; efeito dominó /
\emph{domino effect}; análise de consequências / \emph{consequence analysis};
QRA / \emph{quantitative risk assessment}.
\end{chaves}

\begin{refs}
\begin{itemize}
  \item Cullen, W. D. \emph{The Public Inquiry into the Piper Alpha Disaster}. HMSO, 1990.
  \item Buncefield Major Incident Investigation Board. \emph{The Buncefield Incident 11 December 2005 --- Final Report}, 2008.
  \item Crowl, D. A.; Louvar, J. F. \emph{Chemical Process Safety: Fundamentals with Applications}. 4.\,ed. Pearson, 2019.
  \item Mannan, S. (ed.). \emph{Lees' Loss Prevention in the Process Industries}. 4.\,ed. Butterworth-Heinemann, 2012.
\end{itemize}
\end{refs}

\begin{aprofundar}
\begin{itemize}[leftmargin=*]
  \item \textbf{Kletz, T.} \emph{What Went Wrong? Case Histories of Process Plant Disasters}. 5.\,ed. Butterworth-Heinemann, 2009. Centenas de casos curtos; leia um por dia durante o curso.
  \item \textbf{Hopkins, A.} \emph{Failure to Learn: The BP Texas City Refinery Disaster}. CCH, 2008. Mostra como fatores organizacionais anulam boas análises técnicas.
  \item \textbf{Vídeos do U.S. CSB} (\url{https://www.csb.gov}). Animações de investigações reais; ótimos para visualizar a cadeia liberação--ignição--efeito.
  \item \textbf{Atherton, J.; Gil, F.} \emph{Incidents That Define Process Safety}. CCPS/Wiley, 2008.
\end{itemize}
\end{aprofundar}

% ======================================================================
\renewcommand{\thechapter}{J}
\renewcommand{\theHchapter}{J}
\chapter{Oficina de Julia para iniciantes}

\begin{professor}
Antes de discretizar qualquer equação, você precisa de uma ferramenta em que
confie. Já vi alunos brilhantes em física travarem por semanas porque não
sabiam por que o código estava 100 vezes mais lento que o esperado, ou por
que um vetor ``mudou sozinho''. Esta oficina existe para que isso não
aconteça com você. Não pule: ela vai economizar dezenas de horas nos módulos
seguintes.
\end{professor}

\objetivos
\begin{itemize}
  \item Instalar Julia, usar o REPL e criar ambientes reprodutíveis.
  \item Dominar vetores, matrizes, \emph{broadcasting} e funções.
  \item Entender tipos e \emph{dispatch} múltiplo --- e usá-los para
        organizar esquemas numéricos.
  \item Escrever código rápido: estabilidade de tipo, pré-alocação,
        \emph{benchmarks}.
  \item Testar código numérico com \texttt{Test}.
\end{itemize}

\begin{conceito}{Por que Julia para CFD?}
O ``problema das duas linguagens'': prototipa-se em Python ou MATLAB
(fácil, lento em laços) e reescreve-se em C/Fortran (rápido, trabalhoso).
Julia compila cada função para código de máquina na primeira chamada
(\emph{JIT}, via LLVM), especializada para os tipos dos argumentos. Resultado:
\textbf{laços explícitos são rápidos}, e você escreve um solver de volumes
finitos ``como no livro'', sem truques de vetorização.

O preço: a primeira execução de cada função é mais lenta (tempo de
compilação), e alguns hábitos de Python levam a código lento em Julia.
\end{conceito}

\begin{dicionario}[Vocabulário de Julia]
\begin{itemize}[leftmargin=*,nosep]
  \item \textbf{REPL} --- terminal interativo (\emph{Read-Eval-Print Loop}).
  \item \textbf{JIT} --- compilação ``no momento'' (\emph{just-in-time}).
  \item \textbf{LLVM} --- infraestrutura de compiladores usada também por
        Clang (C/C++) e Rust.
  \item \textbf{Dispatch múltiplo} --- a função executada é escolhida pelos
        tipos de \emph{todos} os argumentos.
  \item \textbf{Estabilidade de tipo} --- o tipo do resultado depende só dos
        tipos das entradas, não dos valores. Condição para código rápido.
  \item \textbf{Broadcasting} --- aplicar uma operação elemento a elemento
        com o ponto: \texttt{f.(x)}, \texttt{a .+ b}.
\end{itemize}
\origem Julia foi criada no MIT por Jeff Bezanson, Stefan Karpinski, Viral
Shah e Alan Edelman, a partir de 2009; anunciada publicamente em 2012
(``\emph{Why We Created Julia}'') e estabilizada na versão 1.0 em 2018.
\end{dicionario}

\begin{conceito}{Instalação, REPL e modos}
Instale pelo \texttt{juliaup} (gerenciador oficial de versões). No REPL:
\begin{itemize}[nosep]
  \item \texttt{]} --- modo de pacotes (\texttt{add}, \texttt{status}, \texttt{activate}).
  \item \texttt{?} --- ajuda: \texttt{?sum} mostra a documentação.
  \item \texttt{;} --- terminal do sistema.
  \item \texttt{Backspace} --- volta ao modo normal.
\end{itemize}
Para scripts: \texttt{julia --project=. meu\_script.jl}. Editor recomendado:
VS Code com a extensão oficial de Julia. Para trabalho interativo, o pacote
\texttt{Revise.jl} recarrega o código ao salvar.
\end{conceito}

\begin{lstlisting}[caption={Primeiros passos: tipos, strings e vetores}]
x = 3            # Int64
y = 3.0          # Float64
typeof(x / 2)    # Float64: "/" sempre da real
div(7, 2)        # 3 (divisao inteira)
2^63             # ESTOURA em silencio (Int64)!
big(2)^63        # BigInt resolve

gas = "metano"
println("LFL do $gas = $(5.0) %vol")

u = [1.0, 2.0, 3.0]  # vetor
u[1], u[end]         # indice comeca em 1
xs = range(0, 1, length=11)
A = zeros(3, 3)      # matriz 3x3
linha = A[2, :]      # COPIA da linha 2
@views l2 = A[2, :]  # visao, sem copia
\end{lstlisting}

\begin{conceito}{Vetores, matrizes e \emph{broadcasting}}
\begin{itemize}
  \item Índices começam em \textbf{1} e o último é \texttt{end}.
  \item Matrizes são armazenadas \textbf{por coluna} (como Fortran):
        em laços duplos, o índice mais interno deve ser o \emph{primeiro}
        (\texttt{A[i,j]} com \texttt{i} no laço interno).
  \item \textbf{Fatiar copia}: \texttt{A[2,:]} aloca memória nova; use
        \texttt{@views} ou \texttt{view(A,2,:)} para evitar.
  \item O \textbf{ponto} aplica qualquer função elemento a elemento, e
        vários pontos numa expressão são \emph{fundidos} num único laço.
  \item \texttt{.=} atribui \emph{dentro} de um vetor existente (sem alocar).
\end{itemize}
\end{conceito}

\begin{lstlisting}[caption={Broadcasting e operacoes in-place}]
x = range(0, 1, length=101)
u = sin.(pi .* x)      # ponto = elemento a elemento
v = similar(u)         # aloca uma vez
v .= u .^ 2 .+ 1       # reusa v, laco unico
@. v = u^2 + 1         # idem, com a macro @.
\end{lstlisting}

\begin{lstlisting}[caption={Funcoes: forma longa, curta, argumentos nomeados e convencao "!"}]
function fluxo_critico(p0, T0, d; Cd=0.85,
                       gam=1.3, M=0.019)
    R = 8.314                  # J/(mol K)
    A = pi * d^2 / 4
    f = (2/(gam+1))^((gam+1)/(gam-1))
    return Cd * A * p0 * sqrt(gam*M/(R*T0) * f)
end

area(d) = pi * d^2 / 4         # forma curta

# "!" = a funcao MODIFICA o argumento
function zera_contorno!(u)
    u[1] = 0.0; u[end] = 0.0
    return u
end

mdot = fluxo_critico(80e5, 313.0, 0.025)
\end{lstlisting}

\begin{conceito}{Tipos e \emph{dispatch} múltiplo}
Em Julia, você define \emph{tipos} (com \texttt{struct}) e escreve várias
versões (\emph{métodos}) da mesma função para tipos diferentes. O compilador
escolhe a versão certa. Para CFD, isso organiza o código de forma natural:
um tipo para cada esquema, cada modelo de turbulência, cada condição de
contorno --- e o laço principal não muda. É o equivalente, mais flexível, a
classes e herança.
\end{conceito}

\begin{lstlisting}[caption={Dispatch: um tipo por esquema numerico}]
abstract type Esquema end
struct Upwind      <: Esquema end
struct LaxWendroff <: Esquema end

function passo!(u, un, C, ::Upwind)
    for j in 2:length(u)
        u[j] = un[j] - C*(un[j] - un[j-1])
    end
end

function passo!(u, un, C, ::LaxWendroff)
    for j in 2:length(u)-1
        u[j] = un[j] - C/2*(un[j+1] - un[j-1]) +
               C^2/2*(un[j+1] - 2un[j] + un[j-1])
    end
end

# o mesmo laco serve para qualquer esquema:
function advecta!(u, C, nt, esq::Esquema)
    un = similar(u)
    for n in 1:nt
        un .= u
        passo!(u, un, C, esq)
    end
    return u
end
# advecta!(u, 0.8, 200, Upwind())
\end{lstlisting}

\begin{conceito}{As sete regras do código rápido}
\begin{enumerate}
  \item Coloque o trabalho \textbf{dentro de funções}; evite variáveis
        globais (ou declare-as \texttt{const}).
  \item Mantenha \textbf{estabilidade de tipo}: não comece uma variável como
        \texttt{0} (Int) se ela vai virar \texttt{Float64}; use \texttt{0.0}.
        Diagnóstico: \texttt{@code\_warntype f(args)} (vermelho = problema).
  \item \textbf{Pré-aloque} vetores fora dos laços de tempo; use \texttt{.=}.
  \item Laços explícitos são rápidos --- não precisa ``vetorizar tudo''.
  \item Respeite a ordem de memória por coluna.
  \item Meça com \texttt{@btime} (pacote \texttt{BenchmarkTools}), não com
        a primeira execução, que inclui compilação. Observe as
        \textbf{alocações}: num laço de tempo bem escrito, devem ser zero.
  \item Só depois de correto e medido, use \texttt{@inbounds} (desliga
        verificação de índices) e paralelismo (\texttt{Threads.@threads}).
\end{enumerate}
\end{conceito}

\begin{lstlisting}[caption={Lento vs. rapido: o mesmo calculo}]
# LENTO: global nao tipada + aloca a cada passo
alfa = 0.4
function difunde_lento(u, n)
    for k in 1:n
        u = u .+ alfa .* [0; diff(diff(u)); 0]
    end
    return u
end

# RAPIDO: argumentos, pre-alocacao, laco
function difunde!(u, tmp, r, n)
    for k in 1:n
        tmp .= u
        @inbounds for j in 2:length(u)-1
            u[j] = tmp[j] + r*(tmp[j+1] -
                   2tmp[j] + tmp[j-1])
        end
    end
    return u
end

using BenchmarkTools
u = rand(10_000); tmp = similar(u)
@btime difunde_lento($u, 100)
@btime difunde!($u, $tmp, 0.4, 100)
\end{lstlisting}

\begin{conceito}{Ambientes, pacotes e testes}
Cada projeto tem um \texttt{Project.toml} (pacotes que você pediu) e um
\texttt{Manifest.toml} (versões exatas de tudo). Versionando ambos no Git,
qualquer pessoa reproduz seu ambiente com \texttt{] instantiate} ---
exatamente o que o critério ``reprodutível'' das tarefas exige.

Testes com o módulo padrão \texttt{Test}: \texttt{@test} verifica uma
condição; \texttt{@testset} agrupa. Em código numérico, compare com
tolerância (\texttt{isapprox}), nunca com \texttt{==}.
\end{conceito}

\begin{lstlisting}[caption={Ambiente e teste do TDMA}]
# terminal:  julia --project=.
# no REPL, tecle ] e depois:
#   add Plots BenchmarkTools JuMP HiGHS
using Test, LinearAlgebra
include("tdma.jl")   # sua funcao tdma(a,b,c,d)

@testset "TDMA vs. solver denso" begin
    n = 50
    a = rand(n); c = rand(n); d = rand(n)
    b = 4 .+ rand(n)   # diagonal dominante
    A = Tridiagonal(a[2:end], b, c[1:end-1])
    @test isapprox(tdma(a, b, c, d), A \ d;
                   rtol=1e-10)
end
\end{lstlisting}

\begin{center}\small
\begin{tabular}{@{}p{0.30\linewidth}p{0.30\linewidth}p{0.30\linewidth}@{}}
\toprule
Aspecto & Python/NumPy & Julia\\ \midrule
1.º índice & 0 & 1\\
Laço explícito & lento & rápido\\
Elemento a elemento & implícito & ponto: \texttt{.+}\\
Fatia & visão & cópia\\
Ordem de memória & por linha & por coluna\\
Potência & \texttt{**} & \texttt{\^{}}\\
\bottomrule
\end{tabular}
\end{center}

\begin{armadilha}
Três erros que vejo todo semestre: (1) \texttt{2\^{}63} retorna um número
negativo sem aviso --- inteiros têm tamanho fixo; (2) \texttt{v = u} não
copia: os dois nomes apontam para o mesmo vetor, e mexer em \texttt{v}
muda \texttt{u} (use \texttt{copy(u)}); (3) medir desempenho na primeira
chamada e concluir que ``Julia é lenta''.
\end{armadilha}

\begin{tarefa}{Ambiente reprodutível}
\enunciado Crie o repositório do curso como projeto Julia
(\texttt{] activate .}), adicione \texttt{Plots}, \texttt{BenchmarkTools} e
\texttt{SpecialFunctions} (para \texttt{erfc}) e refaça a figura da Tarefa
0.1 como script executável com \texttt{julia --project=.}.
\entregaveis\ \texttt{Project.toml} e \texttt{Manifest.toml} versionados;
script; README com o comando único de execução.
\criterio clonar em outra máquina, rodar \texttt{] instantiate} e o script
reproduz a figura.
\reflexao\ Por que fixar versões exatas de pacotes é uma questão de
\emph{segurança} num estudo de engenharia, e não só de conveniência?
\end{tarefa}

\begin{tarefa}{Seu primeiro solver: TDMA}
\enunciado Implemente \texttt{tdma(a, b, c, d)} (e a versão
\texttt{tdma!} sem alocações) para o sistema
$a_ix_{i-1}+b_ix_i+c_ix_{i+1}=d_i$. Teste contra
\texttt{Tridiagonal(...) \textbackslash\ d}. Meça o tempo para
$n=10^3,10^4,10^5,10^6$.
\entregaveis\ Código; \texttt{@testset} passando; gráfico tempo $\times n$
em escala log-log.
\criterio inclinação $\approx1$ (custo linear) e zero alocações na versão
\texttt{tdma!}.
\reflexao\ Por que a eliminação de Gauss genérica custaria $\mathcal O(n^3)$
e o TDMA custa $\mathcal O(n)$? Quando o TDMA pode falhar (dica:
dominância diagonal)?
\end{tarefa}

\begin{tarefa}{Caça às alocações}
\enunciado Use o par \texttt{difunde\_lento}/\texttt{difunde!} da
listagem acima. (a) Meça os dois com \texttt{@btime}. (b) Rode
\texttt{@code\_warntype difunde\_lento(u, 10)} e identifique a instabilidade
de tipo. (c) Corrija a versão lenta em três passos (global $\rightarrow$
argumento; pré-alocação; laço) medindo o ganho de cada um.
\entregaveis\ Tabela com tempo e alocações em cada etapa.
\criterio aceleração final de pelo menos uma ordem de grandeza.
\reflexao\ Qual mudança deu o maior ganho? Escreva a regra com suas
palavras no diário --- ela valerá para todos os solvers do curso.
\end{tarefa}

\begin{tarefa}{Dispatch para esquemas numéricos}
\enunciado Partindo da listagem de dispatch, adicione os tipos
\texttt{Central} e \texttt{TVD\{L\}}, em que \texttt{L} é um tipo de limitador
(\texttt{Minmod}, \texttt{VanLeer}, \texttt{Superbee}) e a função
\texttt{psi(r, ::Minmod)} etc.\ define cada um. Escreva um único laço de
simulação que funcione com qualquer combinação.
\entregaveis\ Módulo \texttt{Esquemas.jl} com testes (ex.: com $C=1$, o
upwind transporta o perfil exatamente uma célula).
\criterio adicionar um novo limitador exige escrever \emph{apenas} uma
função \texttt{psi}.
\reflexao\ Este é o esqueleto do ``torneio'' do Módulo 2. Compare com como
você organizaria isso em Python (classes) ou Rust (\emph{traits}).
\end{tarefa}

\begin{critico}
\begin{itemize}
  \item Em que etapas deste curso Julia é a melhor escolha, e em quais você
        usaria Python (ecossistema, ex.: Cantera) ou Rust (ferramenta
        distribuível, segurança de memória)?
  \item Um código rápido e errado é pior que um código lento e certo. Que
        disciplina de testes você vai adotar antes de otimizar?
\end{itemize}
\end{critico}

\begin{chaves}
REPL; compilação JIT / \emph{just-in-time compilation}; dispatch múltiplo /
\emph{multiple dispatch}; estabilidade de tipo / \emph{type stability};
broadcasting; alocação / \emph{allocation}; pré-alocação /
\emph{preallocation}; ordem por coluna / \emph{column-major}; ambiente /
\emph{environment}; teste unitário / \emph{unit test}.
\end{chaves}

\begin{refs}
\begin{itemize}
  \item Bezanson, J.; Edelman, A.; Karpinski, S.; Shah, V. B. Julia: a fresh approach to numerical computing. \emph{SIAM Rev.}, 59(1):65--98, 2017.
  \item Documentação oficial: \url{https://docs.julialang.org} --- em especial ``\emph{Performance Tips}'' e ``\emph{Noteworthy Differences from other Languages}''.
\end{itemize}
\end{refs}

\begin{aprofundar}
\begin{itemize}[leftmargin=*]
  \item \textbf{Lauwens, B.; Downey, A.} \emph{Think Julia}. O'Reilly, 2019 (versão gratuita on-line). Para quem nunca programou de forma estruturada.
  \item \textbf{Storopoli, J.; Huijzer, R.; Alonso, L.} \emph{Julia Data Science}, 2021 (gratuito on-line). Rápido e prático, escrito em parte por brasileiros.
  \item \textbf{Rackauckas, C.} \emph{Parallel Computing and Scientific Machine Learning} (notas do curso MIT 18.337, gratuitas). Desempenho, paralelismo e EDPs em Julia --- leitura obrigatória depois do Módulo 4.
  \item \textbf{Kochenderfer, M. J.; Wheeler, T. A.} \emph{Algorithms for Optimization}. MIT Press, 2019. Otimização com exemplos em Julia; útil no Módulo 9.
  \item \textbf{Seção ``Performance Tips'' do manual}: releia a cada novo solver.
\end{itemize}
\end{aprofundar}

\renewcommand{\thechapter}{\arabic{chapter}}
\renewcommand{\theHchapter}{\arabic{chapter}}
\setcounter{chapter}{0}

% ======================================================================
\chapter{Equações diferenciais parciais e diferenças finitas}

\objetivos
\begin{itemize}
  \item Classificar EDPs e entender o que cada tipo exige de um método numérico.
  \item Derivar aproximações por diferenças finitas e seu erro de truncamento.
  \item Aplicar os conceitos de consistência, estabilidade e convergência
        (teorema de Lax) e a análise de von Neumann.
  \item Comparar esquemas explícitos e implícitos na prática.
\end{itemize}

\begin{conceito}{Classificação de EDPs de segunda ordem}
Para $A u_{xx} + B u_{xy} + C u_{yy} + \ldots = 0$, o discriminante
$B^2-4AC$ define o tipo:
\begin{itemize}
  \item $<0$: \textbf{elíptica} (ex.: Poisson, $\nabla^2 p = f$). Informação
        se propaga instantaneamente em todas as direções; precisa de
        condições em todo o contorno. É a equação de pressão do Módulo 4.
  \item $=0$: \textbf{parabólica} (ex.: difusão, $u_t = \alpha u_{xx}$).
        Marcha no tempo com suavização.
  \item $>0$: \textbf{hiperbólica} (ex.: ondas, advecção pura, Euler).
        Informação viaja por características com velocidade finita ---
        a física das ondas de choque das explosões (Módulos 12--13).
\end{itemize}
As equações de Navier--Stokes têm caráter misto, e é por isso que os
algoritmos de CFD são combinações engenhosas de técnicas para cada parte.
\end{conceito}

\begin{conceito}{Diferenças finitas e erro de truncamento}
Expandindo em série de Taylor, $u_{i\pm1} = u_i \pm \Delta x\, u_x +
\tfrac{\Delta x^2}{2} u_{xx} \pm \tfrac{\Delta x^3}{6}u_{xxx} + \ldots$, obtém-se
\begin{align*}
  \frac{u_{i+1}-2u_i+u_{i-1}}{\Delta x^2}
  &= u_{xx} + \underbrace{\frac{\Delta x^2}{12}u_{xxxx} + \ldots}_{\tau\;=\;\mathcal{O}(\Delta x^2)}
\end{align*}
onde $\tau$ é o \emph{erro de truncamento local}.
A \textbf{ordem} do esquema é o expoente do termo líder do erro. Ordem alta
não garante precisão em malha grossa: o coeficiente e as derivadas altas
da solução importam.
\end{conceito}

\begin{dicionario}[Dicionário da equação: série de Taylor]
\begin{itemize}[leftmargin=*,nosep]
  \item $u_i$ --- valor discreto da variável no nó $i$, em $x_i=i\,\Delta x$.
  \item $\Delta x$ --- espaçamento da malha [m].
  \item $u_{xx}$, $u_{xxxx}$ --- derivadas parciais de 2.ª e 4.ª ordem em $x$.
  \item $\mathcal O(\Delta x^2)$ --- ``da ordem de $\Delta x^2$'': o erro cai
        4 vezes quando $\Delta x$ cai à metade.
\end{itemize}
\origem A série é de Brook Taylor (1715). A notação $\mathcal O(\cdot)$
(``O grande'') foi introduzida por Paul Bachmann (1894) e popularizada por
Edmund Landau. A ideia de substituir derivadas por quocientes de diferenças
já aparece em Euler (séc.\ XVIII); a primeira solução numérica sistemática de
EDPs por diferenças finitas é de L.\,F. Richardson (1911).
\end{dicionario}


\begin{conceito}{Consistência, estabilidade, convergência}
\begin{itemize}
  \item \textbf{Consistência:} o erro de truncamento $\rightarrow 0$ quando
        $\Delta x, \Delta t \rightarrow 0$.
  \item \textbf{Estabilidade:} erros (de arredondamento, de condição inicial)
        não crescem sem limite.
  \item \textbf{Convergência:} a solução discreta tende à solução exata.
\end{itemize}
\textbf{Teorema da equivalência de Lax:} para um problema linear bem-posto, um
esquema consistente é convergente \emph{se e somente se} é estável.
\medskip

\textbf{Análise de von Neumann.} Substitua $u_j^n = G^n e^{\mathrm{i}kj\Delta x}$
no esquema. Estabilidade exige $|G|\le 1$ para todo número de onda $k$.
Para FTCS na difusão, com $r=\alpha\Delta t/\Delta x^2$:
\begin{gather*}
  u_j^{n+1} = u_j^n + r\,(u_{j+1}^n - 2u_j^n + u_{j-1}^n)\\
  \Downarrow\\
  G = 1 - 4r\sin^2\!\left(\tfrac{k\Delta x}{2}\right)
  \;\Rightarrow\; r \le \tfrac12 .
\end{gather*}
\end{conceito}

\begin{dicionario}[Dicionário da equação: difusão e von Neumann]
\begin{itemize}[leftmargin=*,nosep]
  \item $u_t=\alpha u_{xx}$ --- equação do calor/difusão; $\alpha$ é a
        difusividade [m$^2$/s] (térmica, mássica ou turbulenta).
  \item $u_j^n$ --- valor no nó $j$ e no instante $t_n=n\,\Delta t$.
  \item $r=\alpha\Delta t/\Delta x^2$ --- número de difusão da malha
        (ou ``número de Fourier da malha''), adimensional.
  \item $k$ --- número de onda [rad/m] de um modo de Fourier do erro;
        $k\Delta x\in[0,\pi]$ vai do modo mais suave ao ``serrote''.
  \item $G$ --- fator de amplificação: quanto um modo é multiplicado a
        cada passo. $|G|>1$ $\Rightarrow$ crescimento exponencial.
  \item $\mathrm i=\sqrt{-1}$.
\end{itemize}
\origem A equação do calor é de Joseph Fourier (\emph{Théorie analytique de
la chaleur}, 1822). A análise de estabilidade foi desenvolvida por John von
Neumann em Los Alamos nos anos 1940 e divulgada por Crank \& Nicolson (1947)
e O'Brien, Hyman \& Kaplan (1951). A ideia: como o esquema é linear, o erro é
uma soma de senoides; basta ver se \emph{alguma} delas cresce. FTCS =
\emph{Forward Time, Centered Space}. O teorema de equivalência é de Lax \&
Richtmyer (1956).
\end{dicionario}


\begin{conceito}{Explícito, implícito e Crank--Nicolson}
\begin{itemize}
  \item \textbf{Explícito (FTCS):} barato por passo, estabilidade condicional.
  \item \textbf{Implícito (BTCS):} resolve um sistema linear por passo
        (tridiagonal em 1D, algoritmo de Thomas/TDMA), incondicionalmente
        estável, primeira ordem no tempo.
  \item \textbf{Crank--Nicolson:} média dos dois, segunda ordem no tempo,
        incondicionalmente estável, mas \emph{não amortece} modos de alta
        frequência quando $r$ é grande --- surgem oscilações espúrias.
\end{itemize}
\end{conceito}

\begin{dicionario}[Dicionário da equação: o método $\theta$]
Os três esquemas são casos de uma única fórmula. Sendo $\mathcal L u_j =
u_{j+1}-2u_j+u_{j-1}$:
\[
  u_j^{n+1}-u_j^n = r\left[\theta\,\mathcal L u_j^{n+1}+(1-\theta)\,\mathcal L u_j^{n}\right],
\]
com $\theta=0$ (FTCS, explícito), $\theta=\tfrac12$ (Crank--Nicolson),
$\theta=1$ (BTCS, \emph{Backward Time, Centered Space}). Para $\theta>0$
surge um sistema tridiagonal $a_ju_{j-1}+b_ju_j+c_ju_{j+1}=d_j$.
\origem John Crank e Phyllis Nicolson publicaram o método em 1947, estudando
aquecimento de materiais com reação química. O algoritmo tridiagonal
(eliminação de Gauss especializada, custo $\mathcal O(N)$) é atribuído a
Llewellyn Thomas (1949), daí ``algoritmo de Thomas'' ou TDMA
(\emph{TriDiagonal Matrix Algorithm}).
\end{dicionario}


\begin{lstlisting}[caption={Esqueleto FTCS em Julia (difusao 1D)}]
function ftcs!(u, alpha, dx, dt, nsteps)
    r = alpha*dt/dx^2
    r > 0.5 && @warn "r = $r > 0.5: espere instabilidade"
    un = similar(u)
    for n in 1:nsteps
        un .= u
        for i in 2:length(u)-1
            u[i] = un[i] + r*(un[i+1] - 2un[i] + un[i-1])
        end
        # contornos de Dirichlet: u[1], u[end] mantidos
    end
    return u
end
\end{lstlisting}

\begin{tarefa}{Difusão 1D: três esquemas, um veredito}
\enunciado Resolva $u_t=\alpha u_{xx}$ em $x\in[0,1]$, com $u(0,t)=u(1,t)=0$
e $u(x,0)=\sin(\pi x)$, cuja solução exata é $u=e^{-\alpha\pi^2 t}\sin(\pi x)$.
Implemente FTCS, BTCS (com TDMA escrito por você) e Crank--Nicolson.
\begin{enumerate}
  \item Rode FTCS com $r=0{,}40$; $0{,}50$; $0{,}51$. Documente o que acontece.
  \item Para cada esquema, faça um estudo de convergência: erro na norma $L_2$
        em $t=0{,}1$ versus $\Delta x$ (mantendo $r$ fixo), em escala log-log.
  \item Repita Crank--Nicolson com condição inicial degrau (descontínua) e
        $r=5$. Observe as oscilações.
\end{enumerate}
\entregaveis
\begin{itemize}
  \item Código dos três esquemas e do TDMA, com teste unitário do TDMA.
  \item Gráfico log-log com retas de referência de inclinação 1 e 2.
  \item Tabela: ordem observada $p = \ln(e_1/e_2)/\ln(\Delta x_1/\Delta x_2)$.
  \item Figura das oscilações de Crank--Nicolson.
\end{itemize}
\criterio as ordens observadas coincidem com as teóricas (com tolerância de
$\pm 0{,}1$) na faixa assintótica.
\reflexao\ Por que, mantendo $r$ fixo, o FTCS aparenta ser de segunda ordem
em $\Delta x$? Qual o custo computacional de dobrar a resolução no FTCS
versus no BTCS?
\end{tarefa}

\begin{tarefa}{Von Neumann no papel}
\enunciado Derive à mão o fator de amplificação $G$ do BTCS e do
Crank--Nicolson para a equação de difusão. Mostre que $|G|\le1$ para todo
$r>0$ e determine o limite de $G$ do Crank--Nicolson quando $r\rightarrow\infty$
para o modo de maior frequência ($k\Delta x=\pi$).
\entregaveis\ Derivação digitalizada ou em \LaTeX; gráfico de $G(k\Delta x)$
para os três esquemas e $r\in\{0{,}25;\,1;\,10\}$.
\criterio o gráfico explica, sozinho, as oscilações vistas na Tarefa anterior.
\reflexao\ ``Incondicionalmente estável'' significa ``sempre preciso''? Dê um
contraexemplo com base nos seus resultados.
\end{tarefa}

\begin{professor}
Faço todo aluno ver o FTCS explodir pelo menos uma vez. Parece infantil, mas
a imagem daquele serrote crescendo exponencialmente é o que vai fazer você,
daqui a dois anos, desconfiar imediatamente quando um campo de pressão de
explosão aparecer com oscilações de malha a malha. Instabilidade numérica tem
``cara'', e você precisa reconhecê-la.
\end{professor}

\begin{critico}
\begin{itemize}
  \item O teorema de Lax vale para problemas lineares. As equações de
        Navier--Stokes não são lineares. O que, então, garante a convergência
        dos códigos industriais? (Dica: pense em ``evidência'' em vez de
        ``prova''.)
  \item Uma solução estável e convergente na malha pode estar errada? Em
        que sentido?
\end{itemize}
\end{critico}

\begin{chaves}
equação diferencial parcial / \emph{PDE}; erro de truncamento /
\emph{truncation error}; consistência / \emph{consistency}; estabilidade de
von Neumann / \emph{von Neumann stability}; teorema de equivalência de Lax /
\emph{Lax equivalence theorem}; algoritmo de Thomas / \emph{TDMA};
Crank--Nicolson; ordem de convergência / \emph{order of accuracy}.
\end{chaves}

\begin{refs}
\begin{itemize}
  \item LeVeque, R. J. \emph{Finite Difference Methods for Ordinary and Partial Differential Equations}. SIAM, 2007.
  \item Anderson, J. D. \emph{Computational Fluid Dynamics: The Basics with Applications}. McGraw-Hill, 1995. (Caps. 3--4)
  \item Lax, P. D.; Richtmyer, R. D. Survey of the stability of linear finite difference equations. \emph{Comm. Pure Appl. Math.}, 9:267--293, 1956.
  \item Barba, L. A.; Forsyth, G. F. CFD Python: the 12 steps to Navier--Stokes equations. \emph{J. Open Source Educ.}, 1(9):21, 2018.
\end{itemize}
\end{refs}

\begin{aprofundar}
\begin{itemize}[leftmargin=*]
  \item \textbf{LeVeque (2007)}, caps.\ 9 e 10: a melhor exposição moderna de estabilidade para difusão e advecção, com ênfase em intuição.
  \item \textbf{Strikwerda, J. C.} \emph{Finite Difference Schemes and Partial Differential Equations}. 2.\,ed. SIAM, 2004. Para o rigor da análise de von Neumann e do teorema de Lax.
  \item \textbf{Trefethen, L. N.} \emph{Finite Difference and Spectral Methods for Ordinary and Partial Differential Equations} (notas, 1996, disponíveis gratuitamente na página do autor). O capítulo sobre velocidade de grupo e dispersão numérica é uma joia.
  \item \textbf{Crank \& Nicolson (1947)}, artigo original: leia para ver como um método clássico nasceu de um problema prático.
\end{itemize}
\end{aprofundar}

% ======================================================================
\chapter{Advecção e o dilema da difusão numérica}

\objetivos
\begin{itemize}
  \item Entender por que discretizar advecção é o problema central de CFD.
  \item Derivar e interpretar a \emph{equação modificada}.
  \item Conhecer o teorema de Godunov e os esquemas TVD com limitadores.
  \item Relacionar difusão numérica com erros em concentrações de gás.
\end{itemize}

\begin{conceito}{Advecção linear e a condição CFL}
$u_t + a\,u_x = 0$ transporta o perfil inicial sem deformação com velocidade
$a$. Para esquemas explícitos, a condição de Courant--Friedrichs--Lewy exige
que o domínio de dependência numérico contenha o físico:
\[
  C = \frac{|a|\,\Delta t}{\Delta x} \le 1 \quad\text{(tipicamente).}
\]
Em 3D com velocidades altas (jatos sônicos!), a CFL dita o passo de tempo e,
portanto, o custo da simulação.
\end{conceito}

\begin{dicionario}[Dicionário da equação: CFL]
\begin{itemize}[leftmargin=*,nosep]
  \item $a$ --- velocidade de advecção [m/s]; $C$ --- número de Courant.
  \item \emph{Domínio de dependência}: o conjunto de pontos do instante
        anterior que influenciam um ponto. Fisicamente, a informação anda
        $a\,\Delta t$; o esquema explícito ``enxerga'' só $\Delta x$.
  \item Em 3D: $C=\Delta t\,(|u|/\Delta x+|v|/\Delta y+|w|/\Delta z)$.
  \item Em escoamento compressível (explosões), troque $|u|$ por
        $|u|+c$, com $c$ a velocidade do som: as ondas de pressão também
        transportam informação.
\end{itemize}
\origem Richard Courant, Kurt Friedrichs e Hans Lewy (\emph{Math. Annalen},
1928). Curiosamente, o artigo não tratava de computação: usava diferenças
finitas para \emph{provar a existência} de soluções de EDPs. A condição
virou critério prático com os primeiros computadores, nos anos 1940--50.
\end{dicionario}


\begin{conceito}{Esquemas clássicos e a equação modificada}
\begin{itemize}
  \item \textbf{FTCS} (central no espaço): incondicionalmente \emph{instável}.
  \item \textbf{Upwind} de 1.ª ordem ($a>0$):
        $u_j^{n+1} = u_j^n - C(u_j^n - u_{j-1}^n)$. Estável para $C\le1$.
  \item \textbf{Lax--Wendroff}: 2.ª ordem, estável para $C\le1$, mas gera
        oscilações perto de descontinuidades.
\end{itemize}
A \textbf{equação modificada} é a EDP que o esquema resolve \emph{de fato}.
Para o upwind:
\[
  u_t + a u_x = \underbrace{\frac{a\Delta x}{2}(1-C)}_{\nu_{\text{num}}}\,u_{xx} + \ldots
\]
O upwind resolve uma equação de \emph{advecção-difusão}, com uma viscosidade
artificial $\nu_{\text{num}}$. Esquemas de 2.ª ordem têm termo líder com
$u_{xxx}$: erro \emph{dispersivo} (ondulações), não difusivo.
\end{conceito}

\begin{dicionario}[Dicionário da equação: equação modificada]
\begin{itemize}[leftmargin=*,nosep]
  \item $\nu_{\text{num}}$ --- viscosidade (ou difusividade) numérica
        [m$^2$/s], criada pelo esquema, não pela física.
  \item Termo em $u_{xx}$ --- \emph{difusivo}: achata e alarga picos.
  \item Termo em $u_{xxx}$ --- \emph{dispersivo}: ondas de comprimentos
        diferentes viajam a velocidades diferentes, gerando ondulações.
\end{itemize}
\textbf{Como se deriva:} expanda em Taylor cada termo do esquema discreto em
torno de $(x_j,t_n)$, e use a própria EDP para trocar derivadas temporais
por espaciais ($u_{tt}=a^2u_{xx}$). O que sobra além de $u_t+au_x$ é o erro.
\origem Warming \& Hyett (\emph{J.\ Comput.\ Phys.}, 1974) sistematizaram o
método; a noção de ``viscosidade artificial'' é anterior e intencional em
von Neumann \& Richtmyer (1950), que a \emph{adicionavam} para capturar choques.
\end{dicionario}


\begin{conceito}{Teorema de Godunov e esquemas TVD}
\textbf{Godunov (1959):} esquemas \emph{lineares} que preservam monotonicidade
são, no máximo, de 1.ª ordem. A saída é usar esquemas \emph{não lineares}:
a variação total $TV=\sum_j|u_{j+1}-u_j|$ não deve crescer (TVD). Isso se faz
com \textbf{limitadores de fluxo} $\psi(r)$, com $r$ a razão de gradientes
sucessivos:
\begin{gather*}
  \text{minmod: } \psi=\max(0,\min(1,r)),\\
  \text{van Leer: } \psi=\frac{r+|r|}{1+|r|},\\
  \text{superbee: } \psi=\max(0,\min(2r,1),\min(r,2)).
\end{gather*}
O diagrama de Sweby mostra a região TVD de 2.ª ordem. No OpenFOAM, você verá
isso como \texttt{limitedLinear}, \texttt{vanLeer}, \texttt{linearUpwind} etc.
\end{conceito}

\begin{dicionario}[Dicionário da equação: limitadores]
\begin{itemize}[leftmargin=*,nosep]
  \item $r_j=\dfrac{u_j-u_{j-1}}{u_{j+1}-u_j}$ --- razão de gradientes
        consecutivos: $r\approx1$ em regiões suaves, $r\le0$ em extremos.
  \item $\psi(r)$ --- limitador: mistura um fluxo de baixa ordem (robusto)
        com um de alta ordem (preciso):
        $F_{j+\frac12}=F^{\text{baixa}}+\psi(r_j)\,(F^{\text{alta}}-F^{\text{baixa}})$.
  \item $\psi=0$ $\Rightarrow$ upwind puro; $\psi=1$ $\Rightarrow$ Lax--Wendroff.
  \item $TV$ --- variação total; se não cresce, não nascem novos extremos.
\end{itemize}
\origem Sergei Godunov provou o teorema em 1959 (\emph{Mat.\ Sbornik}). Bram
van Leer (série ``\emph{Towards the ultimate conservative difference scheme}'',
1973--79) e Ami Harten (1983) criaram a teoria TVD; o \emph{superbee} é de Roe
(1985) e o diagrama de Sweby, de 1984.
\end{dicionario}


\begin{tarefa}{O torneio dos esquemas de advecção}
\enunciado Em domínio periódico $[0,1]$ com $a=1$, transporte (i) um pulso
gaussiano e (ii) um pulso quadrado por uma volta completa. Compare upwind,
Lax--Wendroff, e um esquema TVD com minmod, van Leer e superbee. Use
$C\in\{0{,}5;\,0{,}9;\,1{,}0\}$ e 100, 200 e 400 células.
\entregaveis
\begin{itemize}
  \item Figuras sobrepostas: solução exata vs. cada esquema.
  \item Tabela com: erro $L_1$, razão pico numérico/pico exato, $\max u$ e
        $\min u$ (detecta oscilações), variação total final.
  \item Explicação de por que $C=1$ dá resultado ``perfeito'' no upwind e por
        que isso é irrelevante na prática.
\end{itemize}
\criterio você consegue prever, olhando apenas a tabela, qual esquema
escolheria para transportar concentração de gás.
\reflexao\ Concentração de gás é uma grandeza limitada ($0\le Y\le1$). Que
esquemas violaram esses limites? Que consequência física absurda isso teria
num cálculo de inflamabilidade?
\end{tarefa}

\begin{tarefa}{Medindo a viscosidade numérica}
\enunciado Transporte um pulso gaussiano com upwind e meça, ao longo do tempo,
a variância $\sigma^2(t)$ do perfil. Sabendo que a difusão pura faz
$\sigma^2(t)=\sigma_0^2+2\nu t$, estime $\nu_{\text{num}}$ por regressão e
compare com $\frac{a\Delta x}{2}(1-C)$.
\entregaveis\ Gráfico $\sigma^2\times t$, valor estimado vs. teórico para
três malhas.
\criterio diferença menor que 5\% entre estimado e teórico.
\reflexao\ Num jato de gás com $u\sim 50$ m/s em células de 0,5 m, qual seria
$\nu_{\text{num}}$ de um upwind de 1.ª ordem? Compare com a viscosidade
turbulenta típica ($\sim 10^{-2}$ a $1$ m$^2$/s). Quem domina a mistura na
sua simulação: a física ou o esquema?
\end{tarefa}

\begin{armadilha}
Usar upwind de 1.ª ordem ``para convergir'' e esquecer de trocar depois. Em
dispersão, isso achata picos de concentração e altera o tamanho da nuvem
inflamável. Em explosões, espalha a frente de chama e reduz a sobrepressão.
Em ambos, o erro não é conservador por definição: pode ir para qualquer lado.
\end{armadilha}

\begin{critico}
\begin{itemize}
  \item A difusão numérica aumenta ou diminui a distância até o LFL? (Resposta
        honesta: depende. Construa um argumento para cada caso.)
  \item Se o modelo de turbulência já é uma aproximação grosseira, vale a
        pena insistir em esquemas de alta ordem?
\end{itemize}
\end{critico}

\begin{chaves}
advecção / \emph{advection}; número de Courant / \emph{Courant number, CFL};
difusão numérica / \emph{numerical diffusion}; dispersão numérica /
\emph{numerical dispersion}; equação modificada / \emph{modified equation};
limitador de fluxo / \emph{flux limiter}; TVD / \emph{total variation
diminishing}; monotonicidade / \emph{monotonicity}.
\end{chaves}

\begin{refs}
\begin{itemize}
  \item LeVeque, R. J. \emph{Finite Volume Methods for Hyperbolic Problems}. Cambridge University Press, 2002. (Caps. 4, 6)
  \item Sweby, P. K. High resolution schemes using flux limiters for hyperbolic conservation laws. \emph{SIAM J. Numer. Anal.}, 21(5):995--1011, 1984.
  \item Harten, A. High resolution schemes for hyperbolic conservation laws. \emph{J. Comput. Phys.}, 49:357--393, 1983.
  \item Hirsch, C. \emph{Numerical Computation of Internal and External Flows}. 2.\,ed. Butterworth-Heinemann, 2007.
\end{itemize}
\end{refs}

\begin{aprofundar}
\begin{itemize}[leftmargin=*]
  \item \textbf{LeVeque (2002)}, caps.\ 4--6: volumes finitos, limitadores e TVD com rara clareza.
  \item \textbf{Laney, C. B.} \emph{Computational Gasdynamics}. Cambridge, 1998. Excelente ponte entre advecção escalar e Euler.
  \item \textbf{Warming \& Hyett (1974)}, \emph{J.\ Comput.\ Phys.} 14:159--179: o artigo da equação modificada; curto e transformador.
  \item \textbf{Toro (2009)}, cap.\ 13: limitadores TVD em detalhe, preparando o Módulo 12.
\end{itemize}
\end{aprofundar}

% ======================================================================
\chapter{Volumes finitos e a equação de transporte genérica}

\objetivos
\begin{itemize}
  \item Derivar a formulação de volumes finitos a partir da forma integral.
  \item Discretizar a equação geral de transporte com diversos esquemas de face.
  \item Garantir conservação, limitação e convergência de solvers iterativos.
  \item Construir um protótipo de dispersão de escalar passivo.
\end{itemize}

\begin{conceito}{A equação geral de transporte}
Massa de espécie, momento, energia, $k$, $\varepsilon$, fração de mistura ---
quase tudo em CFD tem a forma
\begin{align*}
  \underbrace{\frac{\partial(\rho\phi)}{\partial t}}_{\text{acúmulo}}
  + \underbrace{\nabla\cdot(\rho\mathbf{u}\phi)}_{\text{advecção}}
  &= \underbrace{\nabla\cdot(\Gamma\nabla\phi)}_{\text{difusão}}\\
  &\quad + \underbrace{S_\phi}_{\text{fonte}} .
\end{align*}
Integrando num volume de controle $V_P$ e usando o teorema de Gauss:
\begin{align*}
  \frac{d}{dt}\int_{V_P}\rho\phi\,dV &+ \sum_{f}(\rho\mathbf{u}\phi)_f\cdot\mathbf{S}_f\\
  &= \sum_f(\Gamma\nabla\phi)_f\cdot\mathbf{S}_f + \bar S_\phi V_P .
\end{align*}
O que sai por uma face entra na vizinha: a \textbf{conservação é exata no
nível discreto}. É por isso que FV domina os códigos industriais (OpenFOAM,
Fluent, FLACS, KFX, FDS usa formulação equivalente).
\end{conceito}

\begin{dicionario}[Dicionário da equação: transporte genérico]
\begin{itemize}[leftmargin=*,nosep]
  \item $\phi$ --- grandeza transportada por unidade de massa.
  \item $\rho$ --- massa específica [kg/m$^3$]; $\mathbf u$ --- velocidade [m/s].
  \item $\Gamma$ --- coeficiente de difusão [kg/(m$\cdot$s)].
  \item $S_\phi$ --- fonte por unidade de volume (reação, empuxo\ldots).
  \item $V_P$ --- volume da célula $P$; $\mathbf S_f$ --- vetor-área da
        face $f$ (normal externa $\times$ área) [m$^2$].
\end{itemize}
\begin{center}\small
\begin{tabular}{@{}lll@{}}
\toprule
$\phi$ & $\Gamma$ & equação \\ \midrule
$1$ & --- & continuidade\\
$u_i$ & $\mu$ & momento ($S$ inclui $-\nabla p$, $\rho\mathbf g$)\\
$Y_k$ & $\rho D_k$ & espécie química\\
$h$ & $\lambda/c_p$ & energia (entalpia)\\
\bottomrule
\end{tabular}
\end{center}
\origem O passo-chave é o \emph{teorema da divergência} (Gauss, 1813;
Ostrogradsky, 1826): a integral de um divergente no volume é o fluxo pela
superfície. O método de volumes finitos foi popularizado por Spalding,
Patankar e o grupo do Imperial College nos anos 1970.
\end{dicionario}


\begin{conceito}{Número de Péclet de célula e limitação}
Em 1D com esquema central, a forma discreta $a_P\phi_P=a_E\phi_E+a_W\phi_W+b$
tem $a_E = D_e - F_e/2$, com $F=\rho u$ e $D=\Gamma/\Delta x$. Se o
\textbf{Péclet de célula} $Pe=F/D=\rho u\Delta x/\Gamma > 2$, então $a_E<0$
e surgem oscilações não físicas. Regras de Patankar para soluções limitadas:
coeficientes positivos, $a_P \ge \sum a_{nb}$ (critério de Scarborough), e
linearização da fonte $S = S_C + S_P\phi_P$ com $S_P\le0$.
Esquemas: central, upwind, híbrido, lei de potência, QUICK, TVD.
\end{conceito}

\begin{dicionario}[Dicionário da equação: coeficientes de FV]
\begin{itemize}[leftmargin=*,nosep]
  \item $F=\rho u$ --- fluxo convectivo de massa [kg/(m$^2\cdot$s)].
  \item $D=\Gamma/\Delta x$ --- ``condutância'' difusiva, mesma unidade.
  \item $a_P,a_E,a_W$ --- coeficientes da célula e das vizinhas leste
        (\emph{East}) e oeste (\emph{West}); $b$ --- termo independente.
  \item $Pe=F/D$ --- razão entre transporte por advecção e por difusão
        na escala da célula.
  \item $S_C$, $S_P$ --- partes constante e linear da fonte.
\end{itemize}
\origem O número leva o nome de Jean Claude Eugène Péclet (1793--1857),
físico francês que estudou transferência de calor. As ``quatro regras
básicas'' de limitação são do livro de Patankar (1980); o esquema híbrido é
de Spalding (1972) e o QUICK, de Leonard (1979).
\end{dicionario}


\begin{conceito}{Condições de contorno e solvers lineares}
\textbf{Contornos:} Dirichlet (valor), Neumann (fluxo), Robin (misto),
saída (gradiente nulo --- cuidado com refluxo!). \textbf{Solvers:} Jacobi,
Gauss--Seidel, SOR, gradiente conjugado (matrizes simétricas, ex.: pressão),
BiCGStab (não simétricas, ex.: advecção), multigrid (o que torna 3D viável).
\textbf{Sub-relaxação:} $\phi^{\text{novo}} = \phi^{\text{velho}} +
\lambda(\phi^{\text{calc}}-\phi^{\text{velho}})$. \textbf{Resíduos}:
normalize-os, ou o número ``$10^{-6}$'' não significa nada.
\end{conceito}

\begin{tarefa}{Advecção-difusão 1D estacionária}
\enunciado Resolva $\frac{d}{dx}(\rho u\phi) = \frac{d}{dx}\!\left(\Gamma\frac{d\phi}{dx}\right)$
em $[0,L]$, $\phi(0)=0$, $\phi(L)=1$, com solução exata
\[
  \frac{\phi(x)}{1} = \frac{\exp(Pe_L\,x/L)-1}{\exp(Pe_L)-1}, \qquad Pe_L=\frac{\rho u L}{\Gamma}.
\]
Implemente central, upwind, híbrido e QUICK. Varie o Péclet de célula de
$0{,}1$ a $10$.
\entregaveis
\begin{itemize}
  \item Gráficos $\phi(x)$ numérico vs. exato para $Pe_{\text{célula}}\in\{0{,}5;\,2;\,5\}$.
  \item Mapa ``esquema $\times$ $Pe$'' indicando: preciso, difusivo, oscilatório.
\end{itemize}
\criterio você reproduz os resultados clássicos do Cap. 5 de Versteeg \&
Malalasekera.
\reflexao\ Na camada limite de concentração junto a um jato de gás, o
Péclet de célula é alto ou baixo? O que isso implica para a malha?
\end{tarefa}

\begin{tarefa}{Primeiro protótipo de dispersão 2D}
\enunciado Resolva o transporte transiente 2D de um escalar passivo
(fração mássica $Y$) num campo de vento uniforme $\mathbf{u}=(U,0)$ com
difusividade constante $K$ (pense nela como difusividade turbulenta) e uma
fonte contínua pontual $\dot m$ (kg/s). Use FV com esquema TVD e integração
implícita no tempo. Contornos: entrada $Y=0$, saída gradiente nulo, laterais
simétricas.
\entregaveis
\begin{itemize}
  \item Mapa de contornos de $Y$ em estado estacionário e animação transiente.
  \item \textbf{Balanço global de massa}: gráfico de
        $\int\rho Y\,dV$ + massa que saiu $-$ massa injetada, ao longo do tempo.
        Deve ser zero até o erro de máquina (ou de tolerância do solver).
  \item Comparação com a solução analítica 2D de fonte contínua em vento
        uniforme (busque-a em Seinfeld \& Pandis ou derive-a).
\end{itemize}
\criterio erro de balanço de massa $<10^{-8}$ relativo; concordância com a
analítica longe da fonte.
\reflexao\ Este código é, conceitualmente, o núcleo de qualquer estudo de
dispersão em CFD. Liste tudo o que falta para ele ser ``industrial''
(você vai riscar itens dessa lista ao longo do curso).
\end{tarefa}

\begin{professor}
Guarde a lista da reflexão acima. No fim do curso, eu peço que os alunos a
releiam. É o momento em que percebem quanto avançaram --- e quanto de um
software comercial são decisões de modelagem, não ``mágica''.
\end{professor}

\begin{critico}
\begin{itemize}
  \item Um código que conserva massa exatamente pode ainda assim produzir
        concentrações negativas. Como? Isso é aceitável?
  \item Quando um resíduo cai seis ordens de grandeza, a solução está
        convergida? Proponha um critério melhor para um estudo de dispersão
        (dica: monitore grandezas de interesse, como o volume inflamável).
\end{itemize}
\end{critico}

\begin{chaves}
volumes finitos / \emph{finite volume method}; forma conservativa /
\emph{conservative form}; número de Péclet / \emph{Péclet number};
limitação / \emph{boundedness}; sub-relaxação / \emph{under-relaxation};
resíduo / \emph{residual}; multigrid; QUICK; esquema híbrido /
\emph{hybrid scheme}.
\end{chaves}

\begin{refs}
\begin{itemize}
  \item Versteeg, H. K.; Malalasekera, W. \emph{An Introduction to Computational Fluid Dynamics: The Finite Volume Method}. 2.\,ed. Pearson, 2007.
  \item Patankar, S. V. \emph{Numerical Heat Transfer and Fluid Flow}. Hemisphere, 1980.
  \item Moukalled, F.; Mangani, L.; Darwish, M. \emph{The Finite Volume Method in Computational Fluid Dynamics}. Springer, 2016.
  \item Saad, Y. \emph{Iterative Methods for Sparse Linear Systems}. 2.\,ed. SIAM, 2003.
\end{itemize}
\end{refs}

\begin{aprofundar}
\begin{itemize}[leftmargin=*]
  \item \textbf{Patankar (1980)}, caps.\ 3--5: curto, didático e ainda o melhor texto sobre ``por que'' os coeficientes devem ser positivos.
  \item \textbf{Moukalled, Mangani \& Darwish (2016)}, caps.\ 8--12: mesma teoria em malhas não estruturadas, com a notação do OpenFOAM.
  \item \textbf{Jasak, H.} \emph{Error Analysis and Estimation for the Finite Volume Method with Applications to Fluid Flows}. Tese, Imperial College, 1996. A base teórica do OpenFOAM; leia os caps.\ 3 e 4.
  \item \textbf{Saad (2003)}, caps.\ 6--7: gradiente conjugado e Krylov, para entender o que os solvers lineares fazem.
\end{itemize}
\end{aprofundar}

% ======================================================================
\chapter{Navier--Stokes incompressível e acoplamento pressão-velocidade}

\objetivos
\begin{itemize}
  \item Entender o papel da pressão em escoamentos incompressíveis.
  \item Implementar um método de projeção em malha deslocada.
  \item Conhecer SIMPLE, PISO e a interpolação de Rhie--Chow.
  \item Validar com o problema da cavidade (Ghia et al., 1982).
\end{itemize}

\begin{conceito}{Equações e o ``problema da pressão''}
\begin{gather*}
  \nabla\cdot\mathbf{u}=0,\\
  \frac{\partial\mathbf{u}}{\partial t}+(\mathbf{u}\cdot\nabla)\mathbf{u}
  = -\frac{1}{\rho}\nabla p + \nu\nabla^2\mathbf{u} + \mathbf{g}.
\end{gather*}
Não há equação de evolução para $p$: ela é o \emph{multiplicador de
Lagrange} que impõe $\nabla\cdot\mathbf{u}=0$. Tomando o divergente do
momento surge uma \textbf{equação de Poisson para a pressão} --- elíptica,
global, e geralmente a parte mais cara do cálculo.
\end{conceito}

\begin{dicionario}[Dicionário da equação: Navier--Stokes]
\begin{itemize}[leftmargin=*,nosep]
  \item $\partial\mathbf u/\partial t$ --- aceleração local.
  \item $(\mathbf u\cdot\nabla)\mathbf u$ --- advecção, \emph{não linear}: é
        daqui que nasce a turbulência.
  \item $-\nabla p/\rho$ --- força de pressão; $p$ [Pa].
  \item $\nu\nabla^2\mathbf u$ --- difusão viscosa; $\nu=\mu/\rho$ é a
        viscosidade cinemática [m$^2$/s] (ar: $\approx1{,}5\times10^{-5}$).
  \item $\mathbf g$ --- gravidade [m/s$^2$].
  \item $Re=UL/\nu$ --- razão entre inércia e viscosidade.
\end{itemize}
\textbf{Poisson da pressão}: tomando o divergente do momento e usando
$\nabla\cdot\mathbf u=0$, obtém-se $\nabla^2p=-\rho\,\nabla\cdot[(\mathbf u\cdot\nabla)\mathbf u]$.
\origem Euler (1757) escreveu as equações sem viscosidade. Claude-Louis
Navier (1822) incluiu o atrito a partir de um modelo molecular; George
Stokes (1845) as rederivou como teoria do contínuo. Osborne Reynolds (1883)
mostrou, com corante num tubo, que o adimensional $UL/\nu$ governa a
transição para a turbulência.
\end{dicionario}


\begin{conceito}{Malha deslocada, \emph{checkerboard} e Rhie--Chow}
Em malha colocalizada com interpolação central, um campo de pressão em
``tabuleiro de xadrez'' ($+1,-1,+1,\ldots$) tem gradiente discreto nulo e
é invisível ao momento: soluções espúrias. Remédios:
\begin{itemize}
  \item \textbf{Malha deslocada} (\emph{staggered}, Harlow \& Welch): velocidades
        nas faces, pressão nos centros. Robusto, didático; usado no FDS.
  \item \textbf{Interpolação de Rhie--Chow}: em malha colocalizada, a
        velocidade de face inclui um termo de gradiente de pressão
        compacto. É o que OpenFOAM e Fluent fazem.
\end{itemize}
\end{conceito}

\begin{conceito}{Projeção (Chorin) e família SIMPLE}
\textbf{Projeção}, por passo de tempo:
\begin{enumerate}
  \item Velocidade intermediária sem pressão:
        \[\mathbf{u}^* = \mathbf{u}^n + \Delta t\,\big[-(\mathbf{u}\cdot\nabla)\mathbf{u}+\nu\nabla^2\mathbf{u}\big]^n.\]
  \item Poisson: $\nabla^2 p^{n+1} = \frac{\rho}{\Delta t}\nabla\cdot\mathbf{u}^*$.
  \item Correção: $\mathbf{u}^{n+1}=\mathbf{u}^*-\frac{\Delta t}{\rho}\nabla p^{n+1}$.
\end{enumerate}
\textbf{SIMPLE} (Patankar \& Spalding, 1972): iterativo para regime
permanente, com correção de pressão e sub-relaxação. \textbf{SIMPLEC} e
\textbf{PISO} (Issa, 1986) são variantes; \textbf{PIMPLE} (OpenFOAM) combina
PISO e SIMPLE para transientes com passos grandes.
\end{conceito}

\begin{dicionario}[Dicionário da equação: projeção e siglas]
\begin{itemize}[leftmargin=*,nosep]
  \item $\mathbf u^*$ --- velocidade intermediária, que \emph{não} respeita
        $\nabla\cdot\mathbf u=0$; a pressão é o que a ``projeta'' no espaço
        de campos sem divergente.
  \item SIMPLE --- \emph{Semi-Implicit Method for Pressure-Linked Equations}.
  \item PISO --- \emph{Pressure-Implicit with Splitting of Operators}.
  \item MAC --- \emph{Marker-And-Cell}, a malha deslocada original.
\end{itemize}
\origem A base matemática é a decomposição de Helmholtz--Hodge: todo campo
vetorial é soma de um campo sem divergente e de um gradiente. Alexandre
Chorin (1968) e Roger Temam (1969) criaram o método de projeção de forma
independente. A malha deslocada vem do método MAC de Harlow \& Welch (Los
Alamos, 1965).
\end{dicionario}


\begin{lstlisting}[caption={Estrutura de um passo de projecao (pseudocodigo Julia)}]
function passo!(u, v, p, dt, nu, rho, malha)
    us, vs = predicao(u, v, dt, nu, malha)    # adveccao + difusao
    rhs = (rho/dt) .* divergente(us, vs, malha)
    resolve_poisson!(p, rhs, malha)          # CG ou multigrid
    corrige!(u, v, us, vs, p, dt, rho, malha)
    aplica_contornos!(u, v, malha)
    return maximum(abs, divergente(u, v, malha))   # deve ser ~0
end
\end{lstlisting}

\begin{tarefa}{Cavidade com tampa deslizante}
\enunciado Implemente o método de projeção 2D em malha deslocada para a
cavidade quadrada com tampa móvel ($U=1$, $L=1$). Rode $Re=100$, $400$ e
$1000$ até o regime permanente.
\entregaveis
\begin{itemize}
  \item Linhas de corrente e campo de pressão.
  \item Perfis de $u$ na linha vertical central e $v$ na horizontal central,
        sobrepostos aos dados tabelados de Ghia, Ghia \& Shin (1982).
  \item Histórico de $\max|\nabla\cdot\mathbf{u}|$ e de um ponto monitor.
  \item Estudo com 3 malhas (ex.: $32^2$, $64^2$, $128^2$) --- será reutilizado
        no Módulo 14.
\end{itemize}
\criterio perfis em $Re=100$ e $400$ visualmente sobrepostos aos de Ghia na
malha mais fina; divergente no nível da tolerância do solver de Poisson.
\reflexao\ Quanto do tempo de CPU foi gasto no Poisson? O que mudaria com
multigrid? Por que isso é ainda mais crítico em 3D?
\end{tarefa}

\begin{tarefa}{Escoamento ao redor de um obstáculo}
\enunciado Adapte o código para um canal 2D com um bloco quadrado de lado $H$
(represente-o por células bloqueadas). Rode $Re_H=20$ e $Re_H=150$.
\entregaveis\ Comprimento da zona de recirculação a jusante ($Re_H=20$);
frequência de desprendimento de vórtices e número de Strouhal ($Re_H=150$).
\criterio $St$ na faixa da literatura para cilindro quadrado ($\approx0{,}14$--$0{,}16$).
\reflexao\ Um módulo de FPSO é uma floresta de obstáculos: vasos, tubulações,
racks. O que acontece com uma nuvem de gás presa numa zona de recirculação?
Como isso afeta o tempo de residência e o volume inflamável?
\end{tarefa}

\begin{armadilha}
Declarar convergência porque ``os resíduos caíram''. Em cavidades e em
módulos semiconfinados, a solução pode ser fisicamente transiente
(desprendimento de vórtices); um solver estacionário vai oscilar ou
``convergir'' para um estado artificial. Sempre monitore pontos físicos.
\end{armadilha}

\begin{critico}
\begin{itemize}
  \item Gás vazado tem densidade diferente do ar. Em que condições a hipótese
        incompressível (ou de Boussinesq) ainda serve para dispersão? E para
        explosão?
  \item Por que códigos de explosão usam formulações compressíveis enquanto
        muitos estudos de dispersão usam solvers de baixo Mach?
\end{itemize}
\end{critico}

\begin{chaves}
acoplamento pressão-velocidade / \emph{pressure-velocity coupling}; método
de projeção / \emph{projection method}; malha deslocada / \emph{staggered
grid}; interpolação de Rhie--Chow; SIMPLE; PISO; PIMPLE; equação de Poisson
/ \emph{pressure Poisson equation}; cavidade / \emph{lid-driven cavity}.
\end{chaves}

\begin{refs}
\begin{itemize}
  \item Ferziger, J. H.; Perić, M.; Street, R. L. \emph{Computational Methods for Fluid Dynamics}. 4.\,ed. Springer, 2020.
  \item Chorin, A. J. Numerical solution of the Navier--Stokes equations. \emph{Math. Comp.}, 22:745--762, 1968.
  \item Patankar, S. V.; Spalding, D. B. A calculation procedure for heat, mass and momentum transfer in three-dimensional parabolic flows. \emph{Int. J. Heat Mass Transfer}, 15:1787--1806, 1972.
  \item Rhie, C. M.; Chow, W. L. Numerical study of the turbulent flow past an airfoil with trailing edge separation. \emph{AIAA J.}, 21(11):1525--1532, 1983.
  \item Ghia, U.; Ghia, K. N.; Shin, C. T. High-Re solutions for incompressible flow using the Navier--Stokes equations and a multigrid method. \emph{J. Comput. Phys.}, 48:387--411, 1982.
  \item Issa, R. I. Solution of the implicitly discretised fluid flow equations by operator-splitting. \emph{J. Comput. Phys.}, 62:40--65, 1986.
\end{itemize}
\end{refs}

\begin{aprofundar}
\begin{itemize}[leftmargin=*]
  \item \textbf{Ferziger, Perić \& Street (2020)}, caps.\ 7--8: acoplamento pressão-velocidade em malhas estruturadas e não estruturadas.
  \item \textbf{Harlow, F. H.; Welch, J. E.} Numerical calculation of time-dependent viscous incompressible flow of fluid with free surface. \emph{Phys.\ Fluids}, 8:2182--2189, 1965. O nascimento da malha deslocada.
  \item \textbf{Brown, D. L.; Cortez, R.; Minion, M. L.} Accurate projection methods for the incompressible Navier--Stokes equations. \emph{J.\ Comput.\ Phys.}, 168:464--499, 2001. As sutilezas de condições de contorno na projeção.
  \item \textbf{Barba \& Forsyth (2018)}, ``12 steps'': se travar, refaça os passos 9--12.
\end{itemize}
\end{aprofundar}

% ======================================================================
\part{Turbulência e escoamentos atmosféricos}
% ======================================================================

\chapter{Turbulência: física e modelagem}

\objetivos
\begin{itemize}
  \item Estimar escalas turbulentas e o custo de DNS, LES e RANS.
  \item Entender a hipótese de Boussinesq e os modelos $k$-$\varepsilon$, $k$-$\omega$ e SST.
  \item Usar funções de parede e interpretar $y^+$.
  \item Reconhecer o papel do número de Schmidt turbulento na dispersão.
\end{itemize}

\begin{conceito}{Cascata de energia e escalas de Kolmogorov}
A energia entra nos grandes turbilhões (escala integral $\ell$) e é
transferida a escalas menores até ser dissipada na escala de Kolmogorov:
\[
  \eta=\left(\frac{\nu^3}{\varepsilon}\right)^{1/4},\qquad
  \frac{\ell}{\eta}\sim Re_\ell^{3/4}.
\]
Resolver tudo (DNS) em 3D exige $N\sim Re^{9/4}$ pontos --- inviável para
escoamentos industriais. Daí a modelagem.
\end{conceito}

\begin{dicionario}[Dicionário da equação: escalas turbulentas]
\begin{itemize}[leftmargin=*,nosep]
  \item $\varepsilon$ --- taxa de dissipação de energia cinética turbulenta
        por unidade de massa [m$^2$/s$^3$ = W/kg].
  \item $\eta$ --- escala de Kolmogorov [m]: tamanho dos menores
        turbilhões, onde a viscosidade converte energia em calor.
  \item $\ell$ --- escala integral [m]: tamanho dos turbilhões que contêm a
        energia (no módulo de FPSO, da ordem dos equipamentos).
  \item $u'$ --- flutuação de velocidade (rms) [m/s]; $Re_\ell=u'\ell/\nu$.
\end{itemize}
\textbf{Como se deriva $\eta$:} nas menores escalas só importam $\nu$
[m$^2$/s] e $\varepsilon$ [m$^2$/s$^3$]; a única combinação com dimensão de
comprimento é $(\nu^3/\varepsilon)^{1/4}$. Análise dimensional pura.
\origem A imagem da cascata é de Lewis Fry Richardson (1922). Andrei
Kolmogorov formalizou as hipóteses de similaridade em 1941 (teoria ``K41''),
um dos resultados mais influentes da física do século XX.
\end{dicionario}


\begin{conceito}{RANS, Boussinesq e o modelo $k$-$\varepsilon$}
Decompondo $u_i=\bar u_i+u_i'$ e promediando surge o tensor de Reynolds
$-\rho\overline{u_i'u_j'}$ (problema de fechamento). Com densidade variável
usa-se a média de Favre ($\tilde u = \overline{\rho u}/\bar\rho$).
\textbf{Hipótese de Boussinesq:}
\begin{gather*}
  -\overline{u_i'u_j'} = \nu_t\left(\frac{\partial\bar u_i}{\partial x_j}
  +\frac{\partial\bar u_j}{\partial x_i}\right)-\frac23 k\,\delta_{ij},\\
  \nu_t = C_\mu\frac{k^2}{\varepsilon}.
\end{gather*}
$k$-$\varepsilon$ padrão (Launder \& Spalding, 1974): $C_\mu=0{,}09$,
$C_{1\varepsilon}=1{,}44$, $C_{2\varepsilon}=1{,}92$, $\sigma_k=1{,}0$,
$\sigma_\varepsilon=1{,}3$. Variantes: RNG, \emph{realizable}. Alternativas:
$k$-$\omega$ (Wilcox) e \textbf{SST} (Menter, 1994), melhor em separação e
gradientes adversos.

\textbf{Flutuabilidade}: em escoamentos estratificados (gás denso, plumas
quentes) adiciona-se produção/destruição $G_b = -\frac{\nu_t}{\sigma_t}\,
\mathbf{g}\cdot\frac{\nabla\bar\rho}{\bar\rho}$ nas equações de $k$ (e
$\varepsilon$). Esquecer isso superestima a mistura de nuvens densas.
\end{conceito}

\begin{dicionario}[Dicionário da equação: RANS e $k$-$\varepsilon$]
\begin{itemize}[leftmargin=*,nosep]
  \item $\bar u$ --- média de Reynolds; $\tilde u$ --- média de Favre
        (ponderada pela massa); $u'$ --- flutuação.
  \item $k=\tfrac12\overline{u_i'u_i'}$ --- energia cinética turbulenta [m$^2$/s$^2$].
  \item $\nu_t$ --- viscosidade turbulenta [m$^2$/s]: propriedade do
        \emph{escoamento}, não do fluido; tipicamente $10^2$ a $10^5$ vezes $\nu$.
  \item $\delta_{ij}$ --- delta de Kronecker (1 se $i=j$, 0 caso contrário).
  \item $P_k=\mu_t S^2$ --- produção de $k$ por cisalhamento, $S=\sqrt{2S_{ij}S_{ij}}$.
\end{itemize}
\textbf{Equações completas} (forma de densidade variável):
\begin{align*}
  \frac{\partial(\rho k)}{\partial t}+\nabla\cdot(\rho\mathbf u k)
  &=\nabla\cdot\Big[\Big(\mu+\frac{\mu_t}{\sigma_k}\Big)\nabla k\Big]\\
  &\quad+P_k+G_b-\rho\varepsilon,\\
  \frac{\partial(\rho\varepsilon)}{\partial t}+\nabla\cdot(\rho\mathbf u\varepsilon)
  &=\nabla\cdot\Big[\Big(\mu+\frac{\mu_t}{\sigma_\varepsilon}\Big)\nabla\varepsilon\Big]\\
  &\quad+C_{1\varepsilon}\frac{\varepsilon}{k}\left(P_k+C_{3\varepsilon}G_b\right)
   -C_{2\varepsilon}\rho\frac{\varepsilon^2}{k}.
\end{align*}
\origem Reynolds (1895) propôs a decomposição média + flutuação. Joseph
Boussinesq (1877) postulou que as tensões turbulentas se comportam como as
viscosas, com uma viscosidade ``aparente''. Alexandre Favre (1965) introduziu
a média ponderada pela massa. O $k$-$\varepsilon$ veio de Jones \& Launder
(1972) e Launder \& Spalding (1974); o SST de Florian Menter (1994), na NASA.
\end{dicionario}


\begin{conceito}{Paredes, $y^+$ e anomalias conhecidas}
$y^+ = u_\tau y/\nu$. Lei logarítmica: $u^+ = \frac{1}{\kappa}\ln y^+ + B$,
$\kappa\approx0{,}41$, $B\approx5{,}0$, válida para $30\lesssim y^+\lesssim
300$. Funções de parede exigem a primeira célula nessa faixa; modelos
\emph{low-Re} exigem $y^+\sim1$. Em módulos offshore com milhares de
obstáculos, quase sempre se usam funções de parede --- ou, em códigos do
tipo PDR, os obstáculos nem são resolvidos (Módulo 13).

\textbf{Anomalia do ponto de estagnação:} o $k$-$\varepsilon$ padrão
superestima $k$ onde o escoamento incide em obstáculos, exagerando a mistura
local. Correções do tipo Kato--Launder e limitadores de produção existem.
\end{conceito}

\begin{dicionario}[Dicionário da equação: lei da parede]
\begin{itemize}[leftmargin=*,nosep]
  \item $\tau_w$ --- tensão de cisalhamento na parede [Pa].
  \item $u_\tau=\sqrt{\tau_w/\rho}$ --- velocidade de atrito [m/s].
  \item $y$ --- distância à parede; $y^+$ --- essa distância em ``unidades
        viscosas''; $u^+=\bar u/u_\tau$.
  \item $\kappa\approx0{,}41$ --- constante de von Kármán; $B\approx5{,}0$.
\end{itemize}
\origem Ludwig Prandtl e Theodore von Kármán (anos 1920--30) chegaram ao
perfil logarítmico; Clark Millikan (1938) mostrou que ele resulta de
compatibilizar a lei da parede com a lei do defeito na região de
sobreposição. As funções de parede para CFD são de Launder \& Spalding (1974).
\end{dicionario}


\begin{conceito}{LES e o número de Schmidt turbulento}
\textbf{LES} resolve os grandes turbilhões e modela os menores (Smagorinsky:
$\nu_t=(C_s\Delta)^2|\bar S|$). Captura flutuações e intermitência ---
crucial em incêndios (FDS é LES). Custo de 10 a 100 vezes o RANS.

\textbf{Schmidt turbulento} $Sc_t = \nu_t/D_t$, tipicamente $0{,}7$--$1{,}0$.
A difusividade de espécie é $D_t=\nu_t/Sc_t$: um parâmetro escolhido pelo
usuário que muda diretamente o tamanho da nuvem.
\end{conceito}

\begin{dicionario}[Dicionário da equação: LES e $Sc_t$]
\begin{itemize}[leftmargin=*,nosep]
  \item $\Delta$ --- largura do filtro LES, em geral o tamanho da célula.
  \item $|\bar S|=\sqrt{2\bar S_{ij}\bar S_{ij}}$ --- taxa de deformação
        resolvida [1/s]; $C_s\approx0{,}1$--$0{,}2$.
  \item $D_t$ --- difusividade turbulenta de massa [m$^2$/s].
  \item $Sc=\nu/D$ (molecular) e $Sc_t=\nu_t/D_t$ (turbulento).
\end{itemize}
\origem Joseph Smagorinsky criou seu modelo (1963) para previsão do tempo
--- LES nasceu na meteorologia. Douglas Lilly (1967) estimou $C_s\approx0{,}17$
teoricamente. O número de Schmidt homenageia Ernst Schmidt (1892--1975);
$Sc_t\approx1$ é herança da \emph{analogia de Reynolds} (momento e massa
transportados pelos mesmos turbilhões).
\end{dicionario}


\begin{tarefa}{Ordens de grandeza num módulo de FPSO}
\enunciado Para vento de 5 m/s atravessando um módulo de processo de 20 m de
altura, estime: $Re$, a dissipação $\varepsilon\sim u'^3/\ell$ (adote
intensidade turbulenta de 10--20\%), $\eta$, o número de pontos para DNS e
para uma LES que resolva 80\% da energia (estime $\Delta\sim\ell/10$).
Faça o mesmo para um jato de gás saindo de um furo de 25 mm a velocidade sônica.
\entregaveis\ Planilha ou script com as contas e uma tabela-resumo.
\criterio números coerentes e justificativa das hipóteses.
\reflexao\ Qual dos dois problemas (vento no módulo ou jato perto do furo)
é mais difícil de resolver numericamente, e por quê? Como os códigos
industriais contornam o problema do jato? (Spoiler: Módulo 7.)
\end{tarefa}

\begin{tarefa}{Canal turbulento 1D: do comprimento de mistura ao $k$-$\varepsilon$}
\enunciado Para escoamento plenamente desenvolvido entre placas, a equação
de momento se reduz a uma EDO em $y$:
$\frac{d}{dy}\left[(\nu+\nu_t)\frac{d\bar u}{dy}\right]=\frac{1}{\rho}\frac{dp}{dx}$.
(a) Feche com comprimento de mistura de Prandtl e amortecimento de van
Driest; (b) feche com $k$-$\varepsilon$ com funções de parede. Use
$Re_\tau=590$ e compare com dados de DNS de canal (Moser, Kim \& Mansour,
1999, disponíveis publicamente).
\entregaveis\ Perfil $u^+\times y^+$ em escala log, com lei da parede, DNS
e seus dois modelos; perfil de $\nu_t$.
\criterio a região logarítmica dos dois modelos coincide com a DNS dentro de
5\%.
\reflexao\ Que parte do perfil cada modelo erra mais, e por que isso
importaria (ou não) para a dispersão de um gás liberado junto ao convés?
\end{tarefa}

\begin{tarefa}{OpenFOAM: degrau descendente com dois modelos}
\enunciado Rode o tutorial \texttt{pitzDaily} (\texttt{simpleFoam}) com
$k$-$\varepsilon$ e com $k$-$\omega$ SST. Depois adicione um escalar passivo
injetado a montante e varie $Sc_t\in\{0{,}5;\,0{,}7;\,1{,}0\}$.
\entregaveis\ Comprimento de recolamento para cada modelo; mapa de $y^+$;
curvas de concentração numa linha a jusante para os três $Sc_t$.
\criterio você justifica, por escrito, qual modelo usaria e com qual $Sc_t$,
citando ao menos uma referência.
\reflexao\ A diferença entre modelos de turbulência foi maior ou menor que
a diferença entre valores de $Sc_t$? O que isso diz sobre onde concentrar
esforço de análise de sensibilidade?
\end{tarefa}

\begin{professor}
Quando um aluno me pergunta ``qual é o melhor modelo de turbulência?'', eu
devolvo: ``melhor para quê, validado contra o quê?''. Em dispersão de gás em
módulos congestionados, o $k$-$\varepsilon$ com correções de flutuabilidade
continua sendo a espinha dorsal dos códigos industriais, não por ser o mais
sofisticado, mas porque tem décadas de validação \emph{nesse tipo de
problema}. Sofisticação sem validação é só custo.
\end{professor}

\begin{critico}
\begin{itemize}
  \item Um modelo RANS dá a concentração \emph{média}. A ignição depende da
        concentração \emph{instantânea}. Como reconciliar isso? (Revisite no Módulo 8.)
  \item Se LES é mais fiel, por que a indústria de explosões ainda usa RANS
        em escala massiva?
\end{itemize}
\end{critico}

\begin{chaves}
cascata de energia / \emph{energy cascade}; escala de Kolmogorov; RANS;
média de Favre / \emph{Favre averaging}; hipótese de Boussinesq /
\emph{eddy-viscosity hypothesis}; função de parede / \emph{wall function};
$y^+$; SST; LES; Smagorinsky; número de Schmidt turbulento /
\emph{turbulent Schmidt number}; produção por empuxo / \emph{buoyancy production}.
\end{chaves}

\begin{refs}
\begin{itemize}
  \item Pope, S. B. \emph{Turbulent Flows}. Cambridge University Press, 2000.
  \item Wilcox, D. C. \emph{Turbulence Modeling for CFD}. 3.\,ed. DCW Industries, 2006.
  \item Launder, B. E.; Spalding, D. B. The numerical computation of turbulent flows. \emph{Comput. Methods Appl. Mech. Eng.}, 3:269--289, 1974.
  \item Menter, F. R. Two-equation eddy-viscosity turbulence models for engineering applications. \emph{AIAA J.}, 32(8):1598--1605, 1994.
  \item Moser, R. D.; Kim, J.; Mansour, N. N. Direct numerical simulation of turbulent channel flow up to $Re_\tau=590$. \emph{Phys. Fluids}, 11:943--945, 1999.
  \item Smagorinsky, J. General circulation experiments with the primitive equations. \emph{Mon. Weather Rev.}, 91:99--164, 1963.
\end{itemize}
\end{refs}

\begin{aprofundar}
\begin{itemize}[leftmargin=*]
  \item \textbf{Tennekes, H.; Lumley, J. L.} \emph{A First Course in Turbulence}. MIT Press, 1972. Raciocínio físico por ordens de grandeza --- exatamente o que este módulo exige.
  \item \textbf{Davidson, P. A.} \emph{Turbulence: An Introduction for Scientists and Engineers}. 2.\,ed. Oxford, 2015. Mais acessível que Pope, com boa história.
  \item \textbf{Pope (2000)}, caps.\ 6 (escalas), 7 (camadas limite) e 10--11 (RANS): referência definitiva.
  \item \textbf{Tominaga, Y.; Stathopoulos, T.} Turbulent Schmidt numbers for CFD analysis with various types of flowfield. \emph{Atmos.\ Environ.}, 41:8091--8099, 2007. Revisão essencial sobre $Sc_t$ em dispersão.
\end{itemize}
\end{aprofundar}

% ======================================================================
\chapter{Camada limite atmosférica e dispersão passiva}

\objetivos
\begin{itemize}
  \item Descrever a camada limite atmosférica (CLA) e sua estabilidade.
  \item Implementar o modelo gaussiano de pluma e entender suas hipóteses.
  \item Impor perfis de entrada coerentes com o modelo de turbulência
        (Richards \& Hoxey) e testar a homogeneidade horizontal.
  \item Dimensionar domínios computacionais segundo boas práticas.
\end{itemize}

\begin{conceito}{Estabilidade atmosférica}
A estabilidade controla a mistura vertical. Classes de Pasquill--Gifford:
A (muito instável) a F (muito estável), D neutra. Na teoria de similaridade
de Monin--Obukhov, a escala é o comprimento de Obukhov
\[
  L = -\frac{u_*^3\,\bar\theta}{\kappa\,g\,\overline{w'\theta'}} ,
\]
com $L<0$ instável, $L\rightarrow\infty$ neutra, $L>0$ estável. Em QRA,
combinações típicas são \textbf{D5} (neutra, 5 m/s) e \textbf{F2} (estável,
2 m/s); a segunda costuma gerar as maiores distâncias de dispersão em
terra. Offshore, a estabilidade sobre o mar tem particularidades (diferença
de temperatura ar-mar), e o ``vento dominante'' é dado pela rosa dos ventos
da locação.
\end{conceito}

\begin{dicionario}[Dicionário da equação: estabilidade]
\begin{itemize}[leftmargin=*,nosep]
  \item $u_*$ --- velocidade de atrito do vento junto à superfície [m/s].
  \item $\theta$ --- temperatura potencial [K]: a que o ar teria se
        trazido adiabaticamente a 1000 hPa.
  \item $\overline{w'\theta'}$ --- fluxo turbulento de calor sensível
        (cinemático) [K$\cdot$m/s]; positivo de dia (solo aquece o ar).
  \item $z/L$ --- parâmetro de estabilidade: razão entre produção de
        turbulência por empuxo e por cisalhamento.
\end{itemize}
\textbf{Classes de Pasquill} são escolhidas por tabela a partir de três
observações simples: vento a 10 m, insolação (de dia) e nebulosidade (à noite).
\origem Alexander Obukhov definiu $L$ em 1946; Monin \& Obukhov (1954)
criaram a teoria de similaridade. Frank Pasquill (1961) propôs as classes
para uso prático com observações meteorológicas de rotina; Frank Gifford
(1961) as converteu nas curvas de $\sigma_y$ e $\sigma_z$.
\end{dicionario}


\begin{conceito}{Modelo gaussiano de pluma}
Para fonte contínua $Q$ (kg/s) a altura $H$, vento $u$, com reflexão no solo:
\begin{multline*}
  C(x,y,z)=\frac{Q}{2\pi u\sigma_y\sigma_z}
  \exp\!\left(-\frac{y^2}{2\sigma_y^2}\right)\\
  \times\left[\exp\!\left(-\frac{(z-H)^2}{2\sigma_z^2}\right)
       +\exp\!\left(-\frac{(z+H)^2}{2\sigma_z^2}\right)\right].
\end{multline*}
$\sigma_y(x)$ e $\sigma_z(x)$ dependem da classe de estabilidade (ex.:
correlações de Briggs para campo aberto, classe D:
$\sigma_y=0{,}08x(1+0{,}0001x)^{-1/2}$, $\sigma_z=0{,}06x(1+0{,}0015x)^{-1/2}$;
classe F: $\sigma_y=0{,}04x(1+0{,}0001x)^{-1/2}$, $\sigma_z=0{,}016x(1+0{,}0003x)^{-1}$).
\textbf{Hipóteses:} gás passivo (sem empuxo nem momento), terreno plano sem
obstáculos, vento e estabilidade homogêneos, regime permanente.
\end{conceito}

\begin{dicionario}[Dicionário da equação: pluma gaussiana]
\begin{itemize}[leftmargin=*,nosep]
  \item $C$ --- concentração [kg/m$^3$]; $Q$ --- vazão da fonte [kg/s].
  \item $u$ --- velocidade do vento na altura da liberação [m/s].
  \item $\sigma_y,\sigma_z$ --- desvios-padrão lateral e vertical da pluma [m].
  \item $H$ --- altura efetiva: altura da fonte mais a elevação da pluma.
  \item Termo com $(z+H)$ --- \emph{fonte-imagem}: reflexão total no solo
        (o gás não ``atravessa'' o chão).
  \item Conversão para \%vol: $x=\dfrac{C\,R\,T}{p\,M}\times100$.
\end{itemize}
\textbf{Como se deriva:} a solução fundamental da difusão é uma gaussiana com
variância $\sigma^2=2Kt$; o vento transporta essa nuvem, trocando $t$ por $x/u$.
\origem G.\,I. Taylor (1921) mostrou como a dispersão turbulenta cresce com o
tempo; O.\,G. Sutton (1932, 1947) formulou a pluma; Pasquill--Gifford (1961)
e Briggs (1973) forneceram as curvas de $\sigma$ usadas até hoje.
\end{dicionario}


\begin{conceito}{CFD da CLA: perfis de entrada e homogeneidade}
Para CLA neutra com $k$-$\varepsilon$ (Richards \& Hoxey, 1993):
\begin{gather*}
  u(z)=\frac{u_*}{\kappa}\ln\frac{z+z_0}{z_0},\\
  k=\frac{u_*^2}{\sqrt{C_\mu}},\qquad
  \varepsilon(z)=\frac{u_*^3}{\kappa(z+z_0)} .
\end{gather*}
Problema clássico (Blocken et al., 2007): funções de parede de rugosidade
incompatíveis com $z_0$ fazem o perfil ``decair'' antes de chegar à região
de interesse. Teste: simule o domínio \emph{vazio} e compare perfis na
entrada e no local do obstáculo. Boas práticas de domínio (COST 732): $\sim
5H$ a montante, $\sim15H$ a jusante, topo a $\sim5H$ acima do obstáculo,
razão de bloqueio $<3\%$.
\end{conceito}

\begin{dicionario}[Dicionário da equação: perfis de entrada]
\begin{itemize}[leftmargin=*,nosep]
  \item $z_0$ --- comprimento de rugosidade aerodinâmica [m], da ordem de
        1/30 da altura dos elementos rugosos (mar aberto $\sim10^{-4}$ m;
        área industrial $\sim0{,}5$--$1$ m).
  \item $u_*$ --- obtida invertendo o perfil log com o vento de referência a 10 m.
  \item $C_\mu$ --- a mesma constante do $k$-$\varepsilon$ (0,09).
\end{itemize}
\textbf{Como se deriva:} impõe-se que o perfil log seja solução exata das
equações $k$-$\varepsilon$ com tensão constante e produção = dissipação. O
resultado é um conjunto \emph{coerente} com o modelo de turbulência.
\origem Richards \& Hoxey (1993), engenharia do vento; a classificação de
rugosidades é de Davenport, revisada por Wieringa (1992).
\end{dicionario}


\begin{tarefa}{Calculadora de pluma gaussiana}
\enunciado Implemente o modelo gaussiano com $\sigma$ de Briggs para as seis
classes. Para metano liberado a 2 kg/s, a 10 m de altura, calcule e plote
(i) a concentração no eixo ao nível do solo e (ii) o contorno de 50\% do LFL
em planta, para D5 e F2. Converta corretamente kg/m$^3$ para \%vol.
\entregaveis\ Código com testes (ex.: integral do fluxo de massa num plano
transversal deve igualar $Q$); figuras; tabela de distâncias até LFL e $\frac12$LFL.
\criterio o teste de conservação de fluxo passa em todas as distâncias.
\reflexao\ Metano liberado de um vaso a alta pressão é passivo? Em que
distância ele ``se torna'' passivo? Por que a gaussiana é usada em QRA mesmo
assim?
\end{tarefa}

\begin{tarefa}{Teste de homogeneidade horizontal}
\enunciado No OpenFOAM (ou no seu código 2D com $k$-$\varepsilon$), simule
um domínio vazio de 1000 m $\times$ 200 m com perfis de Richards \& Hoxey,
$u_*$ correspondente a 5 m/s a 10 m e $z_0=0{,}0002$ m (mar aberto) e
$z_0=0{,}1$ m (terra). Compare perfis de $u$, $k$ e $\varepsilon$ na entrada,
no meio e na saída.
\entregaveis\ Perfis sobrepostos; desvio percentual de $u$ e $k$ a 10 m de
altura; descrição da condição de parede usada.
\criterio desvio de $u$ menor que 5\% ao longo do domínio (ou explicação
documentada do porquê não foi possível).
\reflexao\ Se o perfil decai 20\% antes de atingir a plataforma, que vento
você \emph{realmente} simulou? Como isso aparece num relatório de
consultoria se ninguém fizer este teste?
\end{tarefa}

\begin{tarefa}{CFD vs. gaussiana em terreno aberto}
\enunciado Com seu código de dispersão do Módulo 3 estendido para 2D
vertical (ou 3D), use o perfil logarítmico de vento e uma difusividade
turbulenta vertical $K_z(z)=\kappa u_* z/Sc_t$ (teoria K). Compare a
concentração ao nível do solo com a gaussiana classe D.
\entregaveis\ Gráfico comparativo; discussão das diferenças no campo próximo
e distante.
\criterio você explica, com base física, as diferenças observadas.
\reflexao\ Que obstáculos, em uma unidade de processo real, invalidariam a
gaussiana completamente? É aí que o CFD ``paga a conta''.
\end{tarefa}

\begin{critico}
\begin{itemize}
  \item Simular 8 ou 12 direções de vento com igual detalhe é sensato, se a
        rosa dos ventos mostra 40\% de frequência num único setor?
  \item Calmaria (vento quase nulo) é o pior caso para dispersão? Para quais
        geometrias isso é verdade, e para quais não é?
\end{itemize}
\end{critico}

\begin{chaves}
camada limite atmosférica / \emph{atmospheric boundary layer, ABL};
estabilidade atmosférica / \emph{atmospheric stability}; classes de
Pasquill--Gifford; comprimento de Obukhov / \emph{Obukhov length}; rugosidade
aerodinâmica / \emph{aerodynamic roughness length, $z_0$}; pluma gaussiana /
\emph{Gaussian plume}; homogeneidade horizontal / \emph{horizontal homogeneity};
teoria K / \emph{K-theory}; rosa dos ventos / \emph{wind rose}.
\end{chaves}

\begin{refs}
\begin{itemize}
  \item Seinfeld, J. H.; Pandis, S. N. \emph{Atmospheric Chemistry and Physics}. 3.\,ed. Wiley, 2016. (Caps. sobre CLA e dispersão)
  \item Turner, D. B. \emph{Workbook of Atmospheric Dispersion Estimates}. 2.\,ed. Lewis Publishers, 1994.
  \item Richards, P. J.; Hoxey, R. P. Appropriate boundary conditions for computational wind engineering models using the $k$-$\varepsilon$ turbulence model. \emph{J. Wind Eng. Ind. Aerodyn.}, 46--47:145--153, 1993.
  \item Blocken, B.; Stathopoulos, T.; Carmeliet, J. CFD simulation of the atmospheric boundary layer: wall function problems. \emph{Atmos. Environ.}, 41:238--252, 2007.
  \item Franke, J. et al. \emph{Best Practice Guideline for the CFD Simulation of Flows in the Urban Environment}. COST Action 732, 2007.
\end{itemize}
\end{refs}

\begin{aprofundar}
\begin{itemize}[leftmargin=*]
  \item \textbf{Stull, R. B.} \emph{An Introduction to Boundary Layer Meteorology}. Kluwer, 1988. A bíblia da CLA; caps.\ 1, 5 e 9.
  \item \textbf{Arya, S. P.} \emph{Introduction to Micrometeorology}. 2.\,ed. Academic Press, 2001. Mais curto e direto que Stull.
  \item \textbf{Blocken, B.} Computational Fluid Dynamics for urban physics: importance, scales, possibilities, limitations and ten tips and tricks towards accurate and reliable simulations. \emph{Build.\ Environ.}, 91:219--245, 2015. Boas práticas transferíveis para plantas industriais.
  \item \textbf{Turner (1994)}: para fazer contas de pluma à mão e ganhar intuição.
\end{itemize}
\end{aprofundar}

% ======================================================================
\part{Liberação, dispersão e detecção}
% ======================================================================

\chapter{Termodinâmica e modelagem da fonte de liberação}

\objetivos
\begin{itemize}
  \item Calcular vazões de gás por orifício (subcrítica e crítica).
  \item Modelar a despressurização transiente de um vaso.
  \item Entender jatos subexpandidos e modelos de pseudo-fonte.
  \item Reconhecer a fonte como a maior incerteza ``escondida'' do estudo.
\end{itemize}

\begin{conceito}{Escoamento por orifício e bloqueio sônico}
O escoamento se torna \textbf{crítico (bloqueado)} quando
\[
  \frac{p_a}{p_0} \le \left(\frac{2}{\gamma+1}\right)^{\gamma/(\gamma-1)}
  \quad(\approx0{,}53\ \text{para}\ \gamma=1{,}4).
\]
Nesse caso, a vazão mássica independe da pressão a jusante:
\[
  \dot m = C_d A\, p_0\sqrt{\frac{\gamma M}{Z R T_0}
  \left(\frac{2}{\gamma+1}\right)^{(\gamma+1)/(\gamma-1)}} .
\]
$C_d$ tipicamente $0{,}6$--$0{,}85$ (furo de bordas vivas a bocal arredondado);
$Z$ é o fator de compressibilidade (relevante a altas pressões de processo).
Praticamente todo vazamento de gás de processo offshore é sônico na saída.
\end{conceito}

\begin{dicionario}[Dicionário da equação: vazão crítica]
\begin{itemize}[leftmargin=*,nosep]
  \item $p_0,T_0$ --- pressão [Pa] e temperatura [K] de estagnação
        (dentro do equipamento); $p_a$ --- pressão ambiente.
  \item $\gamma=c_p/c_v$ --- razão de calores específicos.
  \item $M$ --- massa molar [kg/mol!]; $R=8{,}314$ J/(mol$\cdot$K).
  \item $Z$ --- fator de compressibilidade (gás real; $Z=1$ ideal).
  \item $A$ --- área do furo [m$^2$]; $C_d$ --- coeficiente de descarga
        (contração da veia e atrito).
\end{itemize}
\textbf{Como se deriva:} escoamento isentrópico + conservação de energia dão
$\dot m(p_{\text{garganta}})$; essa função tem um máximo quando a garganta
atinge Mach 1. Baixar $p_a$ além desse ponto não aumenta a vazão: a
informação não sobe contra um escoamento sônico.
\origem Barré de Saint-Venant e Laurent Wantzel (1839), estudando escoamento
de ar por orifícios --- a ``fórmula de Saint-Venant--Wantzel''.
\end{dicionario}


\begin{conceito}{Despressurização (\emph{blowdown}) transiente}
Balanços de massa e energia no vaso, com vazão pelo furo e/ou pela válvula de
despressurização. Hipóteses limites: isotérmica ou isentrópica (adiabática
reversível). A vazão cai com o tempo, e com ela o tamanho da nuvem. Em QRA,
a duração do vazamento é tão importante quanto a vazão inicial: decide se a
nuvem atinge regime permanente e se o isolamento (ESD) e a
despressurização são eficazes.
\end{conceito}

\begin{dicionario}[Dicionário da equação: modelo mínimo de \emph{blowdown}]
Para gás ideal em vaso de volume fixo $V$:
\begin{gather*}
  \frac{dm}{dt}=-\dot m(p,T),\\
  \text{isotérmico: } p=p_0\frac{m}{m_0},\qquad
  \text{isentrópico: } p=p_0\left(\frac{m}{m_0}\right)^{\gamma},
\end{gather*}
com $T=T_0(p/p_0)^{(\gamma-1)/\gamma}$ no caso isentrópico. O caso real fica
entre os dois (troca de calor com as paredes). Modelos completos acoplam
equação de estado de gás real e balanço de energia com a parede.
\origem Balanço de massa e primeira lei aplicados a volume de controle; os
modelos industriais de despressurização (ex.: Haque et al., 1992) refinam
isso com estratificação líquido-vapor e transferência de calor.
\end{dicionario}


\begin{conceito}{Jatos subexpandidos e pseudo-fonte}
Na saída sônica, $p_{\text{saída}}>p_a$: o jato expande através de um
sistema de choques (disco de Mach) em poucos diâmetros. Resolver isso em CFD
de dispersão exigiria células submilimétricas num domínio de centenas de
metros. Solução: substituir o furo real por uma \textbf{pseudo-fonte}
(\emph{notional nozzle}) já expandida à pressão ambiente, conservando massa
(e, conforme o modelo, momento). No modelo de Birch et al. (1984), assumindo
temperatura de estagnação e velocidade sônica na pseudo-fonte:
\[
  \frac{d_{ps}}{d} = \sqrt{C_d\,\frac{p_0}{p_a}}
  \left(\frac{2}{\gamma+1}\right)^{(\gamma+1)/(4(\gamma-1))}.
\]
Outros modelos (Ewan \& Moodie, 1986; Birch et al., 1987) fazem hipóteses
diferentes sobre temperatura e conservação de momento, e dão diâmetros e
velocidades distintas. A escolha muda o campo próximo do jato.
\end{conceito}

\begin{dicionario}[Dicionário da equação: jato subexpandido]
\begin{itemize}[leftmargin=*,nosep]
  \item $d$ --- diâmetro do furo; $d_{ps}$ --- diâmetro da pseudo-fonte.
  \item $p_0/p_a$ --- razão de pressões; acima de $\approx1{,}9$ (ar) o jato
        é subexpandido.
  \item Disco de Mach: choque normal a $x_M\approx0{,}67\,d\sqrt{p_0/p_a}$
        do furo (correlação de Ashkenas \& Sherman, 1966).
\end{itemize}
\origem O conceito de bocal ``fictício'' equivalente aparece em Thring \&
Newby (1953) para jatos de densidade diferente; para jatos de gás natural a
alta pressão, foi desenvolvido pela British Gas (Birch et al., 1984, 1987).
\end{dicionario}


\begin{conceito}{Escolha de cenários: tamanhos de furo e orientação}
QRAs offshore discretizam furos em faixas (ex.: pequeno, médio, grande,
ruptura), com frequências obtidas de bases como o \emph{Risk Assessment Data
Directory} da IOGP ou do HSE. A orientação do jato (horizontal contra
obstáculo, vertical, para baixo contra o convés) muda drasticamente a nuvem:
um jato que colide com um vaso perde momento e forma nuvem maior e mais
lenta.
\end{conceito}

\begin{tarefa}{Calculadora de fonte e \emph{blowdown}}
\enunciado Para um separador de gás com volume de 30 m$^3$, pressão inicial
de 80 bar(a), 40\,$^\circ$C, gás com $M=19$ kg/kmol e $\gamma=1{,}28$:
(a) calcule a vazão inicial para furos de 5, 25 e 75 mm; (b) integre o
\emph{blowdown} isentrópico e isotérmico (EDOs) até 1 bar(g), com e sem uma
válvula de despressurização adicional; (c) verifique o instante em que o
escoamento deixa de ser crítico.
\entregaveis\ Função de vazão testada; curvas $\dot m(t)$ e $p(t)$; tabela
com massa total liberada e tempo para cair a 50\% da vazão inicial.
\criterio massa liberada integrada concorda com a variação de massa
no vaso (erro $<0{,}1\%$).
\reflexao\ Para o furo de 75 mm, uma simulação de dispersão em regime
permanente com a vazão inicial é conservadora? É realista? Quando seria
preferível usar uma vazão média ou uma simulação transiente?
\end{tarefa}

\begin{tarefa}{Pseudo-fontes em confronto}
\enunciado Derive você mesmo a expressão de Birch (1984) a partir da
conservação de massa. Implemente-a e mais um modelo da literatura (ex.: Ewan
\& Moodie, que assume temperatura ambiente na pseudo-fonte). Para o furo de
25 mm da tarefa anterior, compare $d_{ps}$, $u_{ps}$, $T_{ps}$ e fluxo de
momento.
\entregaveis\ Derivação; tabela comparativa; estimativa do tamanho de célula
necessário para resolver a pseudo-fonte com 3 a 5 células.
\criterio derivação correta e coerência dimensional.
\reflexao\ Se dois modelos ``consagrados'' diferem em 30\% no momento do
jato, quanto isso pode mudar a distância até o LFL? Como você documentaria
essa escolha num relatório?
\end{tarefa}

\begin{armadilha}
Gastar semanas refinando malha e modelo de turbulência enquanto a vazão
de entrada foi calculada com $C_d$ arbitrário, gás ideal a 150 bar e sem
transiente. A fonte é o primeiro termo da cadeia; seu erro se propaga
para tudo.
\end{armadilha}

\begin{critico}
\begin{itemize}
  \item Liberações bifásicas (condensado, \emph{flashing}) formam aerossóis
        e poças. Quais hipóteses deste módulo deixam de valer?
  \item Por que muitas QRAs limitam a duração do vazamento a um valor fixo,
        e quão defensável é isso?
\end{itemize}
\end{critico}

\begin{chaves}
escoamento crítico / \emph{choked flow}; coeficiente de descarga /
\emph{discharge coefficient}; despressurização / \emph{blowdown};
jato subexpandido / \emph{under-expanded jet}; disco de Mach / \emph{Mach
disk}; pseudo-fonte / \emph{notional nozzle, pseudo-source}; escoamento
bifásico / \emph{two-phase release}; \emph{flashing}; frequência de
vazamento / \emph{leak frequency}.
\end{chaves}

\begin{refs}
\begin{itemize}
  \item TNO. \emph{Methods for the Calculation of Physical Effects} (``Yellow Book''), CPR 14E. 3.\,ed., 1997/2005. (Cap. 2: liberações)
  \item Birch, A. D.; Brown, D. R.; Dodson, M. G.; Swaffield, F. The structure and concentration decay of high pressure jets of natural gas. \emph{Combust. Sci. Technol.}, 36:249--261, 1984.
  \item Birch, A. D.; Hughes, D. J.; Swaffield, F. Velocity decay of high pressure jets. \emph{Combust. Sci. Technol.}, 52:161--171, 1987.
  \item Ewan, B. C. R.; Moodie, K. Structure and velocity measurements in underexpanded jets. \emph{Combust. Sci. Technol.}, 45:275--288, 1986.
  \item CCPS. \emph{Guidelines for Chemical Process Quantitative Risk Analysis}. 2.\,ed. AIChE, 2000.
  \item IOGP. \emph{Risk Assessment Data Directory --- Process Release Frequencies}. (Report 434-01, edição vigente.)
\end{itemize}
\end{refs}

\begin{aprofundar}
\begin{itemize}[leftmargin=*]
  \item \textbf{Yellow Book (TNO)}, cap.\ 2: modelos de liberação com exemplos numéricos resolvidos.
  \item \textbf{Franquet, E. et al.} Free underexpanded jets in a quiescent medium: a review. \emph{Prog.\ Aerosp.\ Sci.}, 77:25--53, 2015. Revisão moderna da física do jato subexpandido e das pseudo-fontes.
  \item \textbf{Leung, J. C.} A generalized correlation for one-component homogeneous equilibrium flashing choked flow. \emph{AIChE J.}, 32:1743--1746, 1986. Porta de entrada para liberações bifásicas (método ômega).
  \item \textbf{Crowl \& Louvar (2019)}, cap.\ 4: modelos de fonte em nível introdutório, bom para revisar.
\end{itemize}
\end{aprofundar}

% ======================================================================
\chapter{Jatos, gases densos e nuvens inflamáveis}

\objetivos
\begin{itemize}
  \item Formular o transporte de mistura com densidade variável e empuxo.
  \item Distinguir regimes dominados por momento, empuxo e turbulência ambiente.
  \item Calcular volume inflamável e nuvem estequiométrica equivalente (Q9).
  \item Discutir médias temporais, flutuações e o critério de $\frac12$LFL.
\end{itemize}

\begin{conceito}{Equações de mistura com densidade variável}
Para uma espécie com fração mássica $Y$:
\begin{gather*}
  \frac{\partial(\bar\rho\tilde Y)}{\partial t}+\nabla\cdot(\bar\rho\tilde{\mathbf u}\tilde Y)\\
  =\nabla\cdot\left[\bar\rho\left(D+\frac{\nu_t}{Sc_t}\right)\nabla\tilde Y\right],\\[4pt]
  \frac{1}{\rho}=\frac{RT}{p}\sum_i\frac{Y_i}{M_i}.
\end{gather*}
O momento inclui $(\bar\rho-\rho_a)\mathbf g$. Para diferenças pequenas,
a aproximação de Boussinesq; para gás denso ou hidrogênio, densidade variável
completa.
\end{conceito}

\begin{dicionario}[Dicionário da equação: mistura]
\begin{itemize}[leftmargin=*,nosep]
  \item $Y$ --- fração mássica do gás vazado (0 a 1).
  \item $D$ --- difusividade molecular [m$^2$/s], desprezível frente a
        $\nu_t/Sc_t$ longe da fonte.
  \item $M_i$ --- massa molar de cada componente.
  \item Fração molar (=\,volumétrica para gás ideal), usada para LFL:
        $x_i=\dfrac{Y_i/M_i}{\sum_j Y_j/M_j}$.
\end{itemize}
\textbf{Aproximação de Boussinesq}: a densidade varia \emph{só} no termo de
empuxo $(\rho-\rho_a)\mathbf g$; nos demais termos usa-se $\rho_a$. Válida
para $|\rho-\rho_a|/\rho_a\ll1$ (regra prática: $<10\%$).
\origem A lei de mistura vem de Dalton (pressões parciais, 1801). A
aproximação é de Oberbeck (1879) e Boussinesq (1903), em estudos de convecção.
\end{dicionario}


\begin{conceito}{Jatos, plumas e gás denso}
\begin{itemize}
  \item \textbf{Jato livre de momento}: a concentração no eixo decai como
        $\bar Y_c\propto d_{\text{ef}}/x$, com $d_{\text{ef}}=d\sqrt{\rho_j/\rho_a}$.
  \item \textbf{Número de Froude} densimétrico $Fr=u_j/\sqrt{g' d}$, com
        $g'=g\,|\rho_j-\rho_a|/\rho_a$: separa o campo dominado pelo momento do
        dominado pelo empuxo.
  \item \textbf{Gás denso} (GLP, CO$_2$, gás frio de GNL): espalhamento
        gravitacional, nuvem rasa e larga, mistura vertical suprimida pela
        estratificação. Número de Richardson $Ri=g'h/u_*^2$ alto indica
        comportamento denso. Velocidade da frente de corrente de gravidade:
        $u_f\approx Fr_f\sqrt{g'h}$, com $Fr_f$ da ordem de 0,5 a 1.
\end{itemize}
\end{conceito}

\begin{dicionario}[Dicionário da equação: jatos e gás denso]
\begin{itemize}[leftmargin=*,nosep]
  \item $u_j$, $\rho_j$ --- velocidade e densidade do jato na fonte.
  \item $d_{\text{ef}}$ --- diâmetro efetivo: corrige a diluição para a
        diferença de densidade entre jato e ar.
  \item $g'$ --- gravidade reduzida [m/s$^2$]: o ``peso aparente'' do gás.
  \item $Fr$ --- inércia/empuxo. $Fr\gg1$: momento domina.
  \item $Ri$ --- empuxo/turbulência. $Ri\gg1$: a estratificação suprime a
        mistura vertical.
  \item $h$ --- altura da nuvem; $Fr_f$ --- Froude da frente.
\end{itemize}
\origem William Froude (1810--1879) usou o adimensional para escalar
modelos de navios; L.\,F. Richardson (1920) definiu seu número estudando a
supressão da turbulência atmosférica; o diâmetro efetivo é de Thring \&
Newby (1953).
\end{dicionario}


\begin{conceito}{Inflamabilidade, volume inflamável e Q9}
Limites inferior e superior de inflamabilidade (LFL/UFL; metano
$\approx5$--$15\%$vol). Para misturas, regra de Le Chatelier:
$\mathrm{LFL}_{\text{mist}} = \left(\sum_i y_i/\mathrm{LFL}_i\right)^{-1}$.
\begin{itemize}
  \item \textbf{Volume inflamável}: volume de células com LFL $\le c\le$ UFL.
        Superestima o perigo de explosão, pois mistura pobre perto do LFL
        queima devagar.
  \item \textbf{Nuvem estequiométrica equivalente (ESC, ``Q9'')}: cada célula
        inflamável é ponderada pela reatividade da sua mistura em relação à
        estequiométrica,
        \[
          Q9 = \sum_{\text{células}} V_i\,\frac{(S_L\,E)(c_i)}{(S_L\,E)_{\max}},
        \]
        com $S_L$ a velocidade laminar de chama e $E$ a razão de expansão.
        Usada para ligar dispersão a explosão (Módulos 12--13).
\end{itemize}
\end{conceito}

\begin{dicionario}[Dicionário da equação: inflamabilidade e Q9]
\begin{itemize}[leftmargin=*,nosep]
  \item LFL/UFL --- limites inferior/superior em \%vol no ar.
  \item $y_i$ --- fração molar do componente $i$ \emph{no combustível}
        (sem ar).
  \item $V_i$ --- volume da célula $i$ [m$^3$]; $c_i$ --- concentração local.
  \item $S_L$ --- velocidade laminar [m/s]; $E=\rho_u/\rho_b$ --- razão de
        expansão. O produto $S_L E$ é a taxa de \emph{produção de volume}
        da chama, e é isso que gera sobrepressão.
\end{itemize}
\origem Henry Le Chatelier (1891) propôs a regra de mistura. O conceito de
nuvem equivalente e a família de definições numeradas (``Q5'', ``Q8'',
``Q9'') vieram do desenvolvimento do código FLACS (Gexcon); Q9 se tornou o
padrão de fato (Hansen et al., 2013).
\end{dicionario}


\begin{conceito}{Médias, flutuações e o critério de $\frac12$LFL}
RANS fornece a média de Favre; a concentração instantânea flutua. Bolsões
acima do LFL existem onde a média está abaixo dele. Por isso se usam
critérios conservadores como $\frac12$LFL (ou $\frac{1}{4}$ para alguns fins)
e se discutem razões pico/média. Também por isso detectores e ignição não
``veem'' a média.
\end{conceito}

\begin{tarefa}{Corrente de gravidade (\emph{lock-exchange})}
\enunciado Estenda seu solver 2D do Módulo 4 com a aproximação de Boussinesq
e uma equação de transporte para a densidade (ou fração de gás). Inicialize
um canal com fluido denso à esquerda e leve à direita, separados por uma
comporta removida em $t=0$. Meça a velocidade da frente.
\entregaveis\ Sequência de imagens da frente; gráfico posição da frente
$\times$ tempo; $Fr_f = u_f/\sqrt{g'H}$ medido.
\criterio $Fr_f$ próximo de 0,5 (teoria de Benjamin para troca total em
canal) na fase de velocidade constante.
\reflexao\ Num vazamento de gás denso no convés de uma plataforma, o que
faz o papel da comporta? Que estruturas do convés (soleiras, drenos,
\emph{grating}) mudariam essa física?
\end{tarefa}

\begin{tarefa}{Pós-processador de nuvem inflamável}
\enunciado Escreva uma rotina que recebe um campo 3D de concentração (gere
um sintético a partir da gaussiana do Módulo 6, ou exporte do OpenFOAM em
VTK) e calcula: volume inflamável, volume acima de $\frac12$LFL, Q9, e
extensão máxima da nuvem. Use dados de $S_L(\phi)$ e $E(\phi)$ que você
gerará no Módulo 10 (por ora, use valores tabelados da literatura).
\entregaveis\ Biblioteca de pós-processamento com testes (ex.: campo
uniforme estequiométrico deve dar $Q9=V$); tabela para os cenários D5 e F2.
\criterio testes passam; $Q9<V_{\text{infl}}$ sempre.
\reflexao\ A razão $Q9/V_{\text{infl}}$ variou entre cenários? O que ela
revela sobre a ``qualidade'' da nuvem para uma explosão?
\end{tarefa}

\begin{tarefa}{OpenFOAM: jato numa caixa congestionada}
\enunciado Monte um caso 3D: caixa de 20 m $\times$ 10 m $\times$ 8 m,
aberta em duas faces, com cilindros representando tubulações, vento na
face aberta e um jato de metano a partir da pseudo-fonte do Módulo 7. Use um
solver de mistura com densidade variável e empuxo (consulte, na sua versão
do OpenFOAM, a família de solvers ``\texttt{buoyant}'' / ``\texttt{reacting}''
sem reações). Rode até regime permanente.
\entregaveis\ Isossuperfícies de LFL e $\frac12$LFL; histórico de $V_{\text{infl}}$
e Q9 como critério de convergência; estudo com duas malhas.
\criterio Q9 estabilizado (variação $<2\%$ nas últimas iterações) e diferença
entre malhas documentada.
\reflexao\ Compare a extensão da nuvem com a gaussiana para a mesma vazão.
Em que sentido os obstáculos aumentaram ou reduziram o perigo?
\end{tarefa}

\begin{professor}
O Q9 é uma das ideias mais úteis e mais mal-usadas da área. Útil porque
resume uma nuvem complexa num número que se conecta a explosões. Mal-usado
porque esconde hipóteses fortes: nuvem ``congelada'', homogeneidade
equivalente, ignição no pior momento. Use, mas saiba listar cada hipótese em
voz alta.
\end{professor}

\begin{critico}
\begin{itemize}
  \item Simular dispersão em regime permanente e depois somar uma ignição
        ``no pior instante'' é conservador? Quando não é?
  \item Hidrogênio é muito leve e muito reativo. Como isso muda a física do
        Q9 e dos limites de inflamabilidade?
\end{itemize}
\end{critico}

\begin{chaves}
jato flutuante / \emph{buoyant jet}; número de Froude densimétrico; gás
denso / \emph{dense gas, heavy gas}; corrente de gravidade / \emph{gravity
current}; número de Richardson; LFL/UFL / \emph{lower/upper flammability
limit}; regra de Le Chatelier; volume inflamável / \emph{flammable volume};
nuvem estequiométrica equivalente / \emph{equivalent stoichiometric cloud,
ESC, Q9}; flutuação de concentração / \emph{concentration fluctuation}.
\end{chaves}

\begin{refs}
\begin{itemize}
  \item Britter, R. E.; McQuaid, J. \emph{Workbook on the Dispersion of Dense Gases}. HSE Contract Research Report 17/1988.
  \item Chen, C. J.; Rodi, W. \emph{Vertical Turbulent Buoyant Jets: A Review of Experimental Data}. Pergamon, 1980.
  \item Hansen, O. R.; Gavelli, F.; Davis, S. G.; Middha, P. Equivalent cloud methods used for explosion risk and consequence studies. \emph{J. Loss Prev. Process Ind.}, 26:511--527, 2013.
  \item Benjamin, T. B. Gravity currents and related phenomena. \emph{J. Fluid Mech.}, 31:209--248, 1968.
  \item CCPS. \emph{Guidelines for Vapor Cloud Explosion, Pressure Vessel Burst, BLEVE, and Flash Fire Hazards}. 2.\,ed. AIChE/Wiley, 2010.
\end{itemize}
\end{refs}

\begin{aprofundar}
\begin{itemize}[leftmargin=*]
  \item \textbf{CCPS.} \emph{Guidelines for Use of Vapor Cloud Dispersion Models}. 2.\,ed. AIChE, 1996. Visão de conjunto de modelos e experimentos.
  \item \textbf{Simpson, J. E.} \emph{Gravity Currents in the Environment and the Laboratory}. 2.\,ed. Cambridge, 1997. Fotografias e física de correntes de gravidade.
  \item \textbf{Fischer, H. B. et al.} \emph{Mixing in Inland and Coastal Waters}. Academic Press, 1979, cap.\ 9: a melhor dedução de leis de jatos e plumas por análise dimensional.
  \item \textbf{Hansen et al.\ (2013)}: indispensável antes de usar Q9 em um estudo real.
\end{itemize}
\end{aprofundar}

% ======================================================================
\chapter{Detecção de gás e otimização de posicionamento}

\objetivos
\begin{itemize}
  \item Conhecer tipos de detectores e critérios de desempenho.
  \item Distinguir cobertura geográfica e cobertura por cenários.
  \item Formular o posicionamento como problema de otimização.
  \item Avaliar robustez da solução a incertezas de cenário.
\end{itemize}

\begin{conceito}{Detectores e votação}
\begin{itemize}
  \item \textbf{Pontuais} (catalíticos, infravermelho): medem \%LFL num ponto.
  \item \textbf{Caminho aberto} (\emph{open-path}): integram concentração ao
        longo de um feixe (LFL$\cdot$m) --- natural de calcular em CFD como
        integral de linha.
  \item \textbf{Votação}: 1oo$N$ (alarme), 2oo$N$ (ação executiva, ex.: ESD),
        reduz alarmes espúrios mas exige que a nuvem alcance dois detectores.
\end{itemize}
\end{conceito}

\begin{conceito}{Cobertura geográfica vs. cobertura por cenários}
A ISA TR84.00.07 descreve abordagens baseadas em desempenho:
\begin{itemize}
  \item \textbf{Geográfica}: fração da área/volume em que uma nuvem de
        tamanho-alvo seria detectada (ex.: critérios de nuvem esférica de
        diâmetro prescrito usados em algumas práticas).
  \item \textbf{Por cenários}: fração (ponderada por frequência) de cenários
        de dispersão simulados que seriam detectados, e em quanto tempo.
\end{itemize}
O segundo é o que o CFD viabiliza --- e o que exige mais cenários.
\end{conceito}

\begin{conceito}{Formulação de otimização}
Seja $a_{ij}=1$ se o detector candidato $j$ ``vê'' o cenário $i$ (concentração
acima do limiar), $f_i$ a frequência do cenário. Máxima cobertura com $K$
detectores e votação 1oo$N$:
\begin{gather*}
  \max\ \sum_i f_i z_i\\
  \text{sujeito a}\quad z_i\le\sum_j a_{ij}x_j,\qquad \sum_j x_j\le K,\\
  x_j,\,z_i\in\{0,1\}.
\end{gather*}
Para 2oo$N$: $2z_i\le\sum_j a_{ij}x_j$. Métodos: guloso (rápido, garantia
de aproximação para cobertura submodular), MILP exato, metaheurísticas
(algoritmos genéticos) para objetivos não lineares (ex.: tempo de detecção).
\end{conceito}

\begin{dicionario}[Dicionário da equação: cobertura máxima]
\begin{itemize}[leftmargin=*,nosep]
  \item $i$ --- cenários; $j$ --- posições candidatas a detector.
  \item $x_j\in\{0,1\}$ --- instalo (1) ou não (0) um detector em $j$.
  \item $z_i\in\{0,1\}$ --- cenário $i$ detectado.
  \item $a_{ij}$ --- matriz de visibilidade, calculada dos campos de CFD.
  \item $K$ --- orçamento de detectores; $f_i$ --- frequência do cenário.
\end{itemize}
\origem Problema de localização de máxima cobertura de Church \& ReVelle
(1974), criado para localizar serviços de emergência. A garantia do guloso
(atinge pelo menos $1-1/e\approx63\%$ do ótimo) é de Nemhauser, Wolsey \&
Fisher (1978).
\end{dicionario}


\begin{tarefa}{Otimizador de detectores}
\enunciado Gere um conjunto de ao menos 48 cenários (ex.: 8 direções de
vento $\times$ 3 furos $\times$ 2 orientações), usando seu modelo
gaussiano/jato ou campos do OpenFOAM. Defina uma grade de posições candidatas.
Resolva a máxima cobertura para $K=2,\ldots,12$ com (a) algoritmo guloso e
(b) MILP (JuMP + HiGHS), para 1oo$N$ e 2oo$N$.
\entregaveis\ Curva cobertura $\times$ $K$ para os quatro casos; mapa com
posições ótimas; comparação de tempo de CPU guloso vs. MILP.
\criterio o MILP nunca é pior que o guloso; diferença documentada.
\reflexao\ Onde a curva de cobertura ``satura''? Que argumento de custo-benefício
você apresentaria a um gerente de projeto?
\end{tarefa}

\begin{tarefa}{Robustez: o otimizador superajustado}
\enunciado Otimize com metade dos cenários (ex.: só direções de vento
pares) e avalie a cobertura na outra metade. Depois perturbe cada direção
de vento em $\pm15^\circ$ e reavalie.
\entregaveis\ Tabela de cobertura ``treino'' vs. ``teste''; proposta de
formulação robusta (ex.: maximizar a cobertura do pior subconjunto, ou
exigir redundância).
\criterio você quantifica a perda de cobertura fora da amostra.
\reflexao\ Este é o mesmo fenômeno de \emph{overfitting} de aprendizado de
máquina. Que práticas de ML (validação cruzada, regularização) têm análogo
aqui?
\end{tarefa}

\begin{critico}
\begin{itemize}
  \item Detectores posicionados para nuvens de regime permanente detectam
        a nuvem \emph{a tempo} de evitar a ignição? Como incluir tempo no
        problema?
  \item Um layout ótimo para o projeto continua ótimo depois de dez anos de
        modificações na planta? Que governança isso exige?
\end{itemize}
\end{critico}

\begin{chaves}
detector pontual / \emph{point gas detector}; detector de caminho aberto /
\emph{open-path gas detector, OPGD}; votação / \emph{voting, MooN};
cobertura geográfica / \emph{geographic coverage}; cobertura por cenários /
\emph{scenario coverage}; problema de máxima cobertura / \emph{maximal
covering problem}; programação linear inteira mista / \emph{MILP}.
\end{chaves}

\begin{refs}
\begin{itemize}
  \item ISA. \emph{ISA-TR84.00.07: Guidance on the Evaluation of Fire, Combustible Gas, and Toxic Gas System Effectiveness}. 2010 (e revisões).
  \item Church, R.; ReVelle, C. The maximal covering location problem. \emph{Papers Reg. Sci. Assoc.}, 32:101--118, 1974.
  \item Legg, S. W.; Benavides-Serrano, A. J.; Siirola, J. D.; Watson, J.-P.; Davis, S. G.; Bratteteig, A.; Laird, C. D. A stochastic programming approach for gas detector placement using CFD-based dispersion simulations. \emph{Comput. Chem. Eng.}, 47:194--201, 2012.
\end{itemize}
\end{refs}

\begin{aprofundar}
\begin{itemize}[leftmargin=*]
  \item \textbf{ISA-TR84.00.07}: leia por inteiro; é curto e define o vocabulário da área.
  \item \textbf{Williams, H. P.} \emph{Model Building in Mathematical Programming}. 5.\,ed. Wiley, 2013. Como formular problemas como MILP corretamente.
  \item \textbf{Documentação do JuMP.jl} (\url{https://jump.dev}): tutoriais de cobertura e localização de instalações.
  \item \textbf{Legg et al.\ (2012)}: a formulação estocástica com cenários de CFD que inspirou a Tarefa 9.1.
\end{itemize}
\end{aprofundar}

% ======================================================================
\part{Combustão e incêndios}
% ======================================================================

\chapter{Fundamentos de combustão para segurança}

\objetivos
\begin{itemize}
  \item Calcular estequiometria, temperatura adiabática e velocidade de chama.
  \item Situar chamas industriais no diagrama de regimes de combustão turbulenta.
  \item Conhecer os principais modelos de combustão turbulenta em CFD.
  \item Usar correlações de velocidade turbulenta de queima.
\end{itemize}

\begin{conceito}{Estequiometria e propriedades de chama laminar}
Razão de equivalência $\phi=(F/A)/(F/A)_{\text{esteq}}$. Para metano,
CH$_4$ + 2(O$_2$+3,76N$_2$): $9{,}5\%$vol estequiométrico.
\begin{itemize}
  \item \textbf{Temperatura adiabática} $T_{ad}$: $\sim2200$ K para HC-ar estequiométrico.
  \item \textbf{Velocidade laminar} $S_L$: metano $\approx0{,}37$--$0{,}40$ m/s;
        propano $\approx0{,}43$--$0{,}46$ m/s; hidrogênio $>2$ m/s (máximo em mistura rica).
  \item \textbf{Razão de expansão} $E=\rho_u/\rho_b\approx7$--$8$.
        A velocidade da chama vista por um observador fixo é $\sim E\cdot S_L$
        --- é essa expansão que ``empurra'' o gás e gera sobrepressão.
  \item Espessura de chama $\delta_L\approx\alpha/S_L$ (décimos de mm);
        número de Lewis $Le=\alpha/D$ governa instabilidades
        (hidrogênio pobre: $Le<1$, chamas celulares e aceleração).
\end{itemize}
\end{conceito}

\begin{dicionario}[Dicionário da equação: chama laminar]
\begin{itemize}[leftmargin=*,nosep]
  \item $\phi$ --- razão de equivalência: $<1$ pobre, $>1$ rica.
  \item $F/A$ --- razão combustível/ar (mássica ou molar).
  \item $\rho_u,\rho_b$ --- densidades do gás não queimado e queimado.
  \item $\delta_L$ --- espessura da chama [m]; $\alpha$ --- difusividade
        térmica; $D$ --- difusividade mássica do reagente limitante.
  \item $Le=\alpha/D$ --- número de Lewis.
\end{itemize}
\origem Mallard \& Le Chatelier (1883) propuseram a teoria térmica da chama:
$S_L\sim\sqrt{\alpha/\tau_{\text{reação}}}$. Zeldovich \& Frank-Kamenetskii
(1938) a tornaram quantitativa. O número de Lewis homenageia Warren K. Lewis
(MIT, 1922), um dos fundadores da engenharia química.
\end{dicionario}


\begin{conceito}{Combustão turbulenta: regimes e números adimensionais}
\[
  Da=\frac{\tau_t}{\tau_c}=\frac{\ell/u'}{\delta_L/S_L},\qquad
  Ka=\left(\frac{\delta_L}{\eta}\right)^2 .
\]
O diagrama de Borghi--Peters ($u'/S_L$ vs. $\ell/\delta_L$) separa regimes:
chamas enrugadas, corrugadas, zonas de reação finas, reator homogêneo.
Em explosões em módulos congestionados, $u'/S_L$ cresce muito na esteira
dos obstáculos --- daí a aceleração da chama.

\textbf{Velocidade turbulenta de queima}, exemplos de correlações:
\begin{gather*}
  \text{Gülder (1990):}\quad\frac{S_T}{S_L}=1+0{,}62\left(\frac{u'}{S_L}\right)^{1/2}Re_\ell^{1/4}\\
  \text{Zimont:}\quad S_T = A\,u'^{3/4}S_L^{1/2}\alpha^{-1/4}\ell^{1/4}
\end{gather*}
\end{conceito}

\begin{dicionario}[Dicionário da equação: combustão turbulenta]
\begin{itemize}[leftmargin=*,nosep]
  \item $\tau_t=\ell/u'$ --- tempo de giro dos turbilhões grandes [s].
  \item $\tau_c=\delta_L/S_L$ --- tempo químico da chama laminar [s].
  \item $Da$ --- turbulência/química. $Da\gg1$: chama fina, enrugada.
  \item $Ka$ --- química/menores turbilhões. $Ka>1$: turbilhões entram na
        zona de preaquecimento.
  \item $S_T$ --- velocidade turbulenta de queima [m/s]; $\alpha$ ---
        difusividade térmica [m$^2$/s]; $A\approx0{,}5$ (Zimont).
\end{itemize}
\origem Gerhard Damköhler (1940) fez a primeira teoria de chama turbulenta;
Béla Karlovitz (1953) estudou o estiramento da chama; o diagrama de regimes é
de Borghi (1985) e Peters (1986, 2000). Vladimir Zimont (1979) derivou a
correlação para chamas ``espessadas''.
\end{dicionario}


\begin{conceito}{Modelos de combustão em CFD}
\begin{itemize}
  \item \textbf{Pré-misturada} (explosões): variável de progresso $c$ com
        fechamento por $S_T$ (Zimont/TFC), densidade de superfície de chama,
        modelo de enrugamento de Weller (\texttt{XiFoam}), equação $G$.
  \item \textbf{Não pré-misturada} (incêndios): fração de mistura $Z$,
        chama de Burke--Schumann, \emph{flamelets}.
  \item \textbf{Genéricos}: \emph{Eddy Dissipation Model} (Magnussen \&
        Hjertager, 1977) --- taxa limitada pela mistura, $\dot\omega\propto
        \bar\rho\,\frac{\varepsilon}{k}\min(Y_F, Y_O/s)$; EDC; usados em
        incêndios em RANS e no FDS (variante de mistura rápida).
\end{itemize}
\end{conceito}

\begin{dicionario}[Dicionário da equação: modelos de combustão]
\begin{itemize}[leftmargin=*,nosep]
  \item $Y_F,Y_O$ --- frações mássicas de combustível e oxidante.
  \item $s$ --- massa de oxigênio por massa de combustível (estequiométrica).
  \item $\varepsilon/k$ --- inverso do tempo de mistura turbulenta: a reação
        é tão rápida que quem manda é a mistura.
  \item $Z$ --- fração de mistura: fração da massa local que veio do
        combustível; conservada (não reage), por isso fácil de transportar.
  \item $c$ --- variável de progresso (0 não queimado, 1 queimado).
\end{itemize}
\origem Burke \& Schumann (1928) introduziram a ``chama de folha fina''.
Brian Spalding propôs o \emph{eddy break-up} (1971), refinado por
Magnussen \& Hjertager (1977). Bray, Moss \& Libby (anos 1970--80) formalizaram
a variável de progresso; Henry Weller (anos 1990), o modelo de enrugamento
implementado no OpenFOAM.
\end{dicionario}


\begin{tarefa}{Propriedades de chama com Cantera}
\enunciado Usando Cantera (Python, ou chamado a partir de Julia) com o
mecanismo GRI-Mech 3.0: (a) calcule $T_{ad}(\phi)$ e $E(\phi)$ para metano,
propano e uma composição típica de gás de processo (ex.: 85\% CH$_4$, 8\%
C$_2$H$_6$, 4\% C$_3$H$_8$, 3\% CO$_2$); (b) calcule $S_L(\phi)$ com chama
livre 1D; (c) compare o LFL da mistura pela regra de Le Chatelier com o
ponto onde $S_L$ cai abruptamente; (d) tabele $S_L E$ normalizado ---
o fator de ponderação do Q9 (Módulo 8).
\entregaveis\ Curvas $T_{ad}$, $S_L$, $E$ e $S_L E$ versus $\phi$ e versus
\%vol; tabela exportada para o pós-processador do Módulo 8.
\criterio $S_L$ do metano estequiométrico dentro da faixa experimental.
\reflexao\ A cinética detalhada prevê chama ``morrendo'' perto do LFL de
forma suave ou abrupta? Como a escolha do mecanismo pode afetar Q9?
\end{tarefa}

\begin{tarefa}{Onde vivem as chamas industriais}
\enunciado Desenhe o diagrama de Borghi--Peters e posicione: (i) uma chama
de metano num módulo vazio com $u'=0{,}5$ m/s, $\ell=1$ m; (ii) a mesma chama
na esteira de um rack de tubos com $u'=10$ m/s, $\ell=0{,}1$ m; (iii) um
incêndio em poça de 10 m. Calcule $S_T$ com duas correlações para (i) e (ii).
\entregaveis\ Diagrama anotado; tabela de $Da$, $Ka$, $S_T$.
\criterio pontos corretamente posicionados, com contas reproduzíveis.
\reflexao\ As correlações de $S_T$ concordam entre si? Que fração da
incerteza de uma simulação de explosão você atribuiria a essa escolha?
\end{tarefa}

\begin{critico}
\begin{itemize}
  \item Se $S_T$ depende de $u'$ e a chama gera $u'$ ao empurrar gás pelos
        obstáculos, que tipo de comportamento dinâmico esperar? (Revisite no
        Módulo 13.)
  \item Por que modelos ``limitados pela mistura'' funcionam razoavelmente
        para incêndios e mal para explosões?
\end{itemize}
\end{critico}

\begin{chaves}
razão de equivalência / \emph{equivalence ratio}; temperatura adiabática de
chama / \emph{adiabatic flame temperature}; velocidade laminar de queima /
\emph{laminar burning velocity}; razão de expansão / \emph{expansion ratio};
número de Damköhler, número de Karlovitz; diagrama de regimes /
\emph{regime diagram}; velocidade turbulenta de queima / \emph{turbulent
burning velocity}; fração de mistura / \emph{mixture fraction}; EDM, EDC;
\emph{flamelet}; variável de progresso / \emph{progress variable}.
\end{chaves}

\begin{refs}
\begin{itemize}
  \item Turns, S. R. \emph{An Introduction to Combustion}. 3.\,ed. McGraw-Hill, 2012.
  \item Poinsot, T.; Veynante, D. \emph{Theoretical and Numerical Combustion}. 3.\,ed. Edwards, 2012.
  \item Peters, N. \emph{Turbulent Combustion}. Cambridge University Press, 2000.
  \item Magnussen, B. F.; Hjertager, B. H. On mathematical modeling of turbulent combustion with special emphasis on soot formation and combustion. \emph{Proc. Combust. Inst.}, 16:719--729, 1977.
  \item Gülder, Ö. L. Turbulent premixed flame propagation models for different combustion regimes. \emph{Proc. Combust. Inst.}, 23:743--750, 1990.
  \item Goodwin, D. G. et al. \emph{Cantera: An object-oriented software toolkit for chemical kinetics, thermodynamics, and transport processes}. \url{https://www.cantera.org}.
\end{itemize}
\end{refs}

\begin{aprofundar}
\begin{itemize}[leftmargin=*]
  \item \textbf{Law, C. K.} \emph{Combustion Physics}. Cambridge, 2006. O tratamento mais completo de chamas laminares, estiramento e limites de inflamabilidade.
  \item \textbf{Glassman, I.; Yetter, R. A.; Glumac, N. G.} \emph{Combustion}. 5.\,ed. Academic Press, 2014. Clássico, com bons capítulos sobre detonação e limites.
  \item \textbf{Lewis, B.; von Elbe, G.} \emph{Combustion, Flames and Explosions of Gases}. 3.\,ed. Academic Press, 1987. Fonte histórica rica em dados experimentais.
  \item \textbf{Tutoriais do Cantera} (\url{https://www.cantera.org}): chama livre, equilíbrio e reatores.
\end{itemize}
\end{aprofundar}

% ======================================================================
\chapter{Incêndios: jato de fogo, poça e radiação}

\objetivos
\begin{itemize}
  \item Classificar incêndios de processo e seus fluxos térmicos típicos.
  \item Implementar o modelo de chama sólida com fator de forma.
  \item Entender a equação de transferência radiativa e seus modelos numéricos.
  \item Executar um incêndio em LES (FDS) com critério de resolução.
\end{itemize}

\begin{conceito}{Tipos de incêndio e efeitos}
\begin{itemize}
  \item \textbf{Jato de fogo} (\emph{jet fire}): ignição imediata de vazamento
        pressurizado; alto momento, fluxos de calor muito elevados no
        engolfamento (ordem de grandeza frequentemente citada: 200--350 kW/m$^2$),
        erosivo. Dimensiona proteção passiva (classes J/H).
  \item \textbf{Incêndio em poça} (\emph{pool fire}): queima de líquido
        acumulado; dominado por empuxo; fumaça e fuligem reduzem o poder
        emissivo em diâmetros grandes.
  \item \textbf{\emph{Flash fire}}: queima de nuvem sem sobrepressão
        significativa; efeito letal restrito ao interior da nuvem.
  \item \textbf{BLEVE / bola de fogo}: liberação súbita de líquido
        superaquecido.
\end{itemize}
Limiares usuais de radiação (verifique na norma do seu projeto): $\sim1{,}6$
kW/m$^2$ (exposição prolongada), $\sim5$ kW/m$^2$ (rotas de fuga, exposição
curta), $\sim12{,}5$ kW/m$^2$ (ignição de madeira, dano), $\sim37{,}5$ kW/m$^2$
(dano a equipamentos).
\end{conceito}

\begin{conceito}{Incêndio em poça: correlações e chama sólida}
\begin{itemize}
  \item Taxa de queima (Babrauskas): $\dot m''=\dot m''_\infty(1-e^{-k\beta D})$.
  \item Altura de chama (Thomas, 1963):
        $\frac{L}{D}=42\left(\frac{\dot m''}{\rho_a\sqrt{gD}}\right)^{0{,}61}$;
        (Heskestad): $L=0{,}235\,\dot Q^{2/5}-1{,}02\,D$ ($\dot Q$ em kW, $L$ e $D$ em m).
  \item Poder emissivo de superfície (SEP) decrescente com $D$ pela fumaça,
        ex. Mudan (1984):
        \[SEP=140\,e^{-0{,}12D}+20\,(1-e^{-0{,}12D})\ \ \text{[kW/m}^2\text{]}.\]
\end{itemize}
\textbf{Modelo de chama sólida}: a chama é um cilindro (inclinado pelo vento)
que irradia uniformemente; o fluxo no alvo é
\[
  q''=SEP\cdot F_{12}\cdot\tau_{\text{atm}},
\]
com $F_{12}$ o fator de forma geométrico e $\tau_{\text{atm}}$ a
transmissividade atmosférica (absorção por H$_2$O e CO$_2$).
\end{conceito}

\begin{dicionario}[Dicionário da equação: incêndio em poça]
\begin{itemize}[leftmargin=*,nosep]
  \item $\dot m''$ --- taxa de queima por área [kg/(m$^2\cdot$s)];
        $\dot m''_\infty$ --- valor para poça muito grande.
  \item $k\beta$ --- coeficiente de extinção $\times$ fator de feixe [1/m]:
        quão rápido a chama fica opticamente espessa com o diâmetro.
  \item $\dot Q=\dot m''\,\Delta H_c\,A_{\text{poça}}$ --- taxa de liberação de calor [kW].
  \item $SEP$ --- poder emissivo da superfície da chama [kW/m$^2$].
  \item $F_{12}$ --- fator de forma: fração da radiação que sai da chama e
        atinge o alvo,
        \[
          F_{12}=\frac{1}{A_1}\int_{A_1}\!\int_{A_2}\frac{\cos\theta_1\cos\theta_2}{\pi S^2}\,dA_2\,dA_1 .
        \]
  \item $\tau_{\text{atm}}$ --- fração não absorvida pelo ar úmido.
\end{itemize}
\origem Blinov \& Khudyakov (1957) e Burgess et al.\ (1961) mediram taxas de
queima; Babrauskas (1983) consolidou a correlação. Philip Thomas (1963)
derivou a altura de chama de experimentos com pilhas de madeira; Gunnar
Heskestad (1983), da teoria de plumas. O fator de forma deriva da lei do
cosseno de Lambert (1760); Hottel o sistematizou para engenharia (anos 1930--50).
\end{dicionario}


\begin{conceito}{Equação de transferência radiativa (RTE)}
Para meio absorvente-emissor, não espalhador, cinza:
\[
  \frac{dI(\mathbf r,\hat{\mathbf s})}{ds}=\kappa\left[I_b(T)-I\right],\qquad
  I_b=\frac{\sigma T^4}{\pi}.
\]
Modelos: \textbf{P1} (difusão da radiação, barato, ruim em meio opticamente
fino), \textbf{DOM/FVM} (ordenadas discretas; FDS, Fluent, OpenFOAM
\texttt{fvDOM}), \textbf{DTM} (transferência discreta, usado em códigos de
incêndio offshore). Propriedades: gás cinza, WSGG (soma ponderada de gases
cinzas), contribuição da fuligem (frequentemente dominante em chamas de HC).
Espessura óptica $\tau=\kappa L$ decide o regime.
\end{conceito}

\begin{dicionario}[Dicionário da equação: radiação]
\begin{itemize}[leftmargin=*,nosep]
  \item $I$ --- intensidade radiante [W/(m$^2\cdot$sr)] na direção $\hat{\mathbf s}$.
  \item $\kappa$ --- coeficiente de absorção [1/m]; $\tau=\kappa L$ ---
        espessura óptica ($\ll1$ transparente, $\gg1$ opaco).
  \item $I_b$ --- intensidade de corpo negro; $\sigma=5{,}67\times10^{-8}$
        W/(m$^2\cdot$K$^4$).
  \item $G=\int_{4\pi}I\,d\Omega$ --- radiação incidente [W/m$^2$]. No P1:
        $\nabla\cdot\left(\frac{1}{3\kappa}\nabla G\right)=\kappa\,(G-4\sigma T^4)$.
\end{itemize}
\origem Stefan (1879, experimental) e Boltzmann (1884, teórico) deram a lei
$\sigma T^4$; Kirchhoff (1859), a igualdade entre emissão e absorção; Planck
(1900), o espectro do corpo negro. A RTE moderna está em Chandrasekhar
(\emph{Radiative Transfer}, 1950). O método de transferência discreta é de
Lockwood \& Shah (1981); WSGG, de Hottel \& Sarofim (1967).
\end{dicionario}


\begin{conceito}{LES de incêndio e critério $D^*/\delta x$}
O FDS resolve as equações de baixo Mach com LES e combustão de mistura
rápida. A resolução é avaliada pelo diâmetro característico de fogo
\[
  D^*=\left(\frac{\dot Q}{\rho_\infty c_p T_\infty\sqrt g}\right)^{2/5},
\]
e recomenda-se $D^*/\delta x$ entre 4 (grosso) e 16 (fino), conforme os
estudos de V\&V do NIST/US NRC.
\end{conceito}

\begin{dicionario}[Dicionário da equação: $D^*$]
\begin{itemize}[leftmargin=*,nosep]
  \item $\dot Q$ --- potência do fogo [kW]; $\rho_\infty,c_p,T_\infty$ ---
        propriedades do ar ambiente.
  \item $D^*$ [m] --- escala de comprimento natural de uma pluma de fogo.
  \item $\delta x$ --- tamanho da célula; $D^*/\delta x$ --- ``quantas
        células cabem no fogo''.
\end{itemize}
\textbf{Como se deriva:} análise dimensional da pluma de empuxo: com $\dot Q$,
$g$ e propriedades do ar, a única escala de comprimento é $D^*$ (expoente
$2/5$, o mesmo da correlação de Heskestad).
\origem Teoria de plumas de Morton, Taylor \& Turner (1956) e Zukoski;
critério de resolução consolidado nos relatórios NUREG-1824 (US NRC/EPRI, 2007).
\end{dicionario}


\begin{tarefa}{Calculadora de incêndio em poça}
\enunciado Para uma poça de condensado de 8 m de diâmetro ($\dot m''_\infty$
e $k\beta$ da literatura para hidrocarbonetos leves), calcule $\dot m''$,
$\dot Q$, $L$ (Thomas e Heskestad), $SEP$ (Mudan) e o fluxo radiante a
distâncias de 5 a 60 m, por (i) modelo de fonte pontual
($q''=\chi_r\dot Q\tau/4\pi R^2$) e (ii) chama sólida cilíndrica vertical,
com fator de forma calculado por integração numérica (não use tabela!).
\entregaveis\ Código; gráfico $q''(R)$ dos dois modelos com linhas dos
limiares; tabela de distâncias para 5, 12,5 e 37,5 kW/m$^2$.
\criterio fator de forma numérico validado contra um caso com solução
fechada (ex.: disco paralelo a um elemento).
\reflexao\ Onde os modelos de fonte pontual e chama sólida divergem, e por
quê? Para posicionar uma rota de fuga a 15 m, qual você usaria?
\end{tarefa}

\begin{tarefa}{RTE 1D: da analítica ao P1}
\enunciado Resolva a RTE numa camada plana de gás isotérmico a $T_g$, entre
paredes negras frias, para espessuras ópticas $\tau\in\{0{,}1;\,1;\,10\}$.
Use (a) solução analítica com integrais exponenciais, (b) ordenadas discretas
com $S_4$ e $S_8$, (c) aproximação P1.
\entregaveis\ Fluxo na parede vs. $\tau$ para os três métodos; erro relativo.
\criterio DOM converge para a analítica com o refinamento angular.
\reflexao\ Em que regime o P1 falha, e por que isso importa para chamas de
jato longas e finas vistas de longe?
\end{tarefa}

\begin{tarefa}{FDS: incêndio em poça com estudo de malha}
\enunciado Monte no FDS uma poça de 2 m com HRR prescrito pela taxa de
queima de um hidrocarboneto. Rode com $D^*/\delta x=4$, 8 e 16. Coloque
medidores de fluxo radiante a 5, 10 e 20 m.
\entregaveis\ Altura média de chama (a partir do HRR por unidade de altura)
vs. Heskestad; fluxos radiantes vs. sua calculadora; custo computacional de
cada malha.
\criterio altura de chama dentro de $\pm20\%$ de Heskestad na malha média ou fina.
\reflexao\ O resultado de radiação a 20 m mudou com a malha? E o campo
próximo? O que isso ensina sobre onde investir resolução?
\end{tarefa}

\begin{professor}
Nos meus anos de sala de aula, a pergunta mais reveladora sobre incêndio é:
``qual a fração radiante?''. Quem responde ``0,3'' sem pestanejar ainda não
entendeu que esse número depende de combustível, diâmetro, fuligem e do
próprio modelo. Em CFD de incêndio, a radiação é frequentemente a maior
fonte de erro --- e a menos discutida nos relatórios.
\end{professor}

\begin{critico}
\begin{itemize}
  \item Um jato de fogo atingindo uma estrutura: o fluxo de calor incidente
        é o que importa, ou a temperatura da superfície? Como o CFD
        acoplado com condução na estrutura muda a resposta?
  \item O vento inclina a chama e aumenta a radiação a jusante. Seu
        estudo de incêndio considerou todas as direções da rosa dos ventos?
\end{itemize}
\end{critico}

\begin{chaves}
jato de fogo / \emph{jet fire}; incêndio em poça / \emph{pool fire};
poder emissivo de superfície / \emph{surface emissive power, SEP}; fator de
forma / \emph{view factor}; transmissividade atmosférica / \emph{atmospheric
transmissivity}; fração radiante / \emph{radiative fraction}; equação de
transferência radiativa / \emph{radiative transfer equation, RTE}; ordenadas
discretas / \emph{discrete ordinates}; fuligem / \emph{soot}; taxa de liberação
de calor / \emph{heat release rate, HRR}.
\end{chaves}

\begin{refs}
\begin{itemize}
  \item Drysdale, D. \emph{An Introduction to Fire Dynamics}. 3.\,ed. Wiley, 2011.
  \item Hurley, M. J. (ed.). \emph{SFPE Handbook of Fire Protection Engineering}. 5.\,ed. Springer, 2016.
  \item Modest, M. F. \emph{Radiative Heat Transfer}. 3.\,ed. Academic Press, 2013.
  \item Mudan, K. S. Thermal radiation hazards from hydrocarbon pool fires. \emph{Prog. Energy Combust. Sci.}, 10:59--80, 1984.
  \item Shokri, M.; Beyler, C. L. Radiation from large pool fires. \emph{J. Fire Prot. Eng.}, 1:141--150, 1989.
  \item Chamberlain, G. A. Developments in design methods for predicting thermal radiation from flares. \emph{Chem. Eng. Res. Des.}, 65:299--309, 1987.
  \item McGrattan, K. et al. \emph{Fire Dynamics Simulator: Technical Reference Guide} e \emph{User's Guide}. NIST Special Publication 1018/1019 (versão vigente).
  \item API. \emph{API Standard 521: Pressure-relieving and Depressuring Systems}. (Limiares de radiação.)
\end{itemize}
\end{refs}

\begin{aprofundar}
\begin{itemize}[leftmargin=*]
  \item \textbf{Quintiere, J. G.} \emph{Fundamentals of Fire Phenomena}. Wiley, 2006. Plumas, chamas e escalonamento com abordagem de engenheiro.
  \item \textbf{Howell, J. R.; Mengüç, M. P.; Siegel, R.} \emph{Thermal Radiation Heat Transfer}. 6.\,ed. CRC, 2015. O catálogo de fatores de forma e a teoria da RTE.
  \item \textbf{NUREG-1824} (US NRC/EPRI, 2007 e suplementos): verificação e validação de modelos de incêndio, incluindo FDS; leia para aprender a criticar resultados de fogo.
  \item \textbf{FDS User's Guide}: cap.\ sobre resolução de malha e o critério $D^*/\delta x$.
\end{itemize}
\end{aprofundar}

% ======================================================================
\part{Explosões}
% ======================================================================

\chapter{Explosões: fenomenologia e métodos empíricos}

\objetivos
\begin{itemize}
  \item Distinguir deflagração e detonação e entender o mecanismo de
        realimentação por obstáculos.
  \item Descrever ondas de explosão: sobrepressão, impulso, duração.
  \item Aplicar os métodos TNT, Multi-Energia e Baker--Strehlow--Tang.
  \item Resolver as equações de Euler 1D com um esquema de Godunov.
\end{itemize}

\begin{conceito}{Deflagração, detonação e realimentação}
\begin{itemize}
  \item \textbf{Deflagração}: chama subsônica, propagação por difusão de
        calor e radicais; sobrepressão cresce com a velocidade da chama.
  \item \textbf{Detonação}: onda de choque acoplada à reação, supersônica
        (velocidade de Chapman--Jouguet, $\sim1800$ m/s para HC-ar),
        sobrepressão de $\sim15$--$20$ bar.
  \item \textbf{DDT}: transição deflagração-detonação, possível em geometrias
        altamente congestionadas/confinadas ou com misturas muito reativas.
\end{itemize}
\textbf{Mecanismo de realimentação (Shchelkin)}: a expansão dos produtos
empurra o gás não queimado através dos obstáculos $\rightarrow$ turbulência na
esteira $\rightarrow$ $S_T$ aumenta $\rightarrow$ expansão mais rápida
$\rightarrow$ escoamento mais rápido $\rightarrow$\ldots\ \textbf{Congestionamento}
(densidade de obstáculos) e \textbf{confinamento} (paredes, convés, teto)
controlam a explosão --- não a massa total de combustível.
\end{conceito}

\begin{dicionario}[Dicionário da equação: termos de explosão]
\begin{itemize}[leftmargin=*,nosep]
  \item Velocidade de Chapman--Jouguet --- velocidade de detonação estável,
        em que os produtos saem a Mach 1 relativo à frente.
  \item Congestionamento --- obstáculos \emph{dentro} da nuvem (tubos,
        vasos), medido por razão de bloqueio e densidade de obstáculos.
  \item Confinamento --- superfícies que impedem a expansão (convés, paredes).
\end{itemize}
\origem David Chapman (1899) e Émile Jouguet (1905) criaram a teoria de
detonação; Kirill Shchelkin (anos 1940) mostrou, com espirais em tubos, que
obstáculos aceleram chamas até a detonação.
\end{dicionario}


\begin{conceito}{Ondas de explosão e escalonamento de Sachs}
Parâmetros: sobrepressão de pico $\Delta P$, duração da fase positiva $t_+$,
impulso $i=\int_0^{t_+}\Delta P\,dt$. Escalonamento de Sachs:
\begin{gather*}
  \bar R=\frac{R}{(E/p_0)^{1/3}},\qquad
  \overline{\Delta P}=\frac{\Delta P}{p_0},\\
  \bar t=\frac{t\,c_0}{(E/p_0)^{1/3}},\qquad
  \bar\imath=\frac{i\,c_0}{p_0^{2/3}E^{1/3}} .
\end{gather*}
A resposta estrutural depende da razão entre $t_+$ e o período natural da
estrutura: regime impulsivo, dinâmico ou quase-estático.
\end{conceito}

\begin{dicionario}[Dicionário da equação: escalonamento]
\begin{itemize}[leftmargin=*,nosep]
  \item $R$ --- distância ao centro da explosão [m].
  \item $E$ --- energia de combustão que contribui para a onda [J].
  \item $p_0$, $c_0$ --- pressão e velocidade do som ambientes.
  \item $\Delta P$ --- sobrepressão de pico [Pa]; $t_+$ --- duração da fase
        positiva [s]; $i$ --- impulso [Pa$\cdot$s].
  \item Barras indicam grandezas adimensionais.
\end{itemize}
\textbf{Ideia física:} explosões geometricamente semelhantes geram ondas
semelhantes quando a distância é medida na escala $(E/p_0)^{1/3}$ --- o raio
de uma esfera que conteria a energia $E$ a pressão $p_0$.
\origem Hopkinson (1915) e Cranz (1926) propuseram a lei da raiz cúbica para
explosivos; Robert Sachs (1944, \emph{Ballistic Research Laboratory}, EUA)
incluiu as condições ambientes.
\end{dicionario}


\begin{conceito}{Métodos empíricos de VCE}
\begin{itemize}
  \item \textbf{Equivalência TNT}: $W_{TNT}=\eta\,M\,\Delta H_c/\Delta H_{TNT}$
        com rendimento $\eta\sim1$--$10\%$. Fisicamente inadequado para
        deflagrações (forma de onda e campo próximo errados), ainda citado.
  \item \textbf{Multi-Energia (TNO)}: só a parte congestionada/confinada da
        nuvem gera explosão; energia $E\approx3{,}5$ MJ/m$^3$ $\times$ volume
        estequiométrico da região; curvas de explosão com ``força'' de 1 a 10
        (10 = detonação). A escolha da força é a grande incerteza (guias
        como GAME/GAMES relacionam força com congestionamento e escala).
  \item \textbf{Baker--Strehlow--Tang (BST)}: escolhe-se a velocidade máxima
        de chama (número de Mach) por tabela de reatividade, confinamento e
        congestionamento; curvas de sobrepressão e impulso correspondentes.
\end{itemize}
\end{conceito}

\begin{dicionario}[Dicionário da equação: métodos empíricos]
\begin{itemize}[leftmargin=*,nosep]
  \item $W_{TNT}$ --- massa equivalente de TNT [kg]; $\eta$ --- rendimento
        (fração da energia que vira onda de choque).
  \item $\Delta H_c$ --- calor de combustão ($\sim46$--$50$ MJ/kg para HC);
        $\Delta H_{TNT}\approx4{,}2$--$4{,}7$ MJ/kg.
  \item $3{,}5$ MJ/m$^3$ --- calor liberado por m$^3$ de mistura
        estequiométrica HC-ar, quase igual para todos os hidrocarbonetos.
  \item Força (1--10) --- índice de severidade da fonte no Multi-Energia;
        10 corresponde a detonação.
\end{itemize}
\origem O TNT equivalente vem de dados militares de explosivos sólidos
(compilações como Kingery \& Bulmash, 1984). O Multi-Energia surgiu no TNO
(van den Berg, 1985) depois que se percebeu que nuvens em áreas abertas
quase não geram sobrepressão. O projeto GAME (Eggen, TNO, 1998) correlacionou
a força com congestionamento. O BST é de Baker, Strehlow e Tang (1994--99).
\end{dicionario}


\begin{conceito}{Equações de Euler e o problema de Riemann}
\begin{gather*}
  \frac{\partial\mathbf U}{\partial t}+\frac{\partial\mathbf F(\mathbf U)}{\partial x}=0,\\
  \mathbf U=\begin{pmatrix}\rho\\ \rho u\\ \rho e_t\end{pmatrix},\quad
  \mathbf F=\begin{pmatrix}\rho u\\ \rho u^2+p\\ u(\rho e_t+p)\end{pmatrix}.
\end{gather*}
Descontinuidades obedecem às condições de Rankine--Hugoniot. O método de
Godunov resolve, em cada face, um problema de Riemann (exato ou aproximado:
Rusanov, HLL, HLLC, Roe). Segunda ordem via reconstrução MUSCL com
limitadores --- os mesmos do Módulo 2.
\end{conceito}

\begin{dicionario}[Dicionário da equação: Euler compressível]
\begin{itemize}[leftmargin=*,nosep]
  \item $\mathbf U$ --- vetor de variáveis conservadas; $\mathbf F$ --- fluxos.
  \item $e_t=e+u^2/2$ --- energia total específica; para gás ideal
        $e=p/[(\gamma-1)\rho]$.
  \item $c=\sqrt{\gamma p/\rho}$ --- velocidade do som. As ondas viajam com
        $u-c$, $u$ e $u+c$ (autovalores de $\partial\mathbf F/\partial\mathbf U$).
  \item Rankine--Hugoniot: num choque com velocidade $s$,
        $\mathbf F_R-\mathbf F_L=s\,(\mathbf U_R-\mathbf U_L)$.
\end{itemize}
\origem Euler (1757); Bernhard Riemann (1860) resolveu o ``problema de
Riemann'' para gases; William Rankine (1870) e Pierre-Henri Hugoniot (1887)
derivaram as condições de choque. Godunov (1959) transformou o problema de
Riemann em método numérico; seguiram-se Rusanov (1961), Roe (1981), Harten,
Lax \& van Leer (HLL, 1983) e Toro (HLLC, 1994).
\end{dicionario}


\begin{lstlisting}[caption={Fluxo de Rusanov (Lax--Friedrichs local) para Euler 1D}]
const gam = 1.4
pressao(U) = (gam-1)*(U[3] - 0.5*U[2]^2/U[1])
fluxo(U) = (p = pressao(U); u = U[2]/U[1];
            [U[2], U[2]*u + p, u*(U[3] + p)])
function rusanov(UL, UR)
    cL = sqrt(gam*pressao(UL)/UL[1]); cR = sqrt(gam*pressao(UR)/UR[1])
    smax = max(abs(UL[2]/UL[1]) + cL, abs(UR[2]/UR[1]) + cR)
    return 0.5*(fluxo(UL) + fluxo(UR)) .- 0.5*smax .* (UR .- UL)
end
\end{lstlisting}

\begin{tarefa}{Tubo de choque de Sod}
\enunciado Implemente Godunov de 1.ª ordem (Rusanov e HLL) e uma versão
MUSCL-Hancock de 2.ª ordem. Resolva Sod: $(\rho,u,p)_L=(1;0;1)$,
$(\rho,u,p)_R=(0{,}125;0;0{,}1)$, $\gamma=1{,}4$, até $t=0{,}2$. Implemente
também o solver de Riemann exato (Toro, Cap. 4) para comparação.
\entregaveis\ Perfis de $\rho$, $u$, $p$ numéricos vs. exato; identificação
da onda de rarefação, descontinuidade de contato e choque; estudo de
convergência na norma $L_1$.
\criterio ordem observada próxima de 1 para ambos (soluções descontínuas
limitam a ordem!) e erro do esquema de 2.ª ordem claramente menor.
\reflexao\ Por que a ordem observada é $<2$ mesmo para o esquema de 2.ª
ordem? Que consequência isso tem para o refinamento de malha em explosões?
\end{tarefa}

\begin{tarefa}{Explosão esférica e curvas empíricas}
\enunciado (a) Adapte o solver para simetria esférica (termo fonte
geométrico) e simule a explosão de uma esfera de gás a alta pressão; meça
$\Delta P$ e impulso versus distância e compare, em coordenadas de Sachs,
com uma curva de referência de explosão. (b) Para um módulo de processo com
2000 m$^3$ congestionados preenchidos com nuvem estequiométrica, calcule
$\Delta P(R)$ pelo TNT ($\eta=3\%$ e $10\%$) e pelo Multi-Energia com forças
4, 7 e 10 (digitalize as curvas do Yellow Book para estudo próprio, ou use
ajustes publicados).
\entregaveis\ Gráficos em coordenadas de Sachs; tabela de distâncias para
0,1, 0,3 e 1 bar por método.
\criterio contas reproduzíveis e diferenças entre métodos discutidas.
\reflexao\ A diferença entre força 4 e 7 é maior que a diferença entre TNT
e Multi-Energia? O que isso diz sobre onde está a incerteza --- no método ou
no julgamento do analista?
\end{tarefa}

\begin{critico}
\begin{itemize}
  \item Se a sobrepressão depende do congestionamento e não da massa total,
        por que tantos critérios de aceitação e legislações ainda falam em
        ``massa de inventário''?
  \item Buncefield revelou que vegetação densa pode agir como
        congestionamento. Que outros ``obstáculos invisíveis'' você
        consideraria num FPSO?
\end{itemize}
\end{critico}

\begin{chaves}
deflagração / \emph{deflagration}; detonação / \emph{detonation}; DDT;
explosão de nuvem de vapor / \emph{vapour cloud explosion, VCE};
congestionamento / \emph{congestion}; confinamento / \emph{confinement};
onda de explosão / \emph{blast wave}; impulso / \emph{impulse};
escalonamento de Sachs / \emph{Sachs scaling}; Multi-Energia / \emph{Multi-Energy
Method}; BST; problema de Riemann / \emph{Riemann problem}; MUSCL.
\end{chaves}

\begin{refs}
\begin{itemize}
  \item Bjerketvedt, D.; Bakke, J. R.; van Wingerden, K. Gas explosion handbook. \emph{J. Hazard. Mater.}, 52:1--150, 1997.
  \item van den Berg, A. C. The multi-energy method: a framework for vapour cloud explosion blast prediction. \emph{J. Hazard. Mater.}, 12:1--10, 1985.
  \item Tang, M. J.; Baker, Q. A. A new set of blast curves from vapor cloud explosion. \emph{Process Saf. Prog.}, 18(4):235--240, 1999.
  \item Toro, E. F. \emph{Riemann Solvers and Numerical Methods for Fluid Dynamics}. 3.\,ed. Springer, 2009.
  \item Sod, G. A. A survey of several finite difference methods for systems of nonlinear hyperbolic conservation laws. \emph{J. Comput. Phys.}, 27:1--31, 1978.
  \item CCPS. \emph{Guidelines for Vapor Cloud Explosion, Pressure Vessel Burst, BLEVE, and Flash Fire Hazards}. 2.\,ed., 2010.
  \item Ciccarelli, G.; Dorofeev, S. Flame acceleration and transition to detonation in ducts. \emph{Prog. Energy Combust. Sci.}, 34:499--550, 2008.
\end{itemize}
\end{refs}

\begin{aprofundar}
\begin{itemize}[leftmargin=*]
  \item \textbf{Bjerketvedt, Bakke \& van Wingerden (1997)}: o ``\emph{Gas Explosion Handbook}'' é a melhor introdução escrita por quem fez os experimentos; leia inteiro.
  \item \textbf{Eckhoff, R. K.} \emph{Explosion Hazards in the Process Industries}. 2.\,ed. Gulf, 2016. Visão ampla, incluindo poeiras e BLEVE.
  \item \textbf{Baker, W. E. et al.} \emph{Explosion Hazards and Evaluation}. Elsevier, 1983. Clássico sobre ondas de explosão e resposta estrutural.
  \item \textbf{Toro (2009)}, caps.\ 4 e 10: solver exato e HLL/HLLC --- os que você implementou.
  \item \textbf{Lee, J. H. S.} \emph{The Detonation Phenomenon}. Cambridge, 2008. Para quem quiser ir a fundo em DDT.
\end{itemize}
\end{aprofundar}

% ======================================================================
\chapter{CFD de explosões e análise probabilística}

\objetivos
\begin{itemize}
  \item Formular o problema de explosão como Euler/NS reativo.
  \item Entender a abordagem de porosidade e resistência distribuída (PDR).
  \item Simular aceleração de chama induzida por obstáculos em 1D.
  \item Construir curvas de excedência de sobrepressão e a carga acidental de projeto.
\end{itemize}

\begin{conceito}{Euler/NS reativo com variável de progresso}
Ao sistema de Euler adiciona-se a equação de progresso $c$ (0 = não queimado,
1 = queimado), e a energia libera calor $\Delta h_c$ conforme $c$ avança:
\begin{gather*}
  \frac{\partial(\rho c)}{\partial t}+\nabla\cdot(\rho\mathbf u c)
  =\nabla\cdot(\rho D_t\nabla c)+\dot\omega_c,\\
  \dot\omega_c\sim\rho_u S_T|\nabla c| .
\end{gather*}
O fechamento de $S_T$ (Módulo 10) é o coração do modelo; a turbulência
$u'$ vem do modelo $k$-$\varepsilon$ e da geração nos obstáculos.
\end{conceito}

\begin{dicionario}[Dicionário da equação: explosão reativa]
\begin{itemize}[leftmargin=*,nosep]
  \item $c$ --- variável de progresso; $\rho_u$ --- densidade não queimada.
  \item $\dot\omega_c$ --- taxa de reação [kg/(m$^3\cdot$s)]: ``massa que
        cruza a frente'' por volume, $\rho_uS_T$ vezes a área de chama por
        volume ($|\nabla c|$).
  \item $D_t$ --- difusividade turbulenta; a energia recebe
        $\dot\omega_c\,\Delta h_c$ como fonte de calor.
\end{itemize}
\origem Formulação de variável de progresso de Bray, Moss \& Libby (1977);
fechamento por velocidade turbulenta (TFC) de Zimont; códigos de explosão
offshore usam variações com a frente espessada numericamente para ocupar
algumas células (ex.: o modelo ``$\beta$'' do FLACS).
\end{dicionario}


\begin{conceito}{Porosidade e resistência distribuída (PDR)}
Num módulo real há milhares de tubos menores que a célula. A abordagem PDR
(base do FLACS e de outros códigos de explosão):
\begin{itemize}
  \item cada célula tem \textbf{porosidade volumétrica} $\beta_v$ e de face
        $\beta_f$ (fração aberta ao escoamento);
  \item obstáculos subgrid geram \textbf{arrasto} ($\propto\rho|\mathbf u|\mathbf u$)
        e \textbf{turbulência} adicional, e enrugam a chama;
  \item obstáculos grandes são resolvidos na malha.
\end{itemize}
Consequência: a malha não pode ser refinada livremente --- os submodelos
foram calibrados para certos tamanhos de célula relativos aos obstáculos.
``Convergência de malha'' tem significado diferente aqui.
\end{conceito}

\begin{dicionario}[Dicionário da equação: PDR]
\begin{itemize}[leftmargin=*,nosep]
  \item $\beta_v$ --- fração do volume da célula aberta ao gás (0 a 1).
  \item $\beta_f$ --- fração de cada face aberta ao escoamento.
  \item Arrasto subgrid: $\mathbf F_D\approx-C_D\,\rho\,|\mathbf u|\,\mathbf u\,A_{\text{obs}}/V$,
        com $A_{\text{obs}}$ a área projetada dos obstáculos na célula.
\end{itemize}
\origem A ideia de porosidade em CFD vem de trocadores de calor (Patankar \&
Spalding, 1974; Sha et al., 1982); Bjørn Hjertager, no Chr.\ Michelsen
Institute (Bergen), a adaptou para explosões nos anos 1980 --- semente do FLACS.
\end{dicionario}


\begin{conceito}{Análise probabilística de explosões}
Normas offshore (ex.: NORSOK Z-013, anexo sobre simulação probabilística de
explosões; ISO 13702) pedem:
\begin{enumerate}
  \item Muitos cenários de dispersão (vazão, direção, vento) $\rightarrow$
        distribuição de nuvens (Q9) com frequências.
  \item Modelo de ignição (probabilidade e tempo) $\rightarrow$ frequência
        de explosão por tamanho de nuvem.
  \item Simulações de explosão para um conjunto de nuvens e pontos de ignição
        $\rightarrow$ sobrepressão em alvos.
  \item \textbf{Curva de excedência}: frequência anual de exceder cada nível
        de sobrepressão. A \textbf{carga acidental de projeto (DAL)} é tipicamente
        a correspondente a $10^{-4}$/ano.
\end{enumerate}
\end{conceito}

\begin{dicionario}[Dicionário da equação: excedência]
\[
  F(\Delta P>p)=\sum_{i}f_i\,P_{\text{ign},i}\,
  \Pr\!\left(\Delta P_i>p\right)
\]
\begin{itemize}[leftmargin=*,nosep]
  \item $f_i$ --- frequência do cenário de vazamento $i$ [1/ano].
  \item $P_{\text{ign},i}$ --- probabilidade de ignição (e de ignição
        \emph{retardada}, que gera explosão).
  \item $\Pr(\Delta P_i>p)$ --- fração das simulações (pontos de ignição,
        instantes) em que o alvo vê mais que $p$.
\end{itemize}
\origem Abordagem consolidada na indústria norueguesa após Piper Alpha e
formalizada na NORSOK Z-013; ideia herdada da engenharia sísmica e
estrutural (curvas de perigo).
\end{dicionario}


\begin{tarefa}{Aceleração de chama em tubo com obstáculos}
\enunciado Estenda seu solver de Euler 1D (Módulo 12) com a variável de
progresso $c$, liberação de calor e $S_T=S_L(1+A\,u'/S_L)$ (ou Gülder), onde
$u'$ é aumentado localmente em zonas de obstáculo (modele como
$u'=C_{\text{obs}}|u|$ nessas zonas e um valor de fundo fora delas). Ignite em
uma extremidade fechada de um tubo de 10 m; a outra extremidade aberta.
Compare: sem obstáculos; 5 obstáculos; 10 obstáculos.
\entregaveis\ Diagrama $x$--$t$ da frente de chama; velocidade da chama vs.
posição; histórico de pressão em três sensores; pico de sobrepressão vs.
número de obstáculos.
\criterio você observa e explica o ciclo de realimentação (velocidade da
chama crescente ao atravessar obstáculos).
\reflexao\ Varie $C_{\text{obs}}$ em $\pm50\%$. O pico de pressão muda de forma
linear ou explosiva? O que isso revela sobre a sensibilidade de códigos de
explosão aos seus submodelos?
\end{tarefa}

\begin{tarefa}{Curva de excedência e DAL}
\enunciado Use os cenários da Tarefa do Módulo 9 (ou gere 200 sintéticos por
amostragem de Latin Hypercube), cada um com frequência e Q9. Defina uma
probabilidade de ignição dependente da vazão e uma função de resposta
$\Delta P = f(Q9)$ calibrada com seus resultados do Módulo 12 ou da Tarefa
anterior (com espalhamento aleatório para o ponto de ignição). Construa a
curva de excedência num alvo (ex.: parede corta-fogo).
\entregaveis\ Curva de excedência em escala log; DAL a $10^{-4}$/ano;
intervalo de confiança por \emph{bootstrap}.
\criterio pipeline reprodutível com semente aleatória fixada.
\reflexao\ Que entrada (frequências de vazamento, probabilidade de ignição,
função de resposta) mais move a DAL? Isso deveria mudar onde se investe
esforço de CFD?
\end{tarefa}

\begin{tarefa}{(Opcional avançado) \texttt{XiFoam} em geometria simples}
\enunciado No OpenFOAM, rode um caso de deflagração com \texttt{XiFoam} numa
caixa com obstáculos (ex.: adaptando um tutorial de câmara com obstáculos).
Compare a sobrepressão para nuvem estequiométrica e para $\phi=0{,}8$.
\entregaveis\ Históricos de pressão; campo de $c$ em quatro instantes; tempo
de CPU.
\criterio simulação estável com CFL controlado e conservação de massa global.
\reflexao\ O que seria necessário para usar este solver numa geometria de
módulo real com 5000 tubos? Por que existe a abordagem PDR?
\end{tarefa}

\begin{professor}
Em explosões, mais do que em qualquer outro tema deste curso, a confiança
vem da validação contra experimentos em escala (programas de testes em
escala real conduzidos nas décadas de 1990 e 2000 em geometrias offshore
são a base dos códigos comerciais). Se alguém lhe mostrar um resultado de
explosão em CFD sem dizer contra quais experimentos o código foi validado
\emph{naquele regime de congestionamento}, faça essa pergunta. Educadamente,
mas faça.
\end{professor}

\begin{critico}
\begin{itemize}
  \item A DAL a $10^{-4}$/ano depende de uma cadeia de modelos, cada um com
        incerteza. A incerteza total é maior ou menor que a soma? Como
        comunicá-la a quem projeta a estrutura?
  \item Painéis de alívio, \emph{deluge} e redução de congestionamento:
        quais mitigações o seu modelo 1D consegue representar, e quais não?
\end{itemize}
\end{critico}

\begin{chaves}
variável de progresso / \emph{progress variable}; porosidade e resistência
distribuída / \emph{porosity/distributed resistance, PDR}; aceleração de
chama / \emph{flame acceleration}; realimentação / \emph{positive feedback};
análise probabilística de explosão / \emph{probabilistic explosion analysis};
curva de excedência / \emph{exceedance curve}; carga acidental de projeto /
\emph{design accidental load, DAL}; probabilidade de ignição /
\emph{ignition probability}.
\end{chaves}

\begin{refs}
\begin{itemize}
  \item Arntzen, B. J. \emph{Modelling of Turbulence and Combustion for Simulation of Gas Explosions in Complex Geometries}. Tese de doutorado, NTNU, 1998.
  \item Hjertager, B. H. Simulation of transient compressible turbulent reactive flows. \emph{Combust. Sci. Technol.}, 27:159--170, 1982.
  \item Weller, H. G. et al. Application of a flame-wrinkling LES combustion model to a turbulent mixing layer. \emph{Proc. Combust. Inst.}, 27:899--907, 1998.
  \item Standards Norway. \emph{NORSOK Z-013: Risk and Emergency Preparedness Assessment} (edição vigente).
  \item ISO 13702. \emph{Petroleum and natural gas industries --- Control and mitigation of fires and explosions on offshore production installations}.
  \item Oil \& Gas UK / Energy Institute. \emph{Fire and Explosion Guidance}, 2007.
\end{itemize}
\end{refs}

\begin{aprofundar}
\begin{itemize}[leftmargin=*]
  \item \textbf{Lea, C. J.; Ledin, H. S.} \emph{A Review of the State-of-the-Art in Gas Explosion Modelling}. HSL/2002/02, Health and Safety Laboratory, 2002. Panorama crítico de códigos, incluindo PDR.
  \item \textbf{Arntzen (1998)}, caps.\ 2--4: a física e a numérica por trás do FLACS.
  \item \textbf{Poinsot \& Veynante (2012)}, caps.\ sobre combustão pré-misturada turbulenta.
  \item \textbf{Vinnem, J. E.} \emph{Offshore Risk Assessment}. Springer (edição mais recente). Como explosões entram no QRA offshore.
\end{itemize}
\end{aprofundar}

% ======================================================================
\part{Credibilidade, prática industrial e síntese}
% ======================================================================

\chapter{Verificação, validação e quantificação de incerteza}

\objetivos
\begin{itemize}
  \item Distinguir verificação de código, verificação de solução e validação.
  \item Aplicar o método das soluções manufaturadas (MMS) e o GCI.
  \item Usar métricas estatísticas de avaliação de modelos de dispersão.
  \item Propagar incertezas de entrada com amostragem e análise de sensibilidade.
\end{itemize}

\begin{conceito}{V\&V: as três perguntas}
\begin{itemize}
  \item \textbf{Verificação de código}: ``resolvo as equações corretamente?''
        --- MMS: escolha uma solução analítica arbitrária, calcule o termo
        fonte que ela exige, rode o código com esse termo e confirme a ordem
        de convergência teórica.
  \item \textbf{Verificação de solução}: ``quanto erro numérico há \emph{nesta}
        simulação?'' --- estimativa de erro de discretização.
  \item \textbf{Validação}: ``resolvo as equações \emph{certas}?'' ---
        comparação com experimentos relevantes, com incertezas de ambos os lados.
\end{itemize}
\end{conceito}

\begin{conceito}{Extrapolação de Richardson e GCI}
Com três malhas de razão de refinamento constante $r$ e soluções
$f_1$ (fina), $f_2$, $f_3$ (grossa):
\begin{gather*}
  p=\frac{\ln\left[(f_3-f_2)/(f_2-f_1)\right]}{\ln r},\\
  f_{\text{ext}}=f_1+\frac{f_1-f_2}{r^p-1},\\
  GCI_{\text{fina}}=\frac{F_s\,|(f_2-f_1)/f_1|}{r^p-1},
\end{gather*}
com fator de segurança $F_s=1{,}25$ para três malhas (Roache; Celik et al.,
2008). Atenção: só vale na faixa assintótica e com convergência monotônica.
\end{conceito}

\begin{dicionario}[Dicionário da equação: GCI]
\begin{itemize}[leftmargin=*,nosep]
  \item $f_1,f_2,f_3$ --- grandeza de interesse nas malhas fina, média e grossa.
  \item $h=(V/N)^{1/3}$ --- tamanho representativo de célula em 3D.
  \item $r=h_{\text{grossa}}/h_{\text{fina}}$ --- recomenda-se $r>1{,}3$.
  \item $p$ --- ordem \emph{observada}; compare com a teórica.
  \item $f_{\text{ext}}$ --- estimativa do valor em malha infinita.
\end{itemize}
\origem Lewis Fry Richardson (1911; ``\emph{deferred approach to the
limit}'', 1927). Patrick Roache (1994) propôs o GCI para padronizar relatos de
estudos de malha; Celik et al.\ (2008) o tornaram política editorial do
\emph{Journal of Fluids Engineering}.
\end{dicionario}


\begin{conceito}{Métricas de avaliação de dispersão}
Com $C_o$ observado e $C_p$ previsto (Hanna, Chang e colaboradores):
\begin{gather*}
  FB=\frac{\overline{C_o}-\overline{C_p}}{0{,}5(\overline{C_o}+\overline{C_p})},\qquad
  NMSE=\frac{\overline{(C_o-C_p)^2}}{\overline{C_o}\,\overline{C_p}},\\
  MG=\exp(\overline{\ln C_o}-\overline{\ln C_p}),\\
  VG=\exp\left[\overline{(\ln C_o-\ln C_p)^2}\right],\\
  FAC2=\text{fração com }0{,}5\le C_p/C_o\le2 .
\end{gather*}
Critérios frequentemente citados (Chang \& Hanna, 2004): $|FB|<0{,}3$,
$0{,}7<MG<1{,}3$, $NMSE<1{,}5$, $VG<4$, $FAC2>0{,}5$. Bancos de dados
clássicos: Prairie Grass (passivo), Thorney Island e Kit Fox (gás denso),
Burro/Coyote (GNL). O \emph{Model Evaluation Protocol} para GNL (Ivings et
al., 2007) é referência de procedimento.
\end{conceito}

\begin{dicionario}[Dicionário da equação: métricas]
\begin{itemize}[leftmargin=*,nosep]
  \item $C_o$, $C_p$ --- concentrações observada e prevista, pareadas no
        espaço e no tempo; a barra é a média sobre os pares.
  \item $FB>0$ --- modelo \emph{subestima} em média; $FB<0$, superestima.
  \item $NMSE$ --- espalhamento (dominado pelos valores altos).
  \item $MG$, $VG$ --- versões logarítmicas de viés e espalhamento, adequadas
        quando os dados variam ordens de grandeza.
  \item $FAC2$ --- a mais robusta: insensível a poucos valores extremos.
\end{itemize}
\origem Steven Hanna e colaboradores (1989, 1993) consolidaram o conjunto
avaliando modelos de gases perigosos contra experimentos de campo.
\end{dicionario}


\begin{tarefa}{MMS no seu código de volumes finitos}
\enunciado Escolha $\phi_m(x,y)=\sin(\pi x)\cos(\pi y)$ (ou outra), derive o
termo fonte para a equação de advecção-difusão do Módulo 3 e verifique a
ordem com 4 malhas para cada esquema implementado.
\entregaveis\ Derivação simbólica (pode usar \texttt{Symbolics.jl} ou SymPy);
gráfico log-log; tabela de ordens.
\criterio ordens teóricas recuperadas. Se não forem, você encontrou um bug
--- documente-o e corrija.
\reflexao\ Que tipos de erro a MMS \emph{não} detecta?
\end{tarefa}

\begin{tarefa}{GCI na cavidade}
\enunciado Use as três malhas da Tarefa da cavidade (Módulo 4) para calcular
$p$, $f_{\text{ext}}$ e GCI para: posição do centro do vórtice principal e
$u_{\min}$ na linha vertical central.
\entregaveis\ Tabela no formato de Celik et al. (2008).
\criterio $p$ entre 1 e 2,5 e discussão se a convergência é monotônica.
\reflexao\ Num estudo de dispersão, qual seria sua ``$f$''? Volume
inflamável? Distância ao LFL? Q9? Quais delas convergem de forma monotônica?
\end{tarefa}

\begin{tarefa}{Validação com dados reais: Prairie Grass}
\enunciado Obtenha os dados do experimento Prairie Grass (disponíveis na
literatura e em repositórios de avaliação de modelos). Compare sua
calculadora gaussiana (Módulo 6) com as concentrações observadas nos arcos
de 50 a 800 m, agrupando por classe de estabilidade.
\entregaveis\ Gráfico de dispersão $C_p\times C_o$ log-log com linhas de
fator 2; tabela de $FB$, $NMSE$, $MG$, $VG$, $FAC2$ por estabilidade.
\criterio métricas calculadas e comparadas aos critérios de aceitação.
\reflexao\ Seu modelo ``passou''. Isso o valida para um módulo de FPSO? Defina
\emph{domínio de validade}.
\end{tarefa}

\begin{tarefa}{Propagação de incerteza}
\enunciado Escolha um cenário de dispersão (gaussiana ou seu código 2D) e
trate como incertas: $C_d$, velocidade do vento, classe de estabilidade,
$Sc_t$. Defina distribuições plausíveis, amostre por Latin Hypercube (500
amostras) e calcule a distribuição da distância ao LFL. Estime índices de
Sobol de primeira ordem.
\entregaveis\ Histograma da saída; percentis 5/50/95; gráfico de índices de
Sobol.
\criterio resultados estáveis ao dobrar o número de amostras.
\reflexao\ O valor ``determinístico'' que você calcularia com os parâmetros
nominais está em qual percentil? Ele é conservador?
\end{tarefa}

\begin{critico}
\begin{itemize}
  \item Um modelo calibrado para um experimento é ``validado'' para ele?
  \item Quantos relatórios de CFD de segurança você já leu que apresentavam
        estimativa de erro numérico? Por que isso é raro, e o que mudaria se
        fosse obrigatório?
\end{itemize}
\end{critico}

\begin{chaves}
verificação / \emph{verification}; validação / \emph{validation};
solução manufaturada / \emph{method of manufactured solutions, MMS};
extrapolação de Richardson; índice de convergência de malha / \emph{grid
convergence index, GCI}; viés fracional / \emph{fractional bias}; fator de 2
/ \emph{FAC2}; quantificação de incerteza / \emph{uncertainty quantification,
UQ}; hipercubo latino / \emph{Latin hypercube sampling}; índices de Sobol.
\end{chaves}

\begin{refs}
\begin{itemize}
  \item Roache, P. J. \emph{Verification and Validation in Computational Science and Engineering}. Hermosa, 1998.
  \item Celik, I. B. et al. Procedure for estimation and reporting of uncertainty due to discretization in CFD applications. \emph{J. Fluids Eng.}, 130(7):078001, 2008.
  \item ASME. \emph{V\&V 20-2009: Standard for Verification and Validation in Computational Fluid Dynamics and Heat Transfer}.
  \item Hanna, S. R.; Chang, J. C.; Strimaitis, D. G. Hazardous gas model evaluation with field observations. \emph{Atmos. Environ.}, 27A:2265--2285, 1993.
  \item Chang, J. C.; Hanna, S. R. Air quality model performance evaluation. \emph{Meteorol. Atmos. Phys.}, 87:167--196, 2004.
  \item Ivings, M. J.; Jagger, S. F.; Lea, C. J.; Webber, D. M. \emph{Evaluating Vapor Dispersion Models for Safety Analysis of LNG Facilities}. Fire Protection Research Foundation, 2007.
  \item Saltelli, A. et al. \emph{Global Sensitivity Analysis: The Primer}. Wiley, 2008.
\end{itemize}
\end{refs}

\begin{aprofundar}
\begin{itemize}[leftmargin=*]
  \item \textbf{Oberkampf, W. L.; Roy, C. J.} \emph{Verification and Validation in Scientific Computing}. Cambridge, 2010. O livro de referência; caps.\ 5--8 e 12--13.
  \item \textbf{Roy, C. J.} Review of code and solution verification procedures for computational simulation. \emph{J.\ Comput.\ Phys.}, 205:131--156, 2005.
  \item \textbf{Smith, R. C.} \emph{Uncertainty Quantification: Theory, Implementation, and Applications}. SIAM, 2014.
  \item \textbf{Saltelli et al.\ (2008)}: índices de Sobol com exemplos práticos.
\end{itemize}
\end{aprofundar}

% ======================================================================
\chapter{Do CFD ao QRA: prática industrial}

\objetivos
\begin{itemize}
  \item Situar o CFD na cadeia HAZID $\rightarrow$ cenários $\rightarrow$
        consequências $\rightarrow$ risco.
  \item Decidir quando modelos integrais bastam e quando o CFD se justifica.
  \item Redigir premissas e memorial de cálculo defensáveis.
  \item Reconhecer armadilhas profissionais e éticas.
\end{itemize}

\begin{conceito}{Modelos integrais vs. CFD}
\begin{itemize}
  \item \textbf{Integrais} (ex.: PHAST e similares): rápidos, validados em
        campo aberto, ideais para milhares de cenários de QRA em terra.
        Não ``veem'' obstáculos nem confinamento.
  \item \textbf{CFD}: geometria real, confinamento, ventilação, interação
        entre jatos e estruturas. Caro; exige julgamento em malha, modelos e
        cenários. Essencial offshore e em áreas congestionadas.
  \item \textbf{Híbridos}: usar integrais para triagem e CFD para cenários
        dimensionantes; ou CFD para gerar ``funções de resposta'' usadas
        em milhares de cenários (\emph{surrogates}).
\end{itemize}
Medidas de risco: risco individual (IR), PLL (\emph{potential loss of life}),
curvas F-N; critérios ALARP.
\end{conceito}

\begin{dicionario}[Dicionário da equação: medidas de risco]
\begin{itemize}[leftmargin=*,nosep]
  \item $IR=\sum_i f_i\,P_{\text{fatal},i}(x,y)$ --- risco individual num
        ponto [1/ano]; multiplicado pela fração de presença vira IRPA.
  \item $PLL=\sum_i f_i\,N_i$ --- perda potencial de vidas [fatalidades/ano].
  \item Curva F-N --- frequência $F$ de eventos com $N$ ou mais fatalidades.
  \item ALARP --- \emph{as low as reasonably practicable}.
\end{itemize}
\origem O conceito ALARP vem do direito britânico (caso \emph{Edwards
v.\ National Coal Board}, 1949), incorporado pelo HSE; curvas F-N se
popularizaram com o \emph{Reactor Safety Study} (WASH-1400, 1975) e os
estudos de Canvey Island (1978).
\end{dicionario}


\begin{conceito}{O que um bom relatório de CFD de segurança contém}
\begin{enumerate}
  \item Objetivo e decisão de engenharia que o estudo suporta.
  \item Geometria: fonte (modelo 3D, nível de detalhe), simplificações.
  \item Cenários: critérios de seleção, frequências associadas.
  \item Modelagem da fonte: vazões, transientes, pseudo-fonte.
  \item Modelos físicos e numéricos: turbulência, esquemas, $Sc_t$, combustão,
        radiação --- com justificativa e referência.
  \item Malha e domínio: critérios, estudo de sensibilidade.
  \item Validação do código no regime de interesse.
  \item Resultados com incertezas e limitações explícitas.
  \item Registro de premissas e de mudanças.
\end{enumerate}
\end{conceito}

\begin{tarefa}{Memorial de premissas}
\enunciado Escreva um ``memorial de cálculo'' (máx. 4 páginas) para um
estudo hipotético de dispersão num módulo de separação de FPSO: escopo,
cenários, modelagem da fonte, configuração do CFD, critérios de convergência,
plano de verificação e validação, formato dos resultados.
\entregaveis\ Documento revisável; lista numerada de premissas com
justificativa e referência para cada uma.
\criterio um colega consegue reproduzir o estudo só com o memorial.
\reflexao\ Qual premissa você teria mais dificuldade de defender numa
auditoria? Como reduziria esse risco?
\end{tarefa}

\begin{tarefa}{Integral vs. CFD no mesmo cenário}
\enunciado Para um mesmo vazamento (vazão do Módulo 7), compare: (a) seu
modelo de jato/gaussiana; (b) seu caso OpenFOAM do Módulo 8 sem obstáculos;
(c) com obstáculos. Compare distância ao LFL, volume inflamável e Q9.
\entregaveis\ Tabela comparativa; parágrafo de recomendação técnica.
\criterio diferenças explicadas fisicamente.
\reflexao\ Em que tipo de instalação a diferença entre (a) e (c) mudaria uma
decisão de projeto (ex.: posição de tomada de ar do HVAC, distância de
equipamentos)?
\end{tarefa}

\begin{armadilha}
``Calibrar'' premissas até que o resultado caiba no critério de aceitação.
Ajustar $Sc_t$, ponto de ignição ou malha \emph{depois} de ver o resultado,
sem justificativa física independente, é viés de confirmação --- e em
segurança de processos isso tem nome mais grave. Documente decisões
\emph{antes} de rodar.
\end{armadilha}

\begin{critico}
\begin{itemize}
  \item Quem deve decidir o nível de conservadorismo de um estudo: o analista,
        o cliente, o regulador?
  \item Um resultado de CFD bonito pode reduzir a segurança de uma decisão?
        Como?
\end{itemize}
\end{critico}

\begin{chaves}
análise de perigos / \emph{HAZID}; modelo integral / \emph{integral model};
risco individual / \emph{individual risk}; PLL; curva F-N; ALARP; memorial de
cálculo / \emph{calculation basis}; registro de premissas /
\emph{assumptions register}; modelo substituto / \emph{surrogate model}.
\end{chaves}

\begin{refs}
\begin{itemize}
  \item CCPS. \emph{Guidelines for Chemical Process Quantitative Risk Analysis}. 2.\,ed. AIChE, 2000.
  \item Casal, J. \emph{Evaluation of the Effects and Consequences of Major Accidents in Industrial Plants}. 2.\,ed. Elsevier, 2018.
  \item Mannan, S. (ed.). \emph{Lees' Loss Prevention in the Process Industries}. 4.\,ed., 2012.
  \item Standards Norway. \emph{NORSOK Z-013} (edição vigente).
\end{itemize}
\end{refs}

\begin{aprofundar}
\begin{itemize}[leftmargin=*]
  \item \textbf{Vinnem, J. E.} \emph{Offshore Risk Assessment}, vols.\ 1 e 2. Springer (edição mais recente). Referência para QRA de instalações offshore, inclusive FPSO.
  \item \textbf{Pasman, H. J.} \emph{Risk Analysis and Control for Industrial Processes --- Gas, Oil and Chemicals}. Butterworth-Heinemann, 2015.
  \item \textbf{CCPS (2000)}, caps.\ 4--5: cálculo de risco e apresentação de resultados.
  \item \textbf{HSE.} \emph{Reducing Risks, Protecting People (R2P2)}. HSE Books, 2001. A filosofia ALARP explicada pelo regulador.
\end{itemize}
\end{aprofundar}

% ======================================================================
\chapter{Projeto final integrado: um módulo de FPSO}

\begin{professor}
Este é o momento em que tudo se encontra. Não espero um estudo
comercial: espero um estudo \emph{honesto}, em que cada número tenha origem
rastreável, cada premissa esteja escrita e cada incerteza esteja ao menos
estimada. Já vi projetos finais simples receberem nota máxima e projetos
lindíssimos serem reprovados por não conseguirem responder ``de onde veio
esse número?''.
\end{professor}

\section*{Descrição}
Um módulo de separação simplificado (ex.: 30 m $\times$ 20 m $\times$ 10 m),
com convés inferior e superior, dois vasos horizontais, um rack de tubulação
representado por cilindros e uma parede corta-fogo num dos lados. Gás de
processo com a composição do Módulo 10, a 60--80 bar.

\section*{Etapas e entregáveis}
\begin{enumerate}
  \item \textbf{Cenários} (Mód. 6, 7): 8 direções de vento $\times$ 3 furos
        $\times$ 2 orientações; vazões e durações com \emph{blowdown}.
  \item \textbf{Dispersão CFD} (Mód. 5, 6, 8): ao menos 12 casos no OpenFOAM,
        escolhidos com justificativa; os demais por modelo substituto treinado
        nesses casos (regressão, interpolação ou ML simples).
  \item \textbf{Detecção} (Mód. 9): layout ótimo com 2oo$N$ e análise de robustez.
  \item \textbf{Explosão} (Mód. 12, 13): função de resposta $\Delta P(Q9)$
        a partir do seu solver 1D e/ou Multi-Energia; curva de excedência na
        parede corta-fogo; DAL.
  \item \textbf{Incêndio} (Mód. 11): jato de fogo do furo médio atingindo o
        vaso adjacente (modelo de chama sólida ou FDS); fluxo na parede.
  \item \textbf{V\&V e incerteza} (Mód. 14): estudo de malha em um caso de
        dispersão; incerteza propagada até a DAL.
  \item \textbf{Relatório} (Mód. 15): máx. 25 páginas, no formato de memorial
        e resultados; repositório com \emph{pipeline} reprodutível (um comando
        regenera figuras e tabelas).
  \item \textbf{Opcional}: reescrever em Rust o pós-processador (Q9,
        detecção, excedência), com testes e \emph{benchmark} contra a versão
        Julia --- base para uma ferramenta própria de consequências.
\end{enumerate}

\section*{Rubrica de avaliação}
\begin{center}
\small
\begin{tabularx}{\linewidth}{@{}lcX@{}}
\toprule
Critério & Peso & O que se espera\\
\midrule
Rastreabilidade & 25\% & Todo número tem origem: cálculo, referência ou premissa registrada.\\
Física e modelagem & 20\% & Escolhas justificadas, coerentes com o regime do problema.\\
Qualidade numérica & 20\% & Convergência, malha, conservação, V\&V documentados.\\
Incerteza e limitações & 15\% & Quantificadas quando possível, discutidas sempre.\\
Reprodutibilidade & 10\% & Repositório executa de ponta a ponta.\\
Comunicação & 10\% & Recomendações claras para quem decide.\\
\bottomrule
\end{tabularx}
\end{center}

\begin{critico}
Releia a lista que você escreveu na Tarefa ``Primeiro protótipo de dispersão
2D'' (Módulo 3). Quantos itens foram riscados? Quais continuam abertos, e
quais você descobriu que nem sabia que existiam?
\end{critico}

% ======================================================================
\appendix
\chapter*{Apêndice A --- Glossário português/inglês}
\addcontentsline{toc}{chapter}{Apêndice A --- Glossário português/inglês}
\markboth{Glossário}{}

\begin{center}
\small
\begin{tabular}{@{}p{0.47\linewidth}p{0.47\linewidth}@{}}
\toprule
Português & Inglês\\
\midrule
acoplamento pressão-velocidade & pressure-velocity coupling\\
carga acidental de projeto & design accidental load (DAL)\\
coeficiente de descarga & discharge coefficient\\
congestionamento / confinamento & congestion / confinement\\
corrente de gravidade & gravity current\\
despressurização & blowdown / depressurisation\\
difusão numérica & numerical diffusion\\
escoamento crítico (bloqueado) & choked flow\\
estabilidade atmosférica & atmospheric stability\\
fator de forma & view factor\\
fração de mistura & mixture fraction\\
função de parede & wall function\\
incêndio em poça & pool fire\\
jato de fogo & jet fire\\
jato subexpandido & under-expanded jet\\
limitador de fluxo & flux limiter\\
limite inferior de inflamabilidade & lower flammability limit (LFL/LEL)\\
malha deslocada & staggered grid\\
nuvem estequiométrica equivalente & equivalent stoichiometric cloud (ESC, Q9)\\
número de Schmidt turbulento & turbulent Schmidt number\\
perda de contenção & loss of containment\\
poder emissivo de superfície & surface emissive power (SEP)\\
porosidade/resistência distribuída & porosity/distributed resistance (PDR)\\
pseudo-fonte & notional nozzle / pseudo-source\\
sobrepressão & overpressure\\
solução manufaturada & manufactured solution\\
sub-relaxação & under-relaxation\\
transição deflagração-detonação & deflagration-to-detonation transition (DDT)\\
velocidade laminar de queima & laminar burning velocity\\
verificação e validação & verification and validation (V\&V)\\
\bottomrule
\end{tabular}
\end{center}

\chapter*{Apêndice B --- Checklist de entregáveis}
\addcontentsline{toc}{chapter}{Apêndice B --- Checklist de entregáveis}
\markboth{Checklist}{}

Marque cada tarefa apenas quando código, figuras, tabela e reflexão
estiverem no repositório e no diário.

\begin{center}
\small
\begin{tabular}{@{}lp{0.70\linewidth}c@{}}
\toprule
Mód. & Tarefas & Feito\\
\midrule
0  & Ambiente e primeiro gráfico; autópsia de acidentes & $\square$\\
J  & Ambiente reprodutível; TDMA; caça às alocações; dispatch & $\square$\\
1  & Difusão 1D com três esquemas; von Neumann no papel & $\square$\\
2  & Torneio de esquemas de advecção; viscosidade numérica & $\square$\\
3  & Advecção-difusão 1D; protótipo de dispersão 2D & $\square$\\
4  & Cavidade (Ghia); obstáculo e Strouhal & $\square$\\
5  & Escalas no FPSO; canal 1D; \texttt{pitzDaily} com $Sc_t$ & $\square$\\
6  & Pluma gaussiana; homogeneidade da CLA; CFD vs. gaussiana & $\square$\\
7  & Fonte e \emph{blowdown}; pseudo-fontes & $\square$\\
8  & \emph{Lock-exchange}; pós-processador Q9; jato na caixa & $\square$\\
9  & Otimizador de detectores; robustez & $\square$\\
10 & Cantera; diagrama de regimes & $\square$\\
11 & Poça (chama sólida); RTE 1D; FDS com $D^*/\delta x$ & $\square$\\
12 & Sod; explosão esférica e métodos empíricos & $\square$\\
13 & Aceleração de chama 1D; curva de excedência; (XiFoam) & $\square$\\
14 & MMS; GCI; Prairie Grass; UQ & $\square$\\
15 & Memorial de premissas; integral vs. CFD & $\square$\\
16 & Projeto final integrado & $\square$\\
\bottomrule
\end{tabular}
\end{center}

\chapter*{Apêndice C --- Estante essencial}
\addcontentsline{toc}{chapter}{Apêndice C --- Estante essencial}
\markboth{Estante essencial}{}

Se você só puder ter dez livros sobre a mesa durante este curso:
\begin{enumerate}
  \item Versteeg \& Malalasekera --- \emph{An Introduction to CFD: The Finite Volume Method} (porta de entrada).
  \item Ferziger, Perić \& Street --- \emph{Computational Methods for Fluid Dynamics} (rigor numérico).
  \item LeVeque --- \emph{Finite Volume Methods for Hyperbolic Problems} (esquemas e choques).
  \item Toro --- \emph{Riemann Solvers and Numerical Methods for Fluid Dynamics} (explosões numéricas).
  \item Pope --- \emph{Turbulent Flows} (a física da turbulência).
  \item Poinsot \& Veynante --- \emph{Theoretical and Numerical Combustion} (combustão em CFD).
  \item Drysdale --- \emph{An Introduction to Fire Dynamics} (incêndios).
  \item CCPS --- \emph{Guidelines for Vapor Cloud Explosion, Pressure Vessel Burst, BLEVE, and Flash Fire Hazards}.
  \item TNO --- \emph{Yellow Book} (CPR 14E) (modelos de efeitos físicos).
  \item Mannan --- \emph{Lees' Loss Prevention in the Process Industries} (a enciclopédia).
\end{enumerate}

\chapter*{Apêndice D --- Trilhas de leitura aprofundada}
\addcontentsline{toc}{chapter}{Apêndice D --- Trilhas de leitura aprofundada}
\markboth{Trilhas de leitura}{}

\section*{Trilhas por objetivo}
Leia na ordem indicada; cada livro prepara o seguinte.
\begin{description}[leftmargin=0pt]
  \item[Dominar a numérica.] LeVeque (2007) $\rightarrow$ LeVeque (2002)
        $\rightarrow$ Patankar (1980) $\rightarrow$ Ferziger, Perić \&
        Street (2020) $\rightarrow$ Moukalled et al.\ (2016) $\rightarrow$
        Toro (2009).
  \item[Entender a física.] Tennekes \& Lumley (1972) $\rightarrow$ Pope
        (2000) $\rightarrow$ Stull (1988) $\rightarrow$ Turns (2012)
        $\rightarrow$ Law (2006) $\rightarrow$ Poinsot \& Veynante (2012).
  \item[Segurança aplicada.] Crowl \& Louvar (2019) $\rightarrow$
        \emph{Yellow Book} (TNO) $\rightarrow$ CCPS VCE (2010) $\rightarrow$
        Bjerketvedt et al.\ (1997) $\rightarrow$ Drysdale (2011)
        $\rightarrow$ Vinnem (\emph{Offshore Risk Assessment}).
  \item[Credibilidade.] Roache (1998) $\rightarrow$ Oberkampf \& Roy (2010)
        $\rightarrow$ Saltelli et al.\ (2008).
  \item[Ferramentas.] Think Julia $\rightarrow$ manual de Julia (Performance
        Tips) $\rightarrow$ notas MIT 18.337 $\rightarrow$ OpenFOAM
        \emph{User Guide} $\rightarrow$ FDS \emph{User's Guide}.
\end{description}

\section*{Periódicos para acompanhar}
\begin{itemize}
  \item \emph{Journal of Loss Prevention in the Process Industries} --- o
        principal da área; dispersão, explosões e QRA.
  \item \emph{Process Safety and Environmental Protection} (IChemE).
  \item \emph{Process Safety Progress} (AIChE/CCPS) --- muito aplicado.
  \item \emph{Journal of Hazardous Materials}.
  \item \emph{Fire Safety Journal} e \emph{Fire Technology} --- incêndios.
  \item \emph{Combustion and Flame}; \emph{Proceedings of the Combustion
        Institute} --- combustão fundamental.
  \item \emph{Atmospheric Environment}; \emph{Journal of Wind Engineering
        and Industrial Aerodynamics} --- dispersão e CLA.
  \item \emph{International Journal of Hydrogen Energy} --- segurança de
        hidrogênio, cada vez mais relevante.
\end{itemize}

\section*{Conferências}
IChemE \emph{Hazards} (Reino Unido); ISHPMIE (\emph{International Symposium on
Hazards, Prevention and Mitigation of Industrial Explosions}); AIChE Global
Congress on Process Safety; Mary Kay O'Connor Process Safety Center
International Symposium (Texas A\&M); ICHS (\emph{International Conference on
Hydrogen Safety}). Os anais são uma fonte rica de estudos de caso e de
validação.

\section*{Recursos gratuitos}
\begin{itemize}
  \item MIT OpenCourseWare 2.29 --- \emph{Numerical Fluid Mechanics}.
  \item ``CFD Python: 12 steps to Navier--Stokes'' (Barba).
  \item Guias de usuário, referência técnica e \emph{Validation Guide} do
        FDS (NIST) --- o guia de validação é uma aula de como comparar
        simulação e experimento.
  \item OpenFOAM \emph{User Guide} e o fórum CFD Online.
  \item NUREG-1824 (V\&V de modelos de incêndio, US NRC).
\end{itemize}

\section*{Como ler um artigo técnico}
Use o método das três passadas (Keshav, ``How to read a paper'',
\emph{ACM SIGCOMM Comput.\ Commun.\ Rev.}, 37(3), 2007): (1) título, resumo,
figuras e conclusões --- 10 minutos para decidir se vale ler; (2) leitura
completa sem se prender a derivações, anotando as hipóteses; (3) reproduza
mentalmente (ou em código!) os resultados principais. Para cada artigo,
registre no diário: \emph{qual problema, qual método, contra o que foi
validado, quais as limitações admitidas e quais as não admitidas}.

\begin{professor}
Uma última palavra. CFD aplicado à segurança de processos é uma disciplina
de \emph{humildade disciplinada}: humildade para admitir quanto não sabemos
sobre a física de uma nuvem turbulenta queimando entre tubulações;
disciplina para, mesmo assim, produzir números úteis, rastreáveis e honestos
para quem precisa projetar, operar e voltar para casa em segurança. Bom
estudo --- e mantenha o diário em dia.
\end{professor}

\end{document}
